\documentclass[11pt,a4paper,openany]{book}
\usepackage[utf8]{inputenc}
\usepackage[T1]{fontenc}
\usepackage[turkish]{babel}
\usepackage{geometry}
\geometry{a4paper, margin=2.5cm}
\usepackage{enumitem}
\usepackage{tcolorbox}
\tcbuselibrary{skins}
\usepackage{charter}
\usepackage[scaled]{helvet}
\usepackage{microtype}
\usepackage{amsmath}
\usepackage{xcolor}
\usepackage{titlesec}
\usepackage{fancyhdr}
\usepackage{needspace}
\usepackage{multicol}
\usepackage{longtable}
\usepackage{colortbl}
\usepackage{hyperref}

\definecolor{primarycolor}{HTML}{064E3B} % Deep Emerald Green
\definecolor{secondarycolor}{HTML}{059669} % Modern Green
\definecolor{accentcolor}{HTML}{E6F4EA} % Soft Light Green

\hypersetup{
    colorlinks=true,
    linkcolor=secondarycolor,
    urlcolor=cyan
}

% Sayfa yapısı ayarları
\setlength{\parindent}{0pt}
\setlength{\parskip}{8pt}

% Sayfa numaralarını sağ alta almak için fancyhdr ayarı
\pagestyle{fancy}
\fancyhf{} 
\fancyfoot[R]{\thepage} 
\renewcommand{\headrulewidth}{0pt} 
\renewcommand{\footrulewidth}{0pt} 

% Bölüm başlangıç sayfaları için de sağ alt numara ayarı
\fancypagestyle{plain}{
    \fancyhf{}
    \fancyfoot[R]{\thepage}
    \renewcommand{\headrulewidth}{0pt}
}

% Boş sayfaları kaldırmak için cleardoublepage redefinasyonu
\let\cleardoublepage\clearpage

% Akademik ve Modern Chapter / Section Başlık Stilleri
\titleformat{\chapter}[display]
  {\normalfont\huge\bfseries\sffamily\color{primarycolor}}
  {\filleft\textsf{\color{secondarycolor}\LARGE\chaptername\ \thechapter}}
  {10pt}
  {\filleft\Huge\bfseries}
  [\vspace{1ex}{\color{primarycolor}\hrule height 1.5pt}]

\titleformat{\section}
  {\normalfont\Large\bfseries\sffamily\color{primarycolor}}
  {\thesection}
  {1em}
  {}
  [\vspace{0.5ex}{\color{secondarycolor}\hrule height 0.5pt}]

\titleformat{\subsection}
  {\normalfont\large\bfseries\sffamily\color{secondarycolor}}
  {\thesubsection}
  {1em}
  {}

\tcbset{
    colback=gray!3!white,
    colframe=primarycolor,
    fonttitle=\bfseries\sffamily,
    arc=2mm,
    boxrule=1.2pt,
    left=12pt,
    right=12pt,
    top=8pt,
    bottom=8pt
}

\begin{document}
\shorthandoff{=}

% 
% KAPAK SAYFASI
% 
\begin{titlepage}
    \centering
    \vspace*{2cm}
    {\huge\bfseries\sffamily Dijital Sürdürülebilirlik}\\[0.5cm]
    {\Large\bfseries Akademik Ders Notları \& Açıklamalı Soru Bankası}\\[1.5cm]
    
    \hrule height 2pt
    \vspace{0.8cm}
    {\large Bu ders notu, dijital dönüşüm süreçleri, dijital olgunluk modelleri,}\\[0.2cm]
    {\large büyük veri stratejileri, yapay zekâ karar fabrikaları, ortak malların trajedisi,}\\[0.2cm]
    {\large sürdürülebilir akıllı şehirler ve kalkınma ilkelerini kapsayan kapsamlı bir kaynak kitapçıktır.}\\[0.8cm]
    \hrule height 2pt
    \vspace{1.5cm}
    
    {\Large\bfseries Hazırlayan: Oğuz TAŞDEMİR}\\[1.5cm]
    
\end{titlepage}

\tableofcontents
\newpage

\chapter{Konu Anlatımı ve Ders Özetleri}

\section{Bölüm 1: Dijital Dönüşüm ve Dijitalleşmeye Giriş}
Dijital Dönüşüm (Digital Transformation), küreselleşme, artan rekabet ve hızla değişen müşteri beklentilerine yanıt olarak, iş süreçlerinin dijital teknolojilerle yeniden tasarlanmasıdır. Bu süreç sadece bir teknoloji entegrasyonu değil; kültür, süreç ve iş modeli değişimidir. 

Tarihsel süreçte toplumların gelişimi devrimler ve krizler doğrultusunda şekillenmiştir:
\begin{itemize}
    \item \textbf{İlkel Toplum Dönemi:} Doğal kaynakların kullanımına bağımlı olan bu dönem, kaynakların azalmasıyla bir krize girmiş ve bu kriz insanlığı toprağı işlemeye, yani Tarım Devrimi'ne yönlendirmiştir.
    \item \textbf{Tarım Toplumu Dönemi:} Yerleşik yaşam ve tarım arazilerinin yönetimi önem kazanmıştır. Arazilerin kıtlığı ve verim arayışı Sanayi Devrimi'ni tetiklemiştir.
    \item \textbf{Sanayi Toplumu Dönemi:} Modern makine yatırımları, seri üretim ve standart (tek tip) üretim bu ekonominin ana unsurları olmuştur.
    \item \textbf{Bilgi Toplumu Dönemi:} Bilgisayar destekli üretim (CAD/CAM) ve iletişim teknolojilerindeki hızlı gelişmelerle şekillenmiştir. Global pazarda rekabet edebilmek için ürün ve hizmetlerin \textbf{toplu olarak kişiselleştirilmesi (mass customization)} ön plana çıkmıştır.
\end{itemize}

Dijital dönüşümün sürdürülebilir bir şekilde hayata geçirilmesi dört temel ilkeyle (\textbf{4C}) tanımlanır: Connection (Bağlantı), Coordination (Koordinasyon), Collaboration (İş Birliği) ve Commerce (Ticaret).

\needspace{9cm}
\section{Bölüm 2: Sayısallaştırma, Dijitalleşme ve Dijital Dönüşüm Tarihçesi}
Dijital dönüşüm yolculuğunun kavramsal temelleri üç farklı aşamada ele alınır:
\begin{itemize}
    \item \textbf{Sayısallaştırma (Digitization):} Fiziksel haldeki bilginin dijital forma dönüştürülmesidir. Örneğin; bir kağıt belgeyi tarayıcıdan geçirerek bilgisayarda PDF olarak saklamak, kağıt notları Excel dosyasına işlemek ya da analog VHS kasetlerini DVD disklerine kaydetmek.
    \item \textbf{Dijitalleşme (Digitalization):} Sayısallaştırılmış bilginin iş süreçlerini iyileştirmek, otomatikleştirmek ve veriden yeni gelir elde etmek amacıyla kullanılmasıdır. Örneğin; PDF belgesini buluta yükleyip analiz için paylaşmak, Excel dosyasını Google E-Tabloya dönüştürmek ya da dijital filmleri çevrimiçi yayınlamak (streaming).
    \item \textbf{Dijital Dönüşüm (Digital Transformation):} Teknolojilerin ve verilerin entegrasyonu ile iş modellerinin, kurum kültürünün ve değer zincirinin kökten değiştirilmesidir.
\end{itemize}

\textbf{Sanayi 4.0 ve Teknolojileri:} Dördüncü sanayi devrimi veri ve bilgiye olan bağımlılığı nedeniyle diğer sanayi devrimlerinden ayrılır. Dalenogare ve ark. (2018) tarafından yapılan çalışmada gelişmekte olan ülkeler için öne çıkan kritik teknolojilerden biri \textbf{Simülasyonlar ve Sanal Modellerin Analizi} olarak belirtilmiştir. Bunun yanı sıra IoT, Büyük Veri Analitiği, MES (Üretim Yönetim Sistemleri) ve SCADA sistemleri bu devrimin yapı taşlarıdır.

\needspace{9cm}
\section{Bölüm 3: Dijital Olgunluk Modelleri}
Olgunluk modelleri, bir kuruluşun yeteneklerinin öngörülen, istenen veya mantıklı bir süreç sonucunda nasıl geliştiğini görmeye yardımcı olan yaklaşımlardır. Büyük çoğunluğu Deming'in Kalite Yönetim Sistemleri içerisinde belirlediği \textbf{PUKÖ (Planla-Uygula-Kontrol Et-Önlem Al)} sürecini esas almaktadır.

Literatürde dijital olgunluk modellerine ilişkin farklı tanımlar yapılmıştır:
\begin{itemize}
    \item \textbf{Paulk, Curtis ve Weber (1993):} Yazılım yeteneklerini artırmak için kuruluşların yazılım geliştirme süreçlerini iyileştiren aşamaları tanımlar.
    \item \textbf{Klimko (2001):} Olgunluk durumunu bir işletmede yer alan bir varlığın (fiziki veya örgütsel işlev) zaman içindeki gelişim süreci olarak açıklar.
    \item \textbf{Pullen (2007):} Olgunluk modelini, gelişim sürecinin farklı aşamalarında etkili işlemlerin özelliklerini tanımlayan yapılandırılmış öğeler koleksiyonu olarak betimler.
\end{itemize}

\textbf{IBM Olgunluk Modeli (Berman ve Bell, 2011):} İşletmelerin dijital olgunluk seviyelerini 6 boyutta analiz eder. Bu boyutlardan bazıları şunlardır:
\begin{enumerate}
    \item \textbf{Analitik Görü:} Kestirimci tahmine ve ileri düzey veri analizine dayalı bilgi üretimi.
    \item \textbf{Süreç Entegrasyonu:} Dijital olanaklar kullanılarak süreçlerin optimize edilmesi ve şirketler arasında entegre edilmesi.
    \item \textbf{İş Modeli İnovasyonu:} Her seviyede müşteri değerinin oluşturulmasını hedefleyen yeni iş modellerinin tasarlanması.
    \item \textbf{Müşteri ve Toplum İş Birliği:} Müşteri odaklılığın tüm departmanlara yayılması ve sosyal ağların sisteme entegre edilmesi.
\end{enumerate}

\needspace{9cm}
\section{Bölüm 4: Türkiye'de Dijitalleşme Çalışmaları}
Türkiye'de sanayide dijital dönüşüme yönelik devlet politikaları 3 ana faaliyet alanı ve birime ayrılmıştır:
\begin{enumerate}
    \item \textbf{Politika Geliştirme Süreci:} Sanayi ve Teknoloji Bakanlığı Ar-Ge Teşvikleri Genel Müdürlüğü tarafından üstlenilmiştir.
    \item \textbf{Pilot Uygulama Süreçleri:} Sanayi ve Verimlilik Genel Müdürlüğüne aittir.
    \item \textbf{Dönüşüm Finansmanı:} AB ve Dış İlişkiler Genel Müdürlüğü, TÜBİTAK ve KOSGEB olmak üzere üç ana kaynak belirlenmiştir.
\end{enumerate}

Sanayide dijital dönüşümün yol haritalarını belirlemek üzere TOBB, TİM, TÜSİAD, MÜSİAD, YASED ve TTGV iş birliği ile politika belgeleri oluşturulmuş ve 6 çalışma grubu kurulmuştur: İleri üretim teknikleri, Eğitim, Dijital teknolojiler, Mevzuat ve standardizasyon, Altyapı, Açık İnovasyon.

Devlet nezdinde atılan en önemli adımlar arasında Cumhurbaşkanlığı Dijital Dönüşüm Ofisi ve Dördüncü Sanayi Devrimi Dairesi'nin kurulması, sanayide dijital olgunluk düzeyi anketlerinin yapılması ve öncelikle Ankara ile Bursa'da olmak üzere \textbf{Dijital Model Fabrikaların} kurulması yer almaktadır.

\needspace{9cm}
\section{Bölüm 5: Dijital Dönüşümde Verinin Önemi ve Büyük Veri}
Dijital çağda veri; inovasyonu tetikleyen, stratejik kararları yönlendiren ve işletmeleri dönüştüren yakıttır. Bu bağlamda \textbf{"Veri yeni petroldür"} yaklaşımı benimsenmiştir. 2025 yılında dünya genelinde üretilecek veri miktarının \textbf{175 ZB (Zettabyte)} seviyesine ulaşacağı öngörülmektedir. Şirketlerde verinin yönetilmesi ve stratejik kararlarda kullanılması amacıyla \textbf{CDO (Chief Data Officer)} rolleri hızla yaygınlaşmaktadır.

\textbf{Büyük Verinin 3V Prensibi:} Büyük veri, sadece miktar fazlalığını değil, aynı zamanda verinin üretim hızını ve çeşitliliğini de tanımlar:
\begin{itemize}
    \item \textbf{Volume (Hacim):} Depolanan verinin büyüklüğü.
    \item \textbf{Velocity (Hız):} Verinin üretilme ve işlenme sıklığı.
    \item \textbf{Variety (Çeşitlilik):} Metin, video, ses, sensör sinyalleri gibi farklı formatlardaki veri yapıları.
\end{itemize}

Büyük veri analitiği; tanımlayıcı, tanısal, öngörücü ve yönlendirici analitik olarak sınıflandırılır. Endüstride büyük verinin kullanımına örnek olarak 1.3 milyon cihazdan veri toplayıp analiz eden \textbf{Siemens MindSphere} platformu ve Arçelik bünyesindeki \textbf{Data Lake} (Veri Gölü) çalışmaları gösterilebilir.

\needspace{9cm}
\section{Bölüm 6: Yapay Zekâ ve Dijital Karar Fabrikası}
Yapay zekâ, geleneksel firmaları sınırlayan operasyonel kısıtlamaları ortadan kaldırarak yeni nesil dijital şirketlerin doğmasını sağlamıştır. Bu yeni şirketlerin merkezinde \textbf{"Yapay Zekâ Fabrikası"} (AI Factory) denilen bir karar mekanizması yer almaktadır.

\textbf{Ant Financial Group (Alibaba / Alipay) Örneği:}
\begin{itemize}
    \item Piyasaya girmesinin üzerinden yalnızca 5 yıl geçtikten sonra (2019'da) 1 milyar müşteri barajını aşmıştır.
    \item ABD'nin en büyük bankalarının 10'da birinden daha az çalışanla, bu bankaların 10 katı kadar müşteriye hizmet vermiştir.
    \item 2018 sonlarında piyasa değeri 150 milyar dolara ulaşmıştır.
    \item Kredi onayları, finansal tavsiyeler veya operasyonel faaliyetlerde hiçbir insan temsilci çalışmamakta, tüm süreç yapay zekâ algoritmalarıyla yürütülmektedir.
\end{itemize}

\textbf{Yapay Zekâ Fabrikasının 4 Temel Bileşeni:}
\begin{enumerate}
    \item Sistematik ve ölçeklenebilir şekilde veri toplayan otomatik bir veri hattı (data pipeline).
    \item Geleceğe yönelik kestirimlerde bulunan algoritmik karar motoru.
    \item Algoritmaların test edildiği bir deney platformu.
    \item Bu süreci yazılıma gömen altyapı ve yazılım entegrasyonu.
\end{enumerate}

Yapay zekâ çalışmaları, karmaşıklığına bağlı olarak \textbf{Zayıf Yapay Zekâ} (sadece tek bir dar alanda uzmanlaşmış algoritmalar) ve \textbf{Güçlü Yapay Zekâ} (değişen koşullara hızla adapte olabilen genel insan zekâsı düzeyi) olarak ikiye ayrılır.

\needspace{9cm}
\section{Bölüm 7: Dijital Dönüşümde Doğru ve Yanlışlar}
Araştırmalar, yöneticilerin dijital dönüşümle ilgili pek çok efsaneye inandığını ancak gerçeklerin farklı olduğunu göstermektedir:
\begin{itemize}
    \item \textbf{Efsane 1: Dijital dönüşüm iş modelinde kökten bir yıkım ve devasa yatırım gerektirir.}
        \\ \textit{Gerçek:} Dönüşüm eskiyi çöpe atmak değil, çekirdek değer önermesini dijital araçlarla daha iyi sunmaktır. Örneğin; \textbf{Maersk} nakliye şeffaflığıyla, \textbf{Aeroflot} operasyonel raporlamayı dijitalleştirip NPS skorunu \%44 artırarak, \textbf{G7 Taksi} ise Uber'e karşı kendi çağırma sistemini kurarak değerini artırmıştır.
    \item \textbf{Efsane 2: Fiziksel kanallar tamamen kalkıp sanal kanallara geçilecektir.}
        \\ \textit{Gerçek:} İkisi bir arada var olacaktır (Hibrit Model). \textbf{Galeries Lafayette Champs-Élysées} fiziksel mağazasında stilistler ile teknolojiyi harmanlayarak müşterilerle duygusal bağ kurmaktadır.
    \item \textbf{Efsane 3: Şirketler yenilikçilik için startup satın almalıdır.}
        \\ \textit{Gerçek:} Startup'ları ana firmanın bürokrasisinde boğmak yerine, yarı bağımsız çalışıp sinerji yaratmalarına izin verilmelidir (Örn: \textbf{Avnet}'in Hackster.io'yu satın alıp bağımsız bırakması).
    \item \textbf{Efsane 4: Dönüşüm sadece teknoloji yatırımıdır.}
        \\ \textit{Gerçek:} Dönüşüm insan kaynağını ve kültürü içerir. \textbf{Adobe}'un çalışanların yenilikçi projelerini ödüllendiren \textbf{"Turuncu Kutu"} projesi buna örnektir.
    \item \textbf{Efsane 5: Eski bilgi sistemleri altyapısını tamamen yenileyerek başlamak gerekir.}
        \\ \textit{Gerçek:} Bu çok risklidir. Akıllı şirketler dönüşümü aşamalı yapar ve pazarda test ederek ilerler (Örn: \textbf{Marriott}'un Expedia rekabetinde aşamalı güncellemesi).
\end{itemize}

\needspace{9cm}
\section{Bölüm 8: Sürdürülebilirlik ve Ortak Malların Trajedisi}
Sürdürülebilirlik, gelecek nesillerin kendi ihtiyaçlarını karşılama yeteneğinden ödün vermeden bugünün ihtiyaçlarını karşılamaktır. Çevresel, sosyal ve ekonomik sürdürülebilirlik olmak üzere üç ana boyutu vardır.

\textbf{Ortak Malların Trajedisi (Tragedy of the Commons):} Garrett Hardin tarafından tanımlanan bu kurama göre, bireyler kamuya açık ve ücretsiz olan ortak kaynaklardan (örneğin meralar) kendi faydalarını maksimize etmek için aşırı tüketim yaparlar. Bu durum kaynağın yok olmasına ve uzun vadede herkesin kaybetmesine yol açar (Örn: çobanın merada hayvan sayısını sınırsızca artırması).

\textbf{Kurumsal Sosyal Sorumluluk ve Çevresel Felaketler:} Geleneksel kâr odaklı yönetim anlayışının çevreye zararları büyüktür. 1996'da Nike'ın Pakistan'da çocuk işçi çalıştırdığının ortaya çıkması ve 3 Aralık 1984'te Hindistan Bhopal'deki \textbf{Union Carbide} böcek ilacı fabrikasından sızan \textbf{Metil İzosiyanat} gazının binlerce insanı öldürmesi tarihteki en büyük sorumsuzluk örnekleridir.

\textbf{Yeşil Tedarik Zinciri ve Tersine Lojistik ve Enerji Tasarrufu:} Tersine lojistik, ürünlerin geri kazanımı veya imhası için tüketimden orijine doğru akışıdır. Yeniden üretim süreçleri, sıfırdan üretime göre \textbf{\%85'e varan oranda enerji tasarrufu} sağlar. Geri dönüştürülmüş alüminyum üretmek sıfırdan üretimin sadece \%5'i kadar enerji gerektirir. Lojistikte verimliliği artırmak için 3M, rafların yüksekliğini ayarlayarak araç doluluk oranını \%40 artırmıştır.

\needspace{9cm}
\section{Bölüm 9: Sürdürülebilir Akıllı Şehirler}
Sürdürülebilir akıllı şehirler; bilgi ve iletişim teknolojilerini kullanarak kısıtlı kaynakları daha etkin yöneten, karbon ayak izini azaltan ve vatandaşın yaşam kalitesini artıran kent modelleridir.

Şehirlerin büyümesiyle birlikte krizler de artmaktadır:
\begin{itemize}
    \item Dünyada üretilen toplam enerjinin \textbf{\%75'i şehirlerde tüketilmekte} ve küresel CO2 salınımının \%80'i şehirlerden kaynaklanmaktadır.
    \item Günümüzde dünya nüfusunun \%55'i şehirlerde yaşarken, 2050 yılına kadar bu oranın \textbf{\%68'e ulaşacağı} öngörülmektedir.
\end{itemize}

Kentsel sürdürülebilirlik ekolojik ve çevresel faktörleri korumaya odaklanırken (Neirotti vd. 2014); akıllı şehir kavramı daha çok sosyal, ekonomik ve veri odaklı teknolojik süreçleri içerir. İkisinin birleşimi olan sürdürülebilir akıllı şehirleşmenin temel esası, yerel düzeyde kalkınmayı ve toplumsal dayanışmayı bilgi teknolojileriyle sağlamaktır.

\needspace{9cm}
\section{Bölüm 10: Sürdürülebilir Kalkınma ve Turizm İlkeleri}
Sürdürülebilir kalkınmanın temel hedefi, ekonomik büyümeyi doğal kaynakların yenilenebilme kapasiteleriyle uyumlu şekilde gerçekleştirmektir.

Burns ve Holden (1995) çalışmasına göre sürdürülebilir kalkınmanın bazı ekonomik unsurları: ekonomik faaliyetlerin yeniden dağıtımı ve ikame olanakları olmayan yenilenemeyen kaynakların kullanımına sınırlama getirilmesidir.

BM Çevre Programı (UNEP) ve BM Dünya Turizm Örgütü (UNWTO) sürdürülebilir turizmin 12 temel ilkesini belirlemiştir:
\begin{itemize}
    \item \textbf{Ekonomik Süreklilik ve Yerel Kalkınma:} Turizm gelirlerinin yerel üreticilerle buluşması ve destinasyona katkı sağlaması.
    \item \textbf{İstihdam Kalitesi ve Sosyal Katılım:} Dezavantajlı gruplara fırsatlar sunulması ve istihdam kalitesinin artırılması.
    \item \textbf{Kültürel Zenginlik:} Özgün kültüre, tarihi mirasa ve ayırt edici yerel değerlere saygı duyulması.
    \item \textbf{Biyolojik Çeşitlilik ve Çevresel Etki:} Koruma-kullanma dengesinin gözetilmesi ve atık üretiminin asgari düzeye indirilmesi.
\end{itemize}

\textbf{Kültürel Sürdürülebilirlik:} Kültürel sürdürülebilirlik gelenekçi ve donmuş bir tavır değildir; aksine ülkelerin geçmiş mirasının ve yerleşim koşullarının çok büyük birer kaynak olduğunun fark edilerek yeni kuşaklara aktarılması sürecidir.

\newpage
\chapter{Kelime Kartları ve Hızlı Tekrar}

\needspace{3.5cm}
\begin{tcolorbox}[colback=white, colframe=secondarycolor!50!black, arc=1.2mm, boxrule=0.8pt, top=8pt, bottom=8pt, left=12pt, right=12pt, segmentation style={dashed, secondarycolor!30, line width=0.8pt}]
    \noindent\textbf{\color{red!80!black}\sffamily Soru 1:} Dijital Dönüşüm nedir?
\tcblower
    \noindent\textbf{\color{green!50!black}\sffamily Cevap:} İş süreçlerinin, kültürün ve iş modellerinin dijital teknolojilerle kökten yeniden tasarlanmasıdir.
\end{tcolorbox}
\vspace{6pt}

\needspace{3.5cm}
\begin{tcolorbox}[colback=white, colframe=secondarycolor!50!black, arc=1.2mm, boxrule=0.8pt, top=8pt, bottom=8pt, left=12pt, right=12pt, segmentation style={dashed, secondarycolor!30, line width=0.8pt}]
    \noindent\textbf{\color{red!80!black}\sffamily Soru 2:} Sayısallaştırma (Digitization) nedir?
\tcblower
    \noindent\textbf{\color{green!50!black}\sffamily Cevap:} Fiziksel bilginin dijital forma dönüştürülmesidir (Örn: Kağıdı tarayıp PDF yapmak).
\end{tcolorbox}
\vspace{6pt}

\needspace{3.5cm}
\begin{tcolorbox}[colback=white, colframe=secondarycolor!50!black, arc=1.2mm, boxrule=0.8pt, top=8pt, bottom=8pt, left=12pt, right=12pt, segmentation style={dashed, secondarycolor!30, line width=0.8pt}]
    \noindent\textbf{\color{red!80!black}\sffamily Soru 3:} Dijitalleşme (Digitalization) nedir?
\tcblower
    \noindent\textbf{\color{green!50!black}\sffamily Cevap:} Sayısallaştırılmış bilginin iş süreçlerini iyileştirmek ve otomatikleştirmek için kullanılmasıdır.
\end{tcolorbox}
\vspace{6pt}

\needspace{3.5cm}
\begin{tcolorbox}[colback=white, colframe=secondarycolor!50!black, arc=1.2mm, boxrule=0.8pt, top=8pt, bottom=8pt, left=12pt, right=12pt, segmentation style={dashed, secondarycolor!30, line width=0.8pt}]
    \noindent\textbf{\color{red!80!black}\sffamily Soru 4:} Yaratıcı Yıkım (Creative Destruction) nedir?
\tcblower
    \noindent\textbf{\color{green!50!black}\sffamily Cevap:} Yeni teknolojilerin eski ekonomik yapıları ortadan kaldırarak daha verimli yenilerini kurmasıdir.
\end{tcolorbox}
\vspace{6pt}

\needspace{3.5cm}
\begin{tcolorbox}[colback=white, colframe=secondarycolor!50!black, arc=1.2mm, boxrule=0.8pt, top=8pt, bottom=8pt, left=12pt, right=12pt, segmentation style={dashed, secondarycolor!30, line width=0.8pt}]
    \noindent\textbf{\color{red!80!black}\sffamily Soru 5:} Toplu Kişiselleştirme (Mass Customization) nedir?
\tcblower
    \noindent\textbf{\color{green!50!black}\sffamily Cevap:} Ürün ve hizmetlerin esnek teknolojilerle kişisel tercihlere göre seri üretilebilmesidir.
\end{tcolorbox}
\vspace{6pt}

\needspace{3.5cm}
\begin{tcolorbox}[colback=white, colframe=secondarycolor!50!black, arc=1.2mm, boxrule=0.8pt, top=8pt, bottom=8pt, left=12pt, right=12pt, segmentation style={dashed, secondarycolor!30, line width=0.8pt}]
    \noindent\textbf{\color{red!80!black}\sffamily Soru 6:} Sanayi 4.0 nedir?
\tcblower
    \noindent\textbf{\color{green!50!black}\sffamily Cevap:} Veri ve bilgiye bağımlı olan, siber-fiziksel sistemler üzerine kurulu dördüncü sanayi devrimidir.
\end{tcolorbox}
\vspace{6pt}

\needspace{3.5cm}
\begin{tcolorbox}[colback=white, colframe=secondarycolor!50!black, arc=1.2mm, boxrule=0.8pt, top=8pt, bottom=8pt, left=12pt, right=12pt, segmentation style={dashed, secondarycolor!30, line width=0.8pt}]
    \noindent\textbf{\color{red!80!black}\sffamily Soru 7:} Olgunluk Modeli nedir?
\tcblower
    \noindent\textbf{\color{green!50!black}\sffamily Cevap:} Kuruluş yeteneklerinin öngörülen süreçlerle zaman içinde nasıl geliştiğini izleyen yaklaşımlardır.
\end{tcolorbox}
\vspace{6pt}

\needspace{3.5cm}
\begin{tcolorbox}[colback=white, colframe=secondarycolor!50!black, arc=1.2mm, boxrule=0.8pt, top=8pt, bottom=8pt, left=12pt, right=12pt, segmentation style={dashed, secondarycolor!30, line width=0.8pt}]
    \noindent\textbf{\color{red!80!black}\sffamily Soru 8:} Deming'in PUKÖ döngüsü nedir?
\tcblower
    \noindent\textbf{\color{green!50!black}\sffamily Cevap:} Planla, Uygula, Kontrol Et ve Önlem Al adımlarından oluşan sürekli kalite yönetimi döngüsüdür.
\end{tcolorbox}
\vspace{6pt}

\needspace{3.5cm}
\begin{tcolorbox}[colback=white, colframe=secondarycolor!50!black, arc=1.2mm, boxrule=0.8pt, top=8pt, bottom=8pt, left=12pt, right=12pt, segmentation style={dashed, secondarycolor!30, line width=0.8pt}]
    \noindent\textbf{\color{red!80!black}\sffamily Soru 9:} Büyük Verinin 3V prensibi nedir?
\tcblower
    \noindent\textbf{\color{green!50!black}\sffamily Cevap:} Volume (Hacim), Velocity (Hız) ve Variety (Çeşitlilik) bileşenleridir.
\end{tcolorbox}
\vspace{6pt}

\needspace{3.5cm}
\begin{tcolorbox}[colback=white, colframe=secondarycolor!50!black, arc=1.2mm, boxrule=0.8pt, top=8pt, bottom=8pt, left=12pt, right=12pt, segmentation style={dashed, secondarycolor!30, line width=0.8pt}]
    \noindent\textbf{\color{red!80!black}\sffamily Soru 10:} Yapay Zeka Fabrikası (AI Factory) nedir?
\tcblower
    \noindent\textbf{\color{green!50!black}\sffamily Cevap:} Karar verme sürecini algoritmalarla bilimsel ve sistematik hale getiren yazılımsal yapay zeka merkezidir.
\end{tcolorbox}
\vspace{6pt}

\needspace{3.5cm}
\begin{tcolorbox}[colback=white, colframe=secondarycolor!50!black, arc=1.2mm, boxrule=0.8pt, top=8pt, bottom=8pt, left=12pt, right=12pt, segmentation style={dashed, secondarycolor!30, line width=0.8pt}]
    \noindent\textbf{\color{red!80!black}\sffamily Soru 11:} Zayıf Yapay Zeka nedir?
\tcblower
    \noindent\textbf{\color{green!50!black}\sffamily Cevap:} Tek bir dar görev alanında uzmanlaşmış, etki alanı kısıtlı olan yapay zekadır.
\end{tcolorbox}
\vspace{6pt}

\needspace{3.5cm}
\begin{tcolorbox}[colback=white, colframe=secondarycolor!50!black, arc=1.2mm, boxrule=0.8pt, top=8pt, bottom=8pt, left=12pt, right=12pt, segmentation style={dashed, secondarycolor!30, line width=0.8pt}]
    \noindent\textbf{\color{red!80!black}\sffamily Soru 12:} Güçlü Yapay Zeka nedir?
\tcblower
    \noindent\textbf{\color{green!50!black}\sffamily Cevap:} Etki alanı sınırlarıyla kısıtlanmamış, değişen koşullara adapte olabilen genel insan zekâsı düzeyidir.
\end{tcolorbox}
\vspace{6pt}

\needspace{3.5cm}
\begin{tcolorbox}[colback=white, colframe=secondarycolor!50!black, arc=1.2mm, boxrule=0.8pt, top=8pt, bottom=8pt, left=12pt, right=12pt, segmentation style={dashed, secondarycolor!30, line width=0.8pt}]
    \noindent\textbf{\color{red!80!black}\sffamily Soru 13:} Sürdürülebilirlik nedir?
\tcblower
    \noindent\textbf{\color{green!50!black}\sffamily Cevap:} Gelecek nesillerin kendi ihtiyaçlarını karşılama yeteneğinden ödün vermeden bugünün ihtiyaçlarını karşılayabilmektir.
\end{tcolorbox}
\vspace{6pt}

\needspace{3.5cm}
\begin{tcolorbox}[colback=white, colframe=secondarycolor!50!black, arc=1.2mm, boxrule=0.8pt, top=8pt, bottom=8pt, left=12pt, right=12pt, segmentation style={dashed, secondarycolor!30, line width=0.8pt}]
    \noindent\textbf{\color{red!80!black}\sffamily Soru 14:} Ortak Malların Trajedisi nedir?
\tcblower
    \noindent\textbf{\color{green!50!black}\sffamily Cevap:} Bireylerin serbest ortak kaynakları (meraları vb.) kendi çıkarları için aşırı kullanarak yok etmesidir.
\end{tcolorbox}
\vspace{6pt}

\needspace{3.5cm}
\begin{tcolorbox}[colback=white, colframe=secondarycolor!50!black, arc=1.2mm, boxrule=0.8pt, top=8pt, bottom=8pt, left=12pt, right=12pt, segmentation style={dashed, secondarycolor!30, line width=0.8pt}]
    \noindent\textbf{\color{red!80!black}\sffamily Soru 15:} Tersine Lojistik (Reverse Logistics) nedir?
\tcblower
    \noindent\textbf{\color{green!50!black}\sffamily Cevap:} Ürünlerin değer kazanımı veya imhası için tüketim noktasından çıkış noktasına (orijine) doğru akışıdır.
\end{tcolorbox}
\vspace{6pt}

\needspace{3.5cm}
\begin{tcolorbox}[colback=white, colframe=secondarycolor!50!black, arc=1.2mm, boxrule=0.8pt, top=8pt, bottom=8pt, left=12pt, right=12pt, segmentation style={dashed, secondarycolor!30, line width=0.8pt}]
    \noindent\textbf{\color{red!80!black}\sffamily Soru 16:} Yeşil Tedarik Zinciri Yönetimi nedir?
\tcblower
    \noindent\textbf{\color{green!50!black}\sffamily Cevap:} Satın almadan paketlemeye kadar tüm tedarik zinciri adımlarına çevresel faktörlerin dahil edilmesidir.
\end{tcolorbox}
\vspace{6pt}

\needspace{3.5cm}
\begin{tcolorbox}[colback=white, colframe=secondarycolor!50!black, arc=1.2mm, boxrule=0.8pt, top=8pt, bottom=8pt, left=12pt, right=12pt, segmentation style={dashed, secondarycolor!30, line width=0.8pt}]
    \noindent\textbf{\color{red!80!black}\sffamily Soru 17:} Sürdürülebilir Akıllı Şehir nedir?
\tcblower
    \noindent\textbf{\color{green!50!black}\sffamily Cevap:} Bilgi teknolojilerini kullanarak kısıtlı kaynakları etkin yöneten ve yaşam kalitesini artıran kent modelidir.
\end{tcolorbox}
\vspace{6pt}

\needspace{3.5cm}
\begin{tcolorbox}[colback=white, colframe=secondarycolor!50!black, arc=1.2mm, boxrule=0.8pt, top=8pt, bottom=8pt, left=12pt, right=12pt, segmentation style={dashed, secondarycolor!30, line width=0.8pt}]
    \noindent\textbf{\color{red!80!black}\sffamily Soru 18:} Kentsel Sürdürülebilirlik kentsel alanlarda hangi boyutu öncelikli olarak ele alır?
\tcblower
    \noindent\textbf{\color{green!50!black}\sffamily Cevap:} Ekolojik dengeyi ve kirliliği önlemeyi amaçlayan çevre boyutunu ele alır.
\end{tcolorbox}
\vspace{6pt}

\needspace{3.5cm}
\begin{tcolorbox}[colback=white, colframe=secondarycolor!50!black, arc=1.2mm, boxrule=0.8pt, top=8pt, bottom=8pt, left=12pt, right=12pt, segmentation style={dashed, secondarycolor!30, line width=0.8pt}]
    \noindent\textbf{\color{red!80!black}\sffamily Soru 19:} Sürdürülebilir Kalkınma nedir?
\tcblower
    \noindent\textbf{\color{green!50!black}\sffamily Cevap:} Ekonomik büyümeyi ekolojik dengelerle ve doğal kaynakların yenilenebilme kapasiteleriyle uyumlu yapmaktır.
\end{tcolorbox}
\vspace{6pt}

\needspace{3.5cm}
\begin{tcolorbox}[colback=white, colframe=secondarycolor!50!black, arc=1.2mm, boxrule=0.8pt, top=8pt, bottom=8pt, left=12pt, right=12pt, segmentation style={dashed, secondarycolor!30, line width=0.8pt}]
    \noindent\textbf{\color{red!80!black}\sffamily Soru 20:} Kültürel Sürdürülebilirlik nedir?
\tcblower
    \noindent\textbf{\color{green!50!black}\sffamily Cevap:} Bir toplumun geçmiş mirasını ve kültürel değerlerini koruyarak gelecek kuşaklara aktarabilmesidir.
\end{tcolorbox}
\vspace{6pt}

\newpage
\chapter{Konu Sonu Testleri}
\section{Dijital Sürdürülebilirlik Açıklamalı Soru Bankası}

\needspace{5cm}
\vspace{8pt}
\noindent\textbf{\large\sffamily Soru 1.} \ Joseph Schumpeter tarafından ortaya atılan, yeni teknolojilerin eski ekonomik yapıları yıkarak yerlerine daha verimli sistemler getirmesini ifade eden kavram hangisidir?
\begin{enumerate}[label=\Alph*), leftmargin=25pt, topsep=4pt, itemsep=2pt]
    \item Dijital Olgunluk
    \item Yaratıcı Yıkım (Creative Destruction)
    \item Tersine Lojistik
    \item Ortak Malların Trajedisi
\end{enumerate}
\vspace{4pt}
\begin{tcolorbox}[title=Cevap ve Açıklama, colback=gray!3!white, colframe=primarycolor!90!black, arc=2mm, fonttitle=\bfseries\sffamily, top=6pt, bottom=6pt]
    \textbf{Doğru Cevap: B}
    \vspace{2pt}
    Schumpeter'in 'yaratıcı yıkım' (creative destruction) kavramı, inovasyon odaklı rekabetin ve dijitalleşmenin eski şirket yeteneklerini ve pazar yapılarını ortadan kaldırarak yeni ve verimli modeller getirmesini ifade eder.
\end{tcolorbox}
\vspace{10pt}

\needspace{5cm}
\vspace{8pt}
\noindent\textbf{\large\sffamily Soru 2.} \ Aşağıdakilerden hangisi Tarım Toplumu döneminde ortaya çıkan krizin ve ardından yaşanan evrimin temel sebebidir?
\begin{enumerate}[label=\Alph*), leftmargin=25pt, topsep=4pt, itemsep=2pt]
    \item Tarım arazilerinin aşırı bolluğu ve iş gücü eksikliği
    \item İlkel toplumun doğal kaynakları tüketmesi sonucu yaşanan kriz
    \item Bilgisayar destekli üretim teknolojilerinin yetersiz kalması
    \item Küreselleşmenin getirdiği standart üretim ihtiyacı
\end{enumerate}
\vspace{4pt}
\begin{tcolorbox}[title=Cevap ve Açıklama, colback=gray!3!white, colframe=primarycolor!90!black, arc=2mm, fonttitle=\bfseries\sffamily, top=6pt, bottom=6pt]
    \textbf{Doğru Cevap: B}
    \vspace{2pt}
    İlkel toplum doğal kaynakların kullanımına bağımlıydı. Doğal kaynaklardaki azalma (kıtlık) krize neden olmuş, bu da insanları toprağı ekip biçmeye (Tarım Devrimi'ne) yönlendirmiştir.
\end{tcolorbox}
\vspace{10pt}

\needspace{5cm}
\vspace{8pt}
\noindent\textbf{\large\sffamily Soru 3.} \ Ürün ve hizmetlerin 'toplu olarak kişiselleştirilmesi' (mass customization) hangi toplum evresinin ana ekonomik özelliklerinden biridir?
\begin{enumerate}[label=\Alph*), leftmargin=25pt, topsep=4pt, itemsep=2pt]
    \item Bilgi Toplumu
    \item Sanayi Toplumu
    \item Tarım Toplumu
    \item İlkel Toplum
\end{enumerate}
\vspace{4pt}
\begin{tcolorbox}[title=Cevap ve Açıklama, colback=gray!3!white, colframe=primarycolor!90!black, arc=2mm, fonttitle=\bfseries\sffamily, top=6pt, bottom=6pt]
    \textbf{Doğru Cevap: A}
    \vspace{2pt}
    Bilgi Toplumu döneminde, CAD/CAM gibi esnek üretim ve dijital teknolojiler sayesinde ürün ve hizmetlerin kişiye özel (mass customization) üretilmesi mümkün hale gelmiştir.
\end{tcolorbox}
\vspace{10pt}

\needspace{5cm}
\vspace{8pt}
\noindent\textbf{\large\sffamily Soru 4.} \ Aşağıdakilerden hangisi dijital dönüşümü tanımlayan dört temel ilkeden (4C) biri değildir?
\begin{enumerate}[label=\Alph*), leftmargin=25pt, topsep=4pt, itemsep=2pt]
    \item Connection (Bağlantı)
    \item Collaboration (İş Birliği)
    \item Calculation (Hesaplama)
    \item Coordination (Koordinasyon)
\end{enumerate}
\vspace{4pt}
\begin{tcolorbox}[title=Cevap ve Açıklama, colback=gray!3!white, colframe=primarycolor!90!black, arc=2mm, fonttitle=\bfseries\sffamily, top=6pt, bottom=6pt]
    \textbf{Doğru Cevap: C}
    \vspace{2pt}
    Dijital dönüşümün 4C'si: Connection (Bağlantı), Coordination (Koordinasyon), Collaboration (İş Birliği) ve Commerce (Ticaret) olarak tanımlanır. Calculation bunlardan biri değildir.
\end{tcolorbox}
\vspace{10pt}

\needspace{5cm}
\vspace{8pt}
\noindent\textbf{\large\sffamily Soru 5.} \ Fiziksel haldeki bir belgenin taranarak bilgisayarda PDF olarak saklanması süreci kavramsal olarak aşağıdakilerden hangisidir?
\begin{enumerate}[label=\Alph*), leftmargin=25pt, topsep=4pt, itemsep=2pt]
    \item Dijitalleşme (Digitalization)
    \item Dijital Dönüşüm (Digital Transformation)
    \item Sayısallaştırma (Digitization)
    \item Sanal Gerçeklik (Virtualization)
\end{enumerate}
\vspace{4pt}
\begin{tcolorbox}[title=Cevap ve Açıklama, colback=gray!3!white, colframe=primarycolor!90!black, arc=2mm, fonttitle=\bfseries\sffamily, top=6pt, bottom=6pt]
    \textbf{Doğru Cevap: C}
    \vspace{2pt}
    Sayısallaştırma (digitization), analog veya fiziksel haldeki verinin/bilginin dijital bir formata dönüştürülüp bilgisayar sistemleri tarafından kullanılabilir hale getirilmesi sürecidir.
\end{tcolorbox}
\vspace{10pt}

\needspace{5cm}
\vspace{8pt}
\noindent\textbf{\large\sffamily Soru 6.} \ Sayısallaştırılmış bilginin iş süreçlerini iyileştirmek, otomatikleştirmek ve veri üzerinden gelir yaratmak amacıyla kullanılması süreci hangisidir?
\begin{enumerate}[label=\Alph*), leftmargin=25pt, topsep=4pt, itemsep=2pt]
    \item Sayısallaştırma (Digitization)
    \item Dijitalleşme (Digitalization)
    \item Yapay Zeka Entegrasyonu
    \item Sanayi 3.0
\end{enumerate}
\vspace{4pt}
\begin{tcolorbox}[title=Cevap ve Açıklama, colback=gray!3!white, colframe=primarycolor!90!black, arc=2mm, fonttitle=\bfseries\sffamily, top=6pt, bottom=6pt]
    \textbf{Doğru Cevap: B}
    \vspace{2pt}
    Dijitalleşme (digitalization), sayısallaştırılmış bilginin iş süreçlerini kolaylaştırmak ve gelir sağlamak amacıyla dijital araçlarla kullanılmasıdır.
\end{tcolorbox}
\vspace{10pt}

\needspace{5cm}
\vspace{8pt}
\noindent\textbf{\large\sffamily Soru 7.} \ Dalenogare vd. (2018) çalışmasında yer alan ve gelişmekte olan ülkelerde Sanayi 4.0 kapsamında değerlendirilen beş temel teknoloji başlığından biri hangisidir?
\begin{enumerate}[label=\Alph*), leftmargin=25pt, topsep=4pt, itemsep=2pt]
    \item Geleneksel tarım aletleri
    \item Analog veri saklama sistemleri
    \item Simülasyonlar ve sanal modellerin analizi
    \item Kağıt belgelerin manuel arşivlenmesi
\end{enumerate}
\vspace{4pt}
\begin{tcolorbox}[title=Cevap ve Açıklama, colback=gray!3!white, colframe=primarycolor!90!black, arc=2mm, fonttitle=\bfseries\sffamily, top=6pt, bottom=6pt]
    \textbf{Doğru Cevap: C}
    \vspace{2pt}
    Dalenogare ve ark. çalışmasına göre gelişmekte olan ülkeler için öne çıkan kritik Sanayi 4.0 teknolojilerinden biri 'Simülasyonlar/sanal modellerin analizi'dir.
\end{tcolorbox}
\vspace{10pt}

\needspace{5cm}
\vspace{8pt}
\noindent\textbf{\large\sffamily Soru 8.} \ Kuruluşların yeteneklerinin öngörülen, istenen veya mantıklı bir süreç sonucunda nasıl geliştiğini görmeye yardımcı olan yaklaşımlara ne ad verilir?
\begin{enumerate}[label=\Alph*), leftmargin=25pt, topsep=4pt, itemsep=2pt]
    \item Üretim Modelleri
    \item Olgunluk Modelleri (Maturity Models)
    \item Lojistik Planlama
    \item Karar Ağaçları
\end{enumerate}
\vspace{4pt}
\begin{tcolorbox}[title=Cevap ve Açıklama, colback=gray!3!white, colframe=primarycolor!90!black, arc=2mm, fonttitle=\bfseries\sffamily, top=6pt, bottom=6pt]
    \textbf{Doğru Cevap: B}
    \vspace{2pt}
    Olgunluk modelleri, örgütsel yeteneklerin aşamalı gelişimini ve dijital olgunluk düzeyini ölçen sistematik yapılardır.
\end{tcolorbox}
\vspace{10pt}

\needspace{5cm}
\vspace{8pt}
\noindent\textbf{\large\sffamily Soru 9.} \ Çoğu olgunluk modelinin temelini oluşturan, Deming tarafından geliştirilmiş PUKÖ döngüsünün açılımı nedir?
\begin{enumerate}[label=\Alph*), leftmargin=25pt, topsep=4pt, itemsep=2pt]
    \item Programla-Uygula-Kriter Belirle-Önle
    \item Planla-Uygula-Kontrol Et-Önlem Al
    \item Paylaş-Üret-Kazan-Ölç
    \item Planla-Uyar-Katıl-Öğren
\end{enumerate}
\vspace{4pt}
\begin{tcolorbox}[title=Cevap ve Açıklama, colback=gray!3!white, colframe=primarycolor!90!black, arc=2mm, fonttitle=\bfseries\sffamily, top=6pt, bottom=6pt]
    \textbf{Doğru Cevap: B}
    \vspace{2pt}
    PUKÖ: Planla, Uygula, Kontrol Et ve Önlem Al (PDCA - Plan, Do, Check, Act) aşamalarından oluşan sürekli kalite yönetimi döngüsüdür.
\end{tcolorbox}
\vspace{10pt}

\needspace{5cm}
\vspace{8pt}
\noindent\textbf{\large\sffamily Soru 10.} \ Kurumsal olgunluk gelişimini bir işletmede yer alan bir varlığın (fiziki veya örgütsel işlev) zaman içindeki gelişim süreci üzerinden tanımlayan yazar kimdir?
\begin{enumerate}[label=\Alph*), leftmargin=25pt, topsep=4pt, itemsep=2pt]
    \item Pullen (2007)
    \item Klimko (2001)
    \item Paulk (1993)
    \item Deming (2018)
\end{enumerate}
\vspace{4pt}
\begin{tcolorbox}[title=Cevap ve Açıklama, colback=gray!3!white, colframe=primarycolor!90!black, arc=2mm, fonttitle=\bfseries\sffamily, top=6pt, bottom=6pt]
    \textbf{Doğru Cevap: B}
    \vspace{2pt}
    Klimko (2001), olgunluğu işletmedeki bir varlığın veya fonksiyonel işlevin zaman içerisindeki gelişim süreci olarak betimlemiştir.
\end{tcolorbox}
\vspace{10pt}

\needspace{5cm}
\vspace{8pt}
\noindent\textbf{\large\sffamily Soru 11.} \ IBM Dijital Olgunluk Modeli'nde yer alan, kestirimci tahmine ve ileri düzey veri analizine dayalı bilgi üretimini temsil eden bileşen hangisidir?
\begin{enumerate}[label=\Alph*), leftmargin=25pt, topsep=4pt, itemsep=2pt]
    \item Süreç Entegrasyonu
    \item Analitik Görü (Analytical Insight)
    \item İş Modeli İnovasyonu
    \item Müşteri ve Toplum İş Birliği
\end{enumerate}
\vspace{4pt}
\begin{tcolorbox}[title=Cevap ve Açıklama, colback=gray!3!white, colframe=primarycolor!90!black, arc=2mm, fonttitle=\bfseries\sffamily, top=6pt, bottom=6pt]
    \textbf{Doğru Cevap: B}
    \vspace{2pt}
    IBM olgunluk modelindeki 'Analitik Görü' (Analytical Insight) boyutu, veriyi tahminleme ve ileri veri analitiği teknikleriyle karar almaya dönüştürmeyi ifade eder.
\end{tcolorbox}
\vspace{10pt}

\needspace{5cm}
\vspace{8pt}
\noindent\textbf{\large\sffamily Soru 12.} \ Türkiye'deki dijitalleşme çalışmalarında, politika geliştirme süreci hangi birim tarafından üstlenilmiştir?
\begin{enumerate}[label=\Alph*), leftmargin=25pt, topsep=4pt, itemsep=2pt]
    \item KOSGEB
    \item Sanayi ve Verimlilik Genel Müdürlüğü
    \item Ar-Ge Teşvikleri Genel Müdürlüğü
    \item TÜBİTAK
\end{enumerate}
\vspace{4pt}
\begin{tcolorbox}[title=Cevap ve Açıklama, colback=gray!3!white, colframe=primarycolor!90!black, arc=2mm, fonttitle=\bfseries\sffamily, top=6pt, bottom=6pt]
    \textbf{Doğru Cevap: C}
    \vspace{2pt}
    Devlet politikasında Politika Geliştirme görevi Sanayi ve Teknoloji Bakanlığı Ar-Ge Teşvikleri Genel Müdürlüğü'ne verilmiştir.
\end{tcolorbox}
\vspace{10pt}

\needspace{5cm}
\vspace{8pt}
\noindent\textbf{\large\sffamily Soru 13.} \ Türkiye'de dijitalleşme dönüşüm sürecinde ihtiyaç duyulan finansman için belirlenen üç kurum aşağıdakilerden hangisidir?
\begin{enumerate}[label=\Alph*), leftmargin=25pt, topsep=4pt, itemsep=2pt]
    \item AB ve Dış İlişkiler Gn. Md. - TÜBİTAK - KOSGEB
    \item TOBB - TİM - TÜSİAD
    \item Merkez Bankası - Borsa İstanbul - KOSGEB
    \item Sanayi ve Verimlilik Gn. Md. - MÜSİAD - YASED
\end{enumerate}
\vspace{4pt}
\begin{tcolorbox}[title=Cevap ve Açıklama, colback=gray!3!white, colframe=primarycolor!90!black, arc=2mm, fonttitle=\bfseries\sffamily, top=6pt, bottom=6pt]
    \textbf{Doğru Cevap: A}
    \vspace{2pt}
    Dönüşüm finansmanı için belirlenen resmi üç kaynak: AB ve Dış İlişkiler Genel Müdürlüğü, TÜBİTAK ve KOSGEB'dir.
\end{tcolorbox}
\vspace{10pt}

\needspace{5cm}
\vspace{8pt}
\noindent\textbf{\large\sffamily Soru 14.} \ Türkiye'de sanayinin dijital dönüşümü için kurulan ve ilk uygulamaları Ankara ve Bursa'da gerçekleştirilen merkezler hangileridir?
\begin{enumerate}[label=\Alph*), leftmargin=25pt, topsep=4pt, itemsep=2pt]
    \item Teknoparklar
    \item Dijital Model Fabrikalar
    \item Ar-Ge Merkezleri
    \item KOSGEB Laboratuvarları
\end{enumerate}
\vspace{4pt}
\begin{tcolorbox}[title=Cevap ve Açıklama, colback=gray!3!white, colframe=primarycolor!90!black, arc=2mm, fonttitle=\bfseries\sffamily, top=6pt, bottom=6pt]
    \textbf{Doğru Cevap: B}
    \vspace{2pt}
    Sanayide yalın üretim ve dijital dönüşümü öğretmek üzere Ankara ve Bursa başta olmak üzere 'Dijital Model Fabrikalar' kurulmuştur.
\end{tcolorbox}
\vspace{10pt}

\needspace{5cm}
\vspace{8pt}
\noindent\textbf{\large\sffamily Soru 15.} \ 2025 yılında dünya genelinde üretileceği öngörülen toplam veri miktarı kaç ZettaByte (ZB) seviyesindedir?
\begin{enumerate}[label=\Alph*), leftmargin=25pt, topsep=4pt, itemsep=2pt]
    \item 175 ZB
    \item 59 ZB
    \item 1000 ZB
    \item 10 ZB
\end{enumerate}
\vspace{4pt}
\begin{tcolorbox}[title=Cevap ve Açıklama, colback=gray!3!white, colframe=primarycolor!90!black, arc=2mm, fonttitle=\bfseries\sffamily, top=6pt, bottom=6pt]
    \textbf{Doğru Cevap: A}
    \vspace{2pt}
    Veri stratejisi belgelerinde 2025 yılı için öngörülen küresel veri hacmi 175 ZB (Zettabyte) seviyesindedir.
\end{tcolorbox}
\vspace{10pt}

\needspace{5cm}
\vspace{8pt}
\noindent\textbf{\large\sffamily Soru 16.} \ Şirketlerde verinin stratejik bir varlık olarak konumlandırılmasını yöneten üst düzey yönetici rolü hangisidir?
\begin{enumerate}[label=\Alph*), leftmargin=25pt, topsep=4pt, itemsep=2pt]
    \item CEO (Chief Executive Officer)
    \item CDO (Chief Data Officer)
    \item CFO (Chief Financial Officer)
    \item CTO (Chief Technology Officer)
\end{enumerate}
\vspace{4pt}
\begin{tcolorbox}[title=Cevap ve Açıklama, colback=gray!3!white, colframe=primarycolor!90!black, arc=2mm, fonttitle=\bfseries\sffamily, top=6pt, bottom=6pt]
    \textbf{Doğru Cevap: B}
    \vspace{2pt}
    CDO (Chief Data Officer - Veri Başkanı), şirketlerin veri stratejilerini yöneten ve karar almayı veriye dayandıran liderlik rolüdür.
\end{tcolorbox}
\vspace{10pt}

\needspace{5cm}
\vspace{8pt}
\noindent\textbf{\large\sffamily Soru 17.} \ Aşağıdakilerden hangisi Büyük Veri (Big Data) kavramını tanımlayan 3V bileşenlerinden biri değildir?
\begin{enumerate}[label=\Alph*), leftmargin=25pt, topsep=4pt, itemsep=2pt]
    \item Volume (Hacim)
    \item Velocity (Hız)
    \item Variety (Çeşitlilik)
    \item Validity (Geçerlilik)
\end{enumerate}
\vspace{4pt}
\begin{tcolorbox}[title=Cevap ve Açıklama, colback=gray!3!white, colframe=primarycolor!90!black, arc=2mm, fonttitle=\bfseries\sffamily, top=6pt, bottom=6pt]
    \textbf{Doğru Cevap: D}
    \vspace{2pt}
    Büyük Verinin 3V'si: Volume (Hacim), Velocity (Hız) ve Variety (Çeşitlilik) bileşenleridir. Validity sonradan eklenen farklı V'lerdendir ama temel 3V içinde yoktur.
\end{tcolorbox}
\vspace{10pt}

\needspace{5cm}
\vspace{8pt}
\noindent\textbf{\large\sffamily Soru 18.} \ Yapay Zeka optimizasyonları kapsamında, bulut tabanlı çalışan ve 1.3 milyondan fazla endüstriyel cihazdan veri toplayan Siemens platformu hangisidir?
\begin{enumerate}[label=\Alph*), leftmargin=25pt, topsep=4pt, itemsep=2pt]
    \item Data Lake
    \item MindSphere
    \item Watson AI
    \item Azure IoT
\end{enumerate}
\vspace{4pt}
\begin{tcolorbox}[title=Cevap ve Açıklama, colback=gray!3!white, colframe=primarycolor!90!black, arc=2mm, fonttitle=\bfseries\sffamily, top=6pt, bottom=6pt]
    \textbf{Doğru Cevap: B}
    \vspace{2pt}
    Siemens MindSphere, endüstriyel sahadaki 1.3 milyondan fazla cihazdan akan büyük verileri toplayan ve analiz eden popüler bir platformdur.
\end{tcolorbox}
\vspace{10pt}

\needspace{5cm}
\vspace{8pt}
\noindent\textbf{\large\sffamily Soru 19.} \ Alibaba bünyesindeki Alipay tabanlı mobil ödeme verilerini ve yapay zekayı kullanarak, 5 yılda 1 milyardan fazla müşteriye ulaşan fintech devi hangisidir?
\begin{enumerate}[label=\Alph*), leftmargin=25pt, topsep=4pt, itemsep=2pt]
    \item JPMorgan Chase
    \item Ant Financial Group
    \item Tencent Finance
    \item SoftBank
\end{enumerate}
\vspace{4pt}
\begin{tcolorbox}[title=Cevap ve Açıklama, colback=gray!3!white, colframe=primarycolor!90!black, arc=2mm, fonttitle=\bfseries\sffamily, top=6pt, bottom=6pt]
    \textbf{Doğru Cevap: B}
    \vspace{2pt}
    Ant Financial Group, Alipay verilerini yapay zeka kararlarıyla birleştirerek sadece 5 yılda 1 milyar müşteriye operasyonel çalışan olmadan ulaşmıştır.
\end{tcolorbox}
\vspace{10pt}

\needspace{5cm}
\vspace{8pt}
\noindent\textbf{\large\sffamily Soru 20.} \ Yapay Zeka Fabrikasının (AI Factory) girdisi olan verileri ölçeklenebilir ve otomatik şekilde toplayan, birleştiren ilk bileşeni hangisidir?
\begin{enumerate}[label=\Alph*), leftmargin=25pt, topsep=4pt, itemsep=2pt]
    \item Deney Platformu
    \item Kestirimci Karar Motoru
    \item Veri Hattı (Data Pipeline)
    \item Kullanıcı Arayüzü
\end{enumerate}
\vspace{4pt}
\begin{tcolorbox}[title=Cevap ve Açıklama, colback=gray!3!white, colframe=primarycolor!90!black, arc=2mm, fonttitle=\bfseries\sffamily, top=6pt, bottom=6pt]
    \textbf{Doğru Cevap: C}
    \vspace{2pt}
    Yapay zeka fabrikasının 4 temel bileşeninden ilki, verileri otomatik toplayıp temizleyen yarı otomatik 'Veri Hattı'dır (Data Pipeline).
\end{tcolorbox}
\vspace{10pt}

\needspace{5cm}
\vspace{8pt}
\noindent\textbf{\large\sffamily Soru 21.} \ Kısıtlı tek bir etki alanında (örneğin sadece satranç oynamak veya kredi skoru tahmin etmek) insan gibi karar veren yapay zeka türü hangisidir?
\begin{enumerate}[label=\Alph*), leftmargin=25pt, topsep=4pt, itemsep=2pt]
    \item Güçlü Yapay Zeka
    \item Zayıf Yapay Zeka (Narrow AI)
    \item Genel Yapay Zeka (AGI)
    \item Otonom Yapay Zeka
\end{enumerate}
\vspace{4pt}
\begin{tcolorbox}[title=Cevap ve Açıklama, colback=gray!3!white, colframe=primarycolor!90!black, arc=2mm, fonttitle=\bfseries\sffamily, top=6pt, bottom=6pt]
    \textbf{Doğru Cevap: B}
    \vspace{2pt}
    Zayıf yapay zeka (Dar yapay zeka), belirli tek bir görev alanında uzmanlaşmış, etki alanı sınırlı olan algoritmik yapılardır.
\end{tcolorbox}
\vspace{10pt}

\needspace{5cm}
\vspace{8pt}
\noindent\textbf{\large\sffamily Soru 22.} \ Küresel nakliyatta, verimsizlikleri aşmak ve tüm tedarik zincirinde gerçek zamanlı veri şeffaflığı sağlamak için dijitalleşen dev konteyner şirketi hangisidir?
\begin{enumerate}[label=\Alph*), leftmargin=25pt, topsep=4pt, itemsep=2pt]
    \item Aeroflot
    \item Maersk
    \item G7 Taksi
    \item Patagonia
\end{enumerate}
\vspace{4pt}
\begin{tcolorbox}[title=Cevap ve Açıklama, colback=gray!3!white, colframe=primarycolor!90!black, arc=2mm, fonttitle=\bfseries\sffamily, top=6pt, bottom=6pt]
    \textbf{Doğru Cevap: B}
    \vspace{2pt}
    Maersk nakliye şirketi, konteyner takibindeki verimsizlik ve şeffaflık sorunlarını aşmak amacıyla dijital sensörler ve tek veri kaynağı kullanımıyla dönüşmüştür.
\end{tcolorbox}
\vspace{10pt}

\needspace{5cm}
\vspace{8pt}
\noindent\textbf{\large\sffamily Soru 23.} \ Aeroflot havayolu şirketinin dijital dönüşüm sonrasında Net Tavsiye Skorunda (NPS) elde ettiği artış oranı yüzde kaçtır?
\begin{enumerate}[label=\Alph*), leftmargin=25pt, topsep=4pt, itemsep=2pt]
    \item \%10
    \item \%20
    \item \%44
    \item \%85
\end{enumerate}
\vspace{4pt}
\begin{tcolorbox}[title=Cevap ve Açıklama, colback=gray!3!white, colframe=primarycolor!90!black, arc=2mm, fonttitle=\bfseries\sffamily, top=6pt, bottom=6pt]
    \textbf{Doğru Cevap: C}
    \vspace{2pt}
    Aeroflot, dijitalleşmeyle operasyonlarını ve raporlamalarını düzelterek Net Tavsiye Skorunu \%44 oranında artırmayı başarmıştır.
\end{tcolorbox}
\vspace{10pt}

\needspace{5cm}
\vspace{8pt}
\noindent\textbf{\large\sffamily Soru 24.} \ 1905 yılında kurulan ve Uber'in getirdiği rekabete kendi mobil rezervasyon yazılımını kurarak başarıyla yanıt veren Paris merkezli geleneksel taksi şirketi hangisidir?
\begin{enumerate}[label=\Alph*), leftmargin=25pt, topsep=4pt, itemsep=2pt]
    \item G7 Taksi
    \item Champs-Élysées Cabs
    \item Marriott
    \item Avnet
\end{enumerate}
\vspace{4pt}
\begin{tcolorbox}[title=Cevap ve Açıklama, colback=gray!3!white, colframe=primarycolor!90!black, arc=2mm, fonttitle=\bfseries\sffamily, top=6pt, bottom=6pt]
    \textbf{Doğru Cevap: A}
    \vspace{2pt}
    G7 Taksi şirketi, dijital yıkım tehdidine karşı hızla dijital araçları kendi süreçlerine uygulayarak pazar payını korumayı başarmıştır.
\end{tcolorbox}
\vspace{10pt}

\needspace{5cm}
\vspace{8pt}
\noindent\textbf{\large\sffamily Soru 25.} \ Fiziksel akıllı mağaza (smart store) ile kişisel alışveriş stilistlerini harmanlayarak 'duygusal bağ' kurmayı hedefleyen Paris merkezli lüks perakendeci hangisidir?
\begin{enumerate}[label=\Alph*), leftmargin=25pt, topsep=4pt, itemsep=2pt]
    \item Hackster.io
    \item Galeries Lafayette Champs-Élysées
    \item Walmart
    \item Patagonia
\end{enumerate}
\vspace{4pt}
\begin{tcolorbox}[title=Cevap ve Açıklama, colback=gray!3!white, colframe=primarycolor!90!black, arc=2mm, fonttitle=\bfseries\sffamily, top=6pt, bottom=6pt]
    \textbf{Doğru Cevap: B}
    \vspace{2pt}
    Galeries Lafayette Champs-Élysées akıllı mağazası, fiziksel ve dijitali (stilistler ve veri entegrasyonu) harmanlayan mükemmel bir hibrit mağazacılık örneğidir.
\end{tcolorbox}
\vspace{10pt}

\needspace{5cm}
\vspace{8pt}
\noindent\textbf{\large\sffamily Soru 26.} \ Şirketlerin satın aldığı teknoloji startup'larını kendi yavaşlığı ve bürokrasisi içinde boğmaması için önerilen çalışma şekli hangisidir?
\begin{enumerate}[label=\Alph*), leftmargin=25pt, topsep=4pt, itemsep=2pt]
    \item Startup'ı ana şirkete tamamen entegre etmek
    \item Yarı bağımsız çalışmasına ve sinerji yaratmasına izin vermek
    \item Startup yöneticilerini işten çıkarmak
    \item Startup'ın teknolojisini kullanmayıp kapatmak
\end{enumerate}
\vspace{4pt}
\begin{tcolorbox}[title=Cevap ve Açıklama, colback=gray!3!white, colframe=primarycolor!90!black, arc=2mm, fonttitle=\bfseries\sffamily, top=6pt, bottom=6pt]
    \textbf{Doğru Cevap: B}
    \vspace{2pt}
    Avnet'in Hackster.io örneğinde olduğu gibi, akıllı şirketler satın aldıkları startup'ları bürokrasiden korumak için yarı bağımsız bırakır.
\end{tcolorbox}
\vspace{10pt}

\needspace{5cm}
\vspace{8pt}
\noindent\textbf{\large\sffamily Soru 27.} \ Adobe şirketinin çalışanlarına kendi yenilikçi fikirlerini geliştirip sunma ve 6 ay boyunca bu fikir üzerinde çalışma fırsatı verdiği ödüllü projesinin adı nedir?
\begin{enumerate}[label=\Alph*), leftmargin=25pt, topsep=4pt, itemsep=2pt]
    \item Turuncu Kutu (Orange Box / Kickbox)
    \item Hackathon Günü
    \item Yeşil Kutu
    \item Startup Kuluçka
\end{enumerate}
\vspace{4pt}
\begin{tcolorbox}[title=Cevap ve Açıklama, colback=gray!3!white, colframe=primarycolor!90!black, arc=2mm, fonttitle=\bfseries\sffamily, top=6pt, bottom=6pt]
    \textbf{Doğru Cevap: A}
    \vspace{2pt}
    Adobe'un 'Turuncu Kutu' (Orange Box/Kickbox) projesi, kurum içi girişimciliği ve çalışanların fikir üretmesini destekleyen bir programdır.
\end{tcolorbox}
\vspace{10pt}

\needspace{5cm}
\vspace{8pt}
\noindent\textbf{\large\sffamily Soru 28.} \ Marriott otel zincirinin Expedia gibi çevrimiçi acentelerin (OTA) tehdidi karşısında uyguladığı başarılı strateji aşağıdakilerden hangisidir?
\begin{enumerate}[label=\Alph*), leftmargin=25pt, topsep=4pt, itemsep=2pt]
    \item Tüm fiziksel otelleri kapatıp sanal otele geçmek
    \item Bilgi sistemlerini aşamalı olarak güncelleyip pazarda test etmek
    \item Tüm startup acentelerini satın almak
    \item Fiyatları yarı yarıya düşürmek
\end{enumerate}
\vspace{4pt}
\begin{tcolorbox}[title=Cevap ve Açıklama, colback=gray!3!white, colframe=primarycolor!90!black, arc=2mm, fonttitle=\bfseries\sffamily, top=6pt, bottom=6pt]
    \textbf{Doğru Cevap: B}
    \vspace{2pt}
    Marriott, altyapısını toptan değiştirmek gibi yüksek riskli bir yol yerine, Expedia rekabetine karşı sistemlerini aşamalı güncelleyip pazarda test etmiştir.
\end{tcolorbox}
\vspace{10pt}

\needspace{5cm}
\vspace{8pt}
\noindent\textbf{\large\sffamily Soru 29.} \ Sürdürülebilirlik kavramının en popüler tanımı olan 'Bugünün ihtiyaçlarını gelecek nesillerin yeteneklerinden ödün vermeden karşılamak' hangi raporda veya kurulda yapılmıştır?
\begin{enumerate}[label=\Alph*), leftmargin=25pt, topsep=4pt, itemsep=2pt]
    \item Kyoto Protokolü
    \item Brundtland Raporu (Büyümenin Sınırları / Ortak Geleceğimiz)
    \item Paris İklim Anlaşması
    \item Dartmouth Konferansı
\end{enumerate}
\vspace{4pt}
\begin{tcolorbox}[title=Cevap ve Açıklama, colback=gray!3!white, colframe=primarycolor!90!black, arc=2mm, fonttitle=\bfseries\sffamily, top=6pt, bottom=6pt]
    \textbf{Doğru Cevap: B}
    \vspace{2pt}
    Sürdürülebilirliğin bu klasik tanımı, Birleşmiş Milletler Brundtland Komisyonu'nun 1987 tarihli 'Ortak Geleceğimiz' (Brundtland) raporunda yapılmıştır.
\end{tcolorbox}
\vspace{10pt}

\needspace{5cm}
\vspace{8pt}
\noindent\textbf{\large\sffamily Soru 30.} \ Aşağıdakilerden hangisi sürdürülebilirliğin değerlendirildiği üç ana boyuttan (sütundan) biri değildir?
\begin{enumerate}[label=\Alph*), leftmargin=25pt, topsep=4pt, itemsep=2pt]
    \item Çevresel sürdürülebilirlik
    \item Sosyal sürdürülebilirlik
    \item Kültürel asimilasyon
    \item Ekonomik sürdürülebilirlik
\end{enumerate}
\vspace{4pt}
\begin{tcolorbox}[title=Cevap ve Açıklama, colback=gray!3!white, colframe=primarycolor!90!black, arc=2mm, fonttitle=\bfseries\sffamily, top=6pt, bottom=6pt]
    \textbf{Doğru Cevap: C}
    \vspace{2pt}
    Sürdürülebilirliğin üç ana boyutu: Çevresel, Sosyal ve Ekonomik sürdürülebilirliktir.
\end{tcolorbox}
\vspace{10pt}

\needspace{5cm}
\vspace{8pt}
\noindent\textbf{\large\sffamily Soru 31.} \ Garrett Hardin tarafından ortaya atılan, bireylerin ücretsiz ortak malları kendi çıkarları için aşırı tüketerek kaynağı yok etmesini ifade eden kuram hangisidir?
\begin{enumerate}[label=\Alph*), leftmargin=25pt, topsep=4pt, itemsep=2pt]
    \item Yaratıcı Yıkım
    \item Ortak Malların Trajedisi (Tragedy of the Commons)
    \item Tersine Lojistik
    \item Boyutluluğun Laneti
\end{enumerate}
\vspace{4pt}
\begin{tcolorbox}[title=Cevap ve Açıklama, colback=gray!3!white, colframe=primarycolor!90!black, arc=2mm, fonttitle=\bfseries\sffamily, top=6pt, bottom=6pt]
    \textbf{Doğru Cevap: B}
    \vspace{2pt}
    Ortak Malların Trajedisi (Tragedy of the Commons), serbest ve ortak meraların aşırı hayvan otlatılmasıyla tükenmesi gibi durumları açıklayan kuramdır.
\end{tcolorbox}
\vspace{10pt}

\needspace{5cm}
\vspace{8pt}
\noindent\textbf{\large\sffamily Soru 32.} \ Nike firmasının 1996 yılında karşılaştığı ve marka imajını sarsan, daha sonrasında Kurumsal Sosyal Sorumluluk (KSS) adımları atmasına neden olan skandal nedir?
\begin{enumerate}[label=\Alph*), leftmargin=25pt, topsep=4pt, itemsep=2pt]
    \item Çevreyi kimyasal atıklarla kirletmesi
    \item Pakistan'da çocuk işçi çalıştırıldığının ortaya çıkması
    \item Ürün fiyatlarını aşırı yüksek belirlemesi
    \item Karbon ayak izini saklaması
\end{enumerate}
\vspace{4pt}
\begin{tcolorbox}[title=Cevap ve Açıklama, colback=gray!3!white, colframe=primarycolor!90!black, arc=2mm, fonttitle=\bfseries\sffamily, top=6pt, bottom=6pt]
    \textbf{Doğru Cevap: B}
    \vspace{2pt}
    Nike, 1996 yılında tedarikçilerinin Pakistan'da çocuk işçilere futbol topu diktirdiğinin ifşa olmasıyla büyük bir itibar krizi yaşamış ve sosyal denetim kurallarını sıkılaştırmıştır.
\end{tcolorbox}
\vspace{10pt}

\needspace{5cm}
\vspace{8pt}
\noindent\textbf{\large\sffamily Soru 33.} \ 3 Aralık 1984 tarihinde, Union Carbide firmasına ait böcek ilacı fabrikasından sızan MIC gazının binlerce insanı öldürdüğü endüstriyel çevre felaketi nerede yaşanmıştır?
\begin{enumerate}[label=\Alph*), leftmargin=25pt, topsep=4pt, itemsep=2pt]
    \item Bhopal / Hindistan
    \item Çernobil / Ukrayna
    \item Fukushima / Japonya
    \item Londra / İngiltere
\end{enumerate}
\vspace{4pt}
\begin{tcolorbox}[title=Cevap ve Açıklama, colback=gray!3!white, colframe=primarycolor!90!black, arc=2mm, fonttitle=\bfseries\sffamily, top=6pt, bottom=6pt]
    \textbf{Doğru Cevap: A}
    \vspace{2pt}
    Bhopal felaketi, Hindistan'ın Bhopal kentindeki bir ABD böcek ilacı fabrikasından sızan Metil İzosiyanat gazının yol açtığı tarihin en büyük endüstriyel kazalarından biridir.
\end{tcolorbox}
\vspace{10pt}

\needspace{5cm}
\vspace{8pt}
\noindent\textbf{\large\sffamily Soru 34.} \ Ürünlerin, hammadde ve yarı mamullerin tüketim noktasından orijine (çıkış noktasına) doğru, değer kazanımı veya imha amacıyla akışını yöneten süreç hangisidir?
\begin{enumerate}[label=\Alph*), leftmargin=25pt, topsep=4pt, itemsep=2pt]
    \item Yeşil Satın Alma
    \item Tersine Lojistik (Reverse Logistics)
    \item Intermodal Taşımacılık
    \item Tedarikçi Envanteri
\end{enumerate}
\vspace{4pt}
\begin{tcolorbox}[title=Cevap ve Açıklama, colback=gray!3!white, colframe=primarycolor!90!black, arc=2mm, fonttitle=\bfseries\sffamily, top=6pt, bottom=6pt]
    \textbf{Doğru Cevap: B}
    \vspace{2pt}
    Tersine lojistik (reverse logistics), ürünlerin geri kazanılması, yenilenmesi, geri dönüştürülmesi veya yok edilmesi için tüketimden kaynağa doğru akışını yöneten lojistik disiplinidir.
\end{tcolorbox}
\vspace{10pt}

\needspace{5cm}
\vspace{8pt}
\noindent\textbf{\large\sffamily Soru 35.} \ Yeniden üretim (re-manufacturing) süreçlerinde ihtiyaç duyulan enerji miktarı, hammaddeden sıfırdan üretim yapmaya göre yüzde kaç daha azdır?
\begin{enumerate}[label=\Alph*), leftmargin=25pt, topsep=4pt, itemsep=2pt]
    \item \%10'a kadar
    \item \%50'ye kadar
    \item \%85'e kadar
    \item \%5'e kadar
\end{enumerate}
\vspace{4pt}
\begin{tcolorbox}[title=Cevap ve Açıklama, colback=gray!3!white, colframe=primarycolor!90!black, arc=2mm, fonttitle=\bfseries\sffamily, top=6pt, bottom=6pt]
    \textbf{Doğru Cevap: C}
    \vspace{2pt}
    Tersine lojistik verilerine göre yeniden üretim, ham maden çıkarmaya ve işlemeye göre \%85'e varan oranda enerji tasarrufu sağlar.
\end{tcolorbox}
\vspace{10pt}

\needspace{5cm}
\vspace{8pt}
\noindent\textbf{\large\sffamily Soru 36.} \ Geri dönüştürülmüş alüminyum üretmek, sıfırdan alüminyum üretmek için gereken enerjinin yalnızca yüzde kaçını tüketir?
\begin{enumerate}[label=\Alph*), leftmargin=25pt, topsep=4pt, itemsep=2pt]
    \item \%5'ini
    \item \%25'ini
    \item \%50'sini
    \item \%85'ini
\end{enumerate}
\vspace{4pt}
\begin{tcolorbox}[title=Cevap ve Açıklama, colback=gray!3!white, colframe=primarycolor!90!black, arc=2mm, fonttitle=\bfseries\sffamily, top=6pt, bottom=6pt]
    \textbf{Doğru Cevap: A}
    \vspace{2pt}
    Alüminyumun geri dönüşümü son derece verimlidir; sıfırdan üretimin gerektirdiği devasa elektrik enerjisinin sadece \%5'ini kullanır.
\end{tcolorbox}
\vspace{10pt}

\needspace{5cm}
\vspace{8pt}
\noindent\textbf{\large\sffamily Soru 37.} \ 3M firmasının, araçların yükleme yüksekliğini ayarlayan raflar geliştirerek lojistik süreçte elde ettiği günlük araç doluluk artış oranı yüzde kaçtır?
\begin{enumerate}[label=\Alph*), leftmargin=25pt, topsep=4pt, itemsep=2pt]
    \item \%10
    \item \%20
    \item \%40
    \item \%80
\end{enumerate}
\vspace{4pt}
\begin{tcolorbox}[title=Cevap ve Açıklama, colback=gray!3!white, colframe=primarycolor!90!black, arc=2mm, fonttitle=\bfseries\sffamily, top=6pt, bottom=6pt]
    \textbf{Doğru Cevap: C}
    \vspace{2pt}
    3M, akıllı raf tasarımlarıyla araç yükleme verimliliğini \%40 oranında artırarak daha az seferle aynı miktarda ürün taşımayı başarmıştır.
\end{tcolorbox}
\vspace{10pt}

\needspace{5cm}
\vspace{8pt}
\noindent\textbf{\large\sffamily Soru 38.} \ Dünyada üretilen toplam enerjinin yüzde kaçı şehirlerde tüketilmektedir?
\begin{enumerate}[label=\Alph*), leftmargin=25pt, topsep=4pt, itemsep=2pt]
    \item \%20'si
    \item \%55'i
    \item \%75'i
    \item \%95'i
\end{enumerate}
\vspace{4pt}
\begin{tcolorbox}[title=Cevap ve Açıklama, colback=gray!3!white, colframe=primarycolor!90!black, arc=2mm, fonttitle=\bfseries\sffamily, top=6pt, bottom=6pt]
    \textbf{Doğru Cevap: C}
    \vspace{2pt}
    Kentlerin yoğunlaşması nedeniyle, dünyada üretilen elektrik ve ısınma enerjisi gibi kaynakların \%75'i şehirlerde tüketilmektedir.
\end{tcolorbox}
\vspace{10pt}

\needspace{5cm}
\vspace{8pt}
\noindent\textbf{\large\sffamily Soru 39.} \ Günümüzde dünya nüfusunun yaklaşık \%55'i şehirlerde yaşarken, 2050 yılına kadar bu oranın yüzde kaça ulaşması beklenmektedir?
\begin{enumerate}[label=\Alph*), leftmargin=25pt, topsep=4pt, itemsep=2pt]
    \item \%60
    \item \%68
    \item \%85
    \item \%95
\end{enumerate}
\vspace{4pt}
\begin{tcolorbox}[title=Cevap ve Açıklama, colback=gray!3!white, colframe=primarycolor!90!black, arc=2mm, fonttitle=\bfseries\sffamily, top=6pt, bottom=6pt]
    \textbf{Doğru Cevap: B}
    \vspace{2pt}
    Projeksiyonlara göre hızlı kentleşmeyle birlikte 2050 yılına kadar dünya nüfusunun \%68'i (yaklaşık 3 kişiden 2'si) şehirlerde yaşayacaktır.
\end{tcolorbox}
\vspace{10pt}

\needspace{5cm}
\vspace{8pt}
\noindent\textbf{\large\sffamily Soru 40.} \ Kentsel planlamada 'çevre ve ekoloji' faktörünü merkeze alan kavram 'kentsel sürdürülebilirlik' iken, 'sosyal, ekonomik ve veri odaklı teknolojik katılımı' merkeze alan kavram hangisidir?
\begin{enumerate}[label=\Alph*), leftmargin=25pt, topsep=4pt, itemsep=2pt]
    \item Sanayi 4.0
    \item Akıllı Şehir (Smart City)
    \item Bilgi Toplumu
    \item E-İş Modeli
\end{enumerate}
\vspace{4pt}
\begin{tcolorbox}[title=Cevap ve Açıklama, colback=gray!3!white, colframe=primarycolor!90!black, arc=2mm, fonttitle=\bfseries\sffamily, top=6pt, bottom=6pt]
    \textbf{Doğru Cevap: B}
    \vspace{2pt}
    Neirotti ve ark. çalışmasına göre kentsel sürdürülebilirlik çevreye odaklanırken; akıllı şehir kavramı daha çok sosyal, ekonomik ve ICT tabanlı teknolojik süreçleri içerir.
\end{tcolorbox}
\vspace{10pt}

\needspace{5cm}
\vspace{8pt}
\noindent\textbf{\large\sffamily Soru 41.} \ Ateş (2018) çalışmasına göre kentsel alanlarda karbon ayak izini sürekli azaltan ve kısıtlı kaynakları verimli kullanan yapılara ne ad verilir?
\begin{enumerate}[label=\Alph*), leftmargin=25pt, topsep=4pt, itemsep=2pt]
    \item Geleneksel Şehir
    \item Sürdürülebilir Akıllı Şehir
    \item E-Ticaret Ekosistemi
    \item Sanayi Bölgesi
\end{enumerate}
\vspace{4pt}
\begin{tcolorbox}[title=Cevap ve Açıklama, colback=gray!3!white, colframe=primarycolor!90!black, arc=2mm, fonttitle=\bfseries\sffamily, top=6pt, bottom=6pt]
    \textbf{Doğru Cevap: B}
    \vspace{2pt}
    Ateş'e göre 'Sürdürülebilir Akıllı Şehir', kısıtlı kaynakları etkin kullanıp kentsel karbon ayak izini azaltan ve kalkınmaya yatırım yapan şehirdir.
\end{tcolorbox}
\vspace{10pt}

\needspace{5cm}
\vspace{8pt}
\noindent\textbf{\large\sffamily Soru 42.} \ BM Çevre Programı (UNEP) ve BM Dünya Turizm Örgütü (UNWTO) tarafından belirlenen Sürdürülebilir Turizm İlkelerinden 'Yerel Kalkınma' neyi hedefler?
\begin{enumerate}[label=\Alph*), leftmargin=25pt, topsep=4pt, itemsep=2pt]
    \item Ziyaretçilerin yerel halkla hiçbir şekilde iletişime geçmemesini
    \item Turizm gelirlerinin yerel üretici ve destinasyonda kalma oranını artırmayı
    \item Turist sayısını sınırsız şekilde artırmayı
    \item Yerel kültürü tamamen değiştirmeyi
\end{enumerate}
\vspace{4pt}
\begin{tcolorbox}[title=Cevap ve Açıklama, colback=gray!3!white, colframe=primarycolor!90!black, arc=2mm, fonttitle=\bfseries\sffamily, top=6pt, bottom=6pt]
    \textbf{Doğru Cevap: B}
    \vspace{2pt}
    Yerel kalkınma ilkesi, turistlerin yerelde harcama yapmasını ve bu kazancın destinasyon halkına (yerel üreticilere) doğrudan katkı sağlamasını hedefler.
\end{tcolorbox}
\vspace{10pt}

\needspace{5cm}
\vspace{8pt}
\noindent\textbf{\large\sffamily Soru 43.} \ Sürdürülebilir turizmde, koruma-kullanma dengesini gözeterek endemik türlerin ve yaban hayatının korunmasını destekleyen ilke hangisidir?
\begin{enumerate}[label=\Alph*), leftmargin=25pt, topsep=4pt, itemsep=2pt]
    \item İstihdam Kalitesi
    \item Biyolojik Çeşitlilik (Biodiversity)
    \item Ziyaretçi Memnuniyeti
    \item Yerelden Kontrol
\end{enumerate}
\vspace{4pt}
\begin{tcolorbox}[title=Cevap ve Açıklama, colback=gray!3!white, colframe=primarycolor!90!black, arc=2mm, fonttitle=\bfseries\sffamily, top=6pt, bottom=6pt]
    \textbf{Doğru Cevap: B}
    \vspace{2pt}
    Biyolojik çeşitlilik ilkesi, yaban hayatının, türlerin ve endemizmin zarar görmeden gelecek nesillere aktarılmasını hedefler.
\end{tcolorbox}
\vspace{10pt}

\needspace{5cm}
\vspace{8pt}
\noindent\textbf{\large\sffamily Soru 44.} \ Kültürel sürdürülebilirlik kavramında 'Kültürel Süreklilik' için ana vurgu aşağıdakilerden hangisidir?
\begin{enumerate}[label=\Alph*), leftmargin=25pt, topsep=4pt, itemsep=2pt]
    \item Kültürün hiçbir şekilde değişmemesi ve eski halinde donması
    \item Geçmiş mirasının ve yerleşim ilkelerinin çok büyük birer kaynak olarak fark edilmesi
    \item Küresel kültürün yerel kültürü tamamen yok etmesi
    \item Geleneklerin yasal olarak yasaklanması
\end{enumerate}
\vspace{4pt}
\begin{tcolorbox}[title=Cevap ve Açıklama, colback=gray!3!white, colframe=primarycolor!90!black, arc=2mm, fonttitle=\bfseries\sffamily, top=6pt, bottom=6pt]
    \textbf{Doğru Cevap: B}
    \vspace{2pt}
    Cebeci ve Çakılcıoğlu'na göre kültürel sürdürülebilirlik gelenekçi/donmuş bir tavır değil; geçmişin ve yerleşim ilkelerinin büyük birer değer kaynağı olduğunun fark edilerek yeni kuşaklara aktarılmasıdır.
\end{tcolorbox}
\vspace{10pt}

\needspace{5cm}
\vspace{8pt}
\noindent\textbf{\large\sffamily Soru 45.} \ Avrupa Birliği genelinde yapay zekâ uygulamalarını taşıdıkları risklere göre denetleyen ilk kapsamlı yasal regülasyon hangisidir?
\begin{enumerate}[label=\Alph*), leftmargin=25pt, topsep=4pt, itemsep=2pt]
    \item KVKK
    \item EU AI Act (AB Yapay Zeka Yasası)
    \item GDPR
    \item Brundtland Raporu
\end{enumerate}
\vspace{4pt}
\begin{tcolorbox}[title=Cevap ve Açıklama, colback=gray!3!white, colframe=primarycolor!90!black, arc=2mm, fonttitle=\bfseries\sffamily, top=6pt, bottom=6pt]
    \textbf{Doğru Cevap: B}
    \vspace{2pt}
    EU AI Act, yapay zeka sistemlerini taşıdıkları risklere göre sınıflandırıp denetleyen ilk ve en kapsamlı küresel yapay zeka mevzuatıdır.
\end{tcolorbox}
\vspace{10pt}

\needspace{5cm}
\vspace{8pt}
\noindent\textbf{\large\sffamily Soru 46.} \ Aşağıdakilerden hangisi çevre dostu binaların (yeşil binalar) geleneksel binalara kıyasla sağladığı enerji tasarruf oranını göstermektedir?
\begin{enumerate}[label=\Alph*), leftmargin=25pt, topsep=4pt, itemsep=2pt]
    \item \%5-10
    \item \%24-50
    \item \%90-95
    \item Tasarruf sağlamaz
\end{enumerate}
\vspace{4pt}
\begin{tcolorbox}[title=Cevap ve Açıklama, colback=gray!3!white, colframe=primarycolor!90!black, arc=2mm, fonttitle=\bfseries\sffamily, top=6pt, bottom=6pt]
    \textbf{Doğru Cevap: B}
    \vspace{2pt}
    Çevreci binalar enerji tüketiminde \%24-50 arasında yüksek bir tasarruf sağlayarak karbon salınımını büyük oranda azaltırlar.
\end{tcolorbox}
\vspace{10pt}

\needspace{5cm}
\vspace{8pt}
\noindent\textbf{\large\sffamily Soru 47.} \ Lojistikte, taşıma araçlarının fosil yakıt motor standartlarını (Euro 5, Euro 6 gibi) iyileştirmeyi veya elektrikli araçlar kullanmayı hedefleyen strateji hangisidir?
\begin{enumerate}[label=\Alph*), leftmargin=25pt, topsep=4pt, itemsep=2pt]
    \item Daha Az Taşı
    \item Daha Kısa Taşı
    \item Daha Temiz Taşı
    \item Tersine Taşı
\end{enumerate}
\vspace{4pt}
\begin{tcolorbox}[title=Cevap ve Açıklama, colback=gray!3!white, colframe=primarycolor!90!black, arc=2mm, fonttitle=\bfseries\sffamily, top=6pt, bottom=6pt]
    \textbf{Doğru Cevap: C}
    \vspace{2pt}
    Daha Temiz Taşı stratejisi, Euro 5/6 motorların seçilmesi, fosil yakıt dışı otonom/elektrikli araçların kullanılması gibi karbon emisyonunu düşüren yöntemleri içerir.
\end{tcolorbox}
\vspace{10pt}

\needspace{5cm}
\vspace{8pt}
\noindent\textbf{\large\sffamily Soru 48.} \ İstanbul, Trieste ve Mannheim limanları arasında Ekol Lojistik tarafından gerçekleştirilen ve karayolu yerine Ro-Ro/Tren kullanılan çevre dostu lojistik modeli hangisidir?
\begin{enumerate}[label=\Alph*), leftmargin=25pt, topsep=4pt, itemsep=2pt]
    \item Havayolu Taşımacılığı
    \item İntermodal Taşımacılık
    \item Boru Hattı Lojistiği
    \item Kurye Dağıtımı
\end{enumerate}
\vspace{4pt}
\begin{tcolorbox}[title=Cevap ve Açıklama, colback=gray!3!white, colframe=primarycolor!90!black, arc=2mm, fonttitle=\bfseries\sffamily, top=6pt, bottom=6pt]
    \textbf{Doğru Cevap: B}
    \vspace{2pt}
    Ekol Lojistik'in yürüttüğü bu model, birden fazla taşıma türünü (deniz, demir, kara) birleştiren ve CO2 salınımını devasa oranda düşüren bir İntermodal taşımacılık örneğidir.
\end{tcolorbox}
\vspace{10pt}

\needspace{5cm}
\vspace{8pt}
\noindent\textbf{\large\sffamily Soru 49.} \ Tedarik zincirinde, paketleme malzemelerinin doğada kolayca çözünebilen veya tekrar kullanılabilir (reusable) materyallerden seçilmesini ifade eden kavram hangisidir?
\begin{enumerate}[label=\Alph*), leftmargin=25pt, topsep=4pt, itemsep=2pt]
    \item Çevreci Satın Alma
    \item Yeşil Paketleme (Green Packaging)
    \item Tersine Dağıtım
    \item Ürün Yaşam Döngüsü
\end{enumerate}
\vspace{4pt}
\begin{tcolorbox}[title=Cevap ve Açıklama, colback=gray!3!white, colframe=primarycolor!90!black, arc=2mm, fonttitle=\bfseries\sffamily, top=6pt, bottom=6pt]
    \textbf{Doğru Cevap: B}
    \vspace{2pt}
    Yeşil paketleme (green packaging), doğaya zarar vermeyen, plastikten ziyade kağıt veya biyolojik çözünür hammaddelerden üretilen paketleme sürecidir.
\end{tcolorbox}
\vspace{10pt}

\needspace{5cm}
\vspace{8pt}
\noindent\textbf{\large\sffamily Soru 50.} \ Sürdürülebilir kalkınmanın gerçekleşmesinde ekolojik dengenin korunmasının en temel nedeni aşağıdakilerden hangisidir?
\begin{enumerate}[label=\Alph*), leftmargin=25pt, topsep=4pt, itemsep=2pt]
    \item Ekonomik büyümenin süreklilik arz edebilmesinin ekolojik dengeye bağlı olması
    \item Teknoloji yatırımlarının maliyetini düşürmek
    \item Şehir nüfusunun hızla köylere geri dönmesini sağlamak
    \item Sanayi üretimini tamamen durdurmak
\end{enumerate}
\vspace{4pt}
\begin{tcolorbox}[title=Cevap ve Açıklama, colback=gray!3!white, colframe=primarycolor!90!black, arc=2mm, fonttitle=\bfseries\sffamily, top=6pt, bottom=6pt]
    \textbf{Doğru Cevap: A}
    \vspace{2pt}
    Ekonomik büyümenin sürdürülebilir olması ve devamlılığı için ekolojik dengenin bozulmaması, su, hava ve toprak gibi kıt kaynakların korunması zorunludur.
\end{tcolorbox}
\vspace{10pt}

\needspace{5cm}
\vspace{8pt}
\noindent\textbf{\large\sffamily Soru 51.} \ Toplumların tarihsel gelişim evrelerinde yerleşik yaşama geçiş ve tarım arazilerinin yönetimi hangi devrimle başlamıştır?
\begin{enumerate}[label=\Alph*), leftmargin=25pt, topsep=4pt, itemsep=2pt]
    \item Sanayi Devrimi
    \item Tarım Devrimi
    \item Bilgi Devrimi
    \item Dijital Devrim
\end{enumerate}
\vspace{4pt}
\begin{tcolorbox}[title=Cevap ve Açıklama, colback=gray!3!white, colframe=primarycolor!90!black, arc=2mm, fonttitle=\bfseries\sffamily, top=6pt, bottom=6pt]
    \textbf{Doğru Cevap: B}
    \vspace{2pt}
    Tarım Devrimi ile insanlar göçebe yaşamdan yerleşik yaşama geçmiş, toprağı işlemeye ve tarım arazilerini yönetmeye başlamışlardır.
\end{tcolorbox}
\vspace{10pt}

\needspace{5cm}
\vspace{8pt}
\noindent\textbf{\large\sffamily Soru 52.} \ Sanayi Toplumu evresinin en belirgin ekonomik üretim özelliği aşağıdakilerden hangisidir?
\begin{enumerate}[label=\Alph*), leftmargin=25pt, topsep=4pt, itemsep=2pt]
    \item Toplu kişiselleştirme (Mass customization)
    \item Standart ve seri üretim (Mass production)
    \item Esnek el sanatları üretimi
    \item Hizmet odaklı bulut üretimi
\end{enumerate}
\vspace{4pt}
\begin{tcolorbox}[title=Cevap ve Açıklama, colback=gray!3!white, colframe=primarycolor!90!black, arc=2mm, fonttitle=\bfseries\sffamily, top=6pt, bottom=6pt]
    \textbf{Doğru Cevap: B}
    \vspace{2pt}
    Sanayi Toplumu döneminde fabrikaların kurulmasıyla birlikte standart, tek tip ve seri üretim (mass production) ekonominin ana unsuru haline gelmiştir.
\end{tcolorbox}
\vspace{10pt}

\needspace{5cm}
\vspace{8pt}
\noindent\textbf{\large\sffamily Soru 53.} \ Bilgisayar destekli üretim (CAD/CAM) teknolojilerinin yaygınlaşması ve küreselleşmenin hızlanması hangi toplum evresine aittir?
\begin{enumerate}[label=\Alph*), leftmargin=25pt, topsep=4pt, itemsep=2pt]
    \item İlkel Toplum
    \item Tarım Toplumu
    \item Sanayi Toplumu
    \item Bilgi Toplumu
\end{enumerate}
\vspace{4pt}
\begin{tcolorbox}[title=Cevap ve Açıklama, colback=gray!3!white, colframe=primarycolor!90!black, arc=2mm, fonttitle=\bfseries\sffamily, top=6pt, bottom=6pt]
    \textbf{Doğru Cevap: D}
    \vspace{2pt}
    Bilgi Toplumu evresinde bilgisayar destekli üretim teknolojileri (CAD/CAM) ve iletişim teknolojilerindeki gelişmeler üretimi ve küresel rekabeti yönlendirmiştir.
\end{tcolorbox}
\vspace{10pt}

\needspace{5cm}
\vspace{8pt}
\noindent\textbf{\large\sffamily Soru 54.} \ Dijital dönüşümün 4C ilkelerinden 'Collaboration' (İş Birliği) kavramı aşağıdakilerden hangisini ifade eder?
\begin{enumerate}[label=\Alph*), leftmargin=25pt, topsep=4pt, itemsep=2pt]
    \item Şirket içindeki departmanların ve dış ekosistemlerin entegre çalışmasını
    \item Sadece verilerin bir sunucuya bağlanmasını
    \item Finansal işlemlerin dijitalleştirilmesini
    \item Otomatik üretim bantlarının koordinasyonunu
\end{enumerate}
\vspace{4pt}
\begin{tcolorbox}[title=Cevap ve Açıklama, colback=gray!3!white, colframe=primarycolor!90!black, arc=2mm, fonttitle=\bfseries\sffamily, top=6pt, bottom=6pt]
    \textbf{Doğru Cevap: A}
    \vspace{2pt}
    Collaboration (İş Birliği), dijital araçları kullanarak kurum içi departmanların ve dış paydaşların iş birliği içinde değer üretmesini ifade eder.
\end{tcolorbox}
\vspace{10pt}

\needspace{5cm}
\vspace{8pt}
\noindent\textbf{\large\sffamily Soru 55.} \ Dijital dönüşümde 'Connection' (Bağlantı) ilkesi aşağıdakilerden hangisini doğrudan sağlar?
\begin{enumerate}[label=\Alph*), leftmargin=25pt, topsep=4pt, itemsep=2pt]
    \item Fiyatların ucuzlamasını
    \item Sistemlerin, sensörlerin ve verilerin birbirleriyle haberleşmesini
    \item Rekabet yasalarının esnetilmesini
    \item İş gücünün tamamen elenmesini
\end{enumerate}
\vspace{4pt}
\begin{tcolorbox}[title=Cevap ve Açıklama, colback=gray!3!white, colframe=primarycolor!90!black, arc=2mm, fonttitle=\bfseries\sffamily, top=6pt, bottom=6pt]
    \textbf{Doğru Cevap: B}
    \vspace{2pt}
    Connection (Bağlantı), nesnelerin interneti (IoT), API'ler ve ağlar aracılığıyla veri ve sistemlerin birbirine bağlanmasını temsil eder.
\end{tcolorbox}
\vspace{10pt}

\needspace{5cm}
\vspace{8pt}
\noindent\textbf{\large\sffamily Soru 56.} \ İlkel toplumun çöküşüne ve insanların Tarım Toplumu'na geçmesine neden olan kriz aşağıdakilerden hangisidir?
\begin{enumerate}[label=\Alph*), leftmargin=25pt, topsep=4pt, itemsep=2pt]
    \item Makinelerin yetersiz kalması krizidir
    \item Doğal kaynakların kontrolsüzce tüketilmesi sonucu yaşanan kıtlıktır
    \item Yazılımların çökmesi krizidir
    \item Küresel ticaret savaşları krizidir
\end{enumerate}
\vspace{4pt}
\begin{tcolorbox}[title=Cevap ve Açıklama, colback=gray!3!white, colframe=primarycolor!90!black, arc=2mm, fonttitle=\bfseries\sffamily, top=6pt, bottom=6pt]
    \textbf{Doğru Cevap: B}
    \vspace{2pt}
    İlkel toplum doğal kaynakların toplayıcılığına bağımlıydı. Doğal kaynakların azalması bir gıda/kaynak krizine yol açmış ve insanlığı toprağı ekip biçmeye (Tarım Devrimi) yöneltmiştir.
\end{tcolorbox}
\vspace{10pt}

\needspace{5cm}
\vspace{8pt}
\noindent\textbf{\large\sffamily Soru 57.} \ Sanayi Toplumu'ndan Bilgi Toplumu'na evrilmeyi tetikleyen temel piyasa koşulu aşağıdakilerden hangisidir?
\begin{enumerate}[label=\Alph*), leftmargin=25pt, topsep=4pt, itemsep=2pt]
    \item Standart üretimlerin pazarda doygunluğa ulaşması ve kişiselleştirilmiş ürün talebi
    \item Tarım arazilerinin tamamen çölleşmesi
    \item El işçiliğinin tamamen yasaklanması
    \item Para birimlerinin dijitalleşmesi
\end{enumerate}
\vspace{4pt}
\begin{tcolorbox}[title=Cevap ve Açıklama, colback=gray!3!white, colframe=primarycolor!90!black, arc=2mm, fonttitle=\bfseries\sffamily, top=6pt, bottom=6pt]
    \textbf{Doğru Cevap: A}
    \vspace{2pt}
    Sanayi Toplumu'nun tek tip standart üretimi pazarı doyurmuş, rekabet edebilmek için esnek ve kişiselleştirilmiş (mass customization) ürünler sunma ihtiyacı Bilgi Toplumu'nu doğurmuştur.
\end{tcolorbox}
\vspace{10pt}

\needspace{5cm}
\vspace{8pt}
\noindent\textbf{\large\sffamily Soru 58.} \ Dijital dönüşümün nedenlerinden biri olan 'sektörel sınırların bulanıklaşması' ne anlama gelmektedir?
\begin{enumerate}[label=\Alph*), leftmargin=25pt, topsep=4pt, itemsep=2pt]
    \item Şirketlerin fiziksel mağazalarını tamamen kapatması
    \item Geleneksel şirketlerin dijital teknolojiler kullanarak farklı sektörlerde (finans, lojistik vb.) hizmet vermeye başlaması
    \item Şirketlerin yalnızca tek bir ürüne odaklanması
    \item Devletlerin ithalat vergilerini kaldırması
\end{enumerate}
\vspace{4pt}
\begin{tcolorbox}[title=Cevap ve Açıklama, colback=gray!3!white, colframe=primarycolor!90!black, arc=2mm, fonttitle=\bfseries\sffamily, top=6pt, bottom=6pt]
    \textbf{Doğru Cevap: B}
    \vspace{2pt}
    Sektörel sınırların bulanıklaşması, örneğin bir perakende veya nakliye şirketinin dijital altyapısıyla finansal ödeme sistemleri (fintech) sunarak bankalarla rekabet etmesidir.
\end{tcolorbox}
\vspace{10pt}

\needspace{5cm}
\vspace{8pt}
\noindent\textbf{\large\sffamily Soru 59.} \ Schumpeter'in ortaya attığı 'Yaratıcı Yıkım' kavramı dijitalleşme sürecinde şirketler için neyi öngörür?
\begin{enumerate}[label=\Alph*), leftmargin=25pt, topsep=4pt, itemsep=2pt]
    \item Devletin tüm şirketleri satın alacağını
    \item Teknolojik yenilik yapmayan şirketlerin pazar dışı kalarak yok olacağını
    \item Üretimde hiçbir zaman hata yapılmayacağını
    \item Şirketlerin birleşerek tek bir dev haline geleceğini
\end{enumerate}
\vspace{4pt}
\begin{tcolorbox}[title=Cevap ve Açıklama, colback=gray!3!white, colframe=primarycolor!90!black, arc=2mm, fonttitle=\bfseries\sffamily, top=6pt, bottom=6pt]
    \textbf{Doğru Cevap: B}
    \vspace{2pt}
    Yaratıcı Yıkım (Creative Destruction), yeni teknolojilerin eski verimsiz iş modellerini yıkarak yerlerine daha verimli sistemler getirmesidir. Dönüşmeyen şirketler elenir.
\end{tcolorbox}
\vspace{10pt}

\needspace{5cm}
\vspace{8pt}
\noindent\textbf{\large\sffamily Soru 60.} \ Aşağıdakilerden hangisi dijital dönüşümün 4C ilkelerinden 'Coordination' (Koordinasyon) kavramını açıklar?
\begin{enumerate}[label=\Alph*), leftmargin=25pt, topsep=4pt, itemsep=2pt]
    \item İş süreçlerinin dijital araçlarla uyumlu, akıcı ve senkronize bir şekilde yönetilmesi
    \item Sadece bulut depolama alanı satın alınması
    \item Startup firmaların satın alınması
    \item Çocuk işçi çalıştırılmasının engellenmesi
\end{enumerate}
\vspace{4pt}
\begin{tcolorbox}[title=Cevap ve Açıklama, colback=gray!3!white, colframe=primarycolor!90!black, arc=2mm, fonttitle=\bfseries\sffamily, top=6pt, bottom=6pt]
    \textbf{Doğru Cevap: A}
    \vspace{2pt}
    Coordination (Koordinasyon), dijitalleşen süreçler ve departmanlar arasındaki faaliyetlerin uyumlu ve verimli şekilde senkronize edilmesidir.
\end{tcolorbox}
\vspace{10pt}

\needspace{5cm}
\vspace{8pt}
\noindent\textbf{\large\sffamily Soru 61.} \ Bir kütüphanedeki eski el yazması kitapların sayfalarının fotoğraflanarak bilgisayar ortamına PDF olarak aktarılması işlemi hangisidir?
\begin{enumerate}[label=\Alph*), leftmargin=25pt, topsep=4pt, itemsep=2pt]
    \item Dijitalleşme (Digitalization)
    \item Dijital Dönüşüm (Digital Transformation)
    \item Sayısallaştırma (Digitization)
    \item Sanal Gerçeklik
\end{enumerate}
\vspace{4pt}
\begin{tcolorbox}[title=Cevap ve Açıklama, colback=gray!3!white, colframe=primarycolor!90!black, arc=2mm, fonttitle=\bfseries\sffamily, top=6pt, bottom=6pt]
    \textbf{Doğru Cevap: C}
    \vspace{2pt}
    Analog/fiziksel verinin bilgisayarın işleyebileceği dijital formatlara (PDF, excel vb.) dönüştürülmesi sayısallaştırmadır (digitization).
\end{tcolorbox}
\vspace{10pt}

\needspace{5cm}
\vspace{8pt}
\noindent\textbf{\large\sffamily Soru 62.} \ Taranarak PDF yapılmış faturaların bir bulut yazılımına yüklenmesi ve otomatik onay süreçlerine dahil edilerek iş sürecinin hızlandırılması hangisidir?
\begin{enumerate}[label=\Alph*), leftmargin=25pt, topsep=4pt, itemsep=2pt]
    \item Sayısallaştırma (Digitization)
    \item Dijitalleşme (Digitalization)
    \item Dijital Dönüşüm (Digital Transformation)
    \item Yapay Zekâ Fabrikası
\end{enumerate}
\vspace{4pt}
\begin{tcolorbox}[title=Cevap ve Açıklama, colback=gray!3!white, colframe=primarycolor!90!black, arc=2mm, fonttitle=\bfseries\sffamily, top=6pt, bottom=6pt]
    \textbf{Doğru Cevap: B}
    \vspace{2pt}
    Sayısallaştırılmış verinin mevcut iş süreçlerini iyileştirmek, otomatikleştirmek ve optimize etmek için kullanılması dijitalleşmedir (digitalization).
\end{tcolorbox}
\vspace{10pt}

\needspace{5cm}
\vspace{8pt}
\noindent\textbf{\large\sffamily Soru 63.} \ Tüm müşteri etkileşim verilerini toplayarak tamamen veriye dayalı, kişiselleştirilmiş yeni bir yayıncılık iş modeli (Netflix vb.) oluşturmak hangisidir?
\begin{enumerate}[label=\Alph*), leftmargin=25pt, topsep=4pt, itemsep=2pt]
    \item Sayısallaştırma
    \item Dijitalleşme
    \item Dijital Dönüşüm (Digital Transformation)
    \item Sanayi 3.0
\end{enumerate}
\vspace{4pt}
\begin{tcolorbox}[title=Cevap ve Açıklama, colback=gray!3!white, colframe=primarycolor!90!black, arc=2mm, fonttitle=\bfseries\sffamily, top=6pt, bottom=6pt]
    \textbf{Doğru Cevap: C}
    \vspace{2pt}
    Teknoloji ve veri entegrasyonu ile şirketin iş modelinin, değer önerisinin ve operasyonel yapısının kökten değiştirilmesi dijital dönüşümdür.
\end{tcolorbox}
\vspace{10pt}

\needspace{5cm}
\vspace{8pt}
\noindent\textbf{\large\sffamily Soru 64.} \ Aşağıdakilerden hangisi sayısallaştırma (digitization) kavramına örnek olarak gösterilemez?
\begin{enumerate}[label=\Alph*), leftmargin=25pt, topsep=4pt, itemsep=2pt]
    \item Kağıt belgelerin taranarak PDF yapılması
    \item Analog ses kayıtlarının MP3 formatına dönüştürülmesi
    \item Depodaki ürün listesinin Excel dosyasına manuel olarak işlenmesi
    \item Müşteri şikayetlerine göre otomatik indirim kodları tanımlayan algoritma yazılması
\end{enumerate}
\vspace{4pt}
\begin{tcolorbox}[title=Cevap ve Açıklama, colback=gray!3!white, colframe=primarycolor!90!black, arc=2mm, fonttitle=\bfseries\sffamily, top=6pt, bottom=6pt]
    \textbf{Doğru Cevap: D}
    \vspace{2pt}
    Otomatik indirim kodu tanımlayan algoritma yazılması süreç iyileştirme ve otomasyon içerdiğinden dijitalleşme/dijital dönüşüm kapsamındadır, basit bir sayısallaştırma değildir.
\end{tcolorbox}
\vspace{10pt}

\needspace{5cm}
\vspace{8pt}
\noindent\textbf{\large\sffamily Soru 65.} \ Dalenogare ve ark. (2018) çalışmasına göre gelişmekte olan ülkelerde Sanayi 4.0 uygulamalarında öne çıkan kritik teknoloji hangisidir?
\begin{enumerate}[label=\Alph*), leftmargin=25pt, topsep=4pt, itemsep=2pt]
    \item Manuel veri yedekleme sistemleri
    \item Simülasyonlar ve sanal modellerin analizi
    \item Geleneksel kağıt arşivleme
    \item Standart montaj hattı robotları
\end{enumerate}
\vspace{4pt}
\begin{tcolorbox}[title=Cevap ve Açıklama, colback=gray!3!white, colframe=primarycolor!90!black, arc=2mm, fonttitle=\bfseries\sffamily, top=6pt, bottom=6pt]
    \textbf{Doğru Cevap: B}
    \vspace{2pt}
    Dalenogare ve ark. (2018) çalışmasında gelişmekte olan ülkeler için öne çıkan Sanayi 4.0 teknolojilerinden biri 'Simülasyonlar ve Sanal Modellerin Analizi'dir.
\end{tcolorbox}
\vspace{10pt}

\needspace{5cm}
\vspace{8pt}
\noindent\textbf{\large\sffamily Soru 66.} \ Üretim süreçlerinin anlık izlenmesini, takibini ve siber-fiziksel entegrasyonunu sağlayan MES (Manufacturing Execution System) hangi devrin ürünüdür?
\begin{enumerate}[label=\Alph*), leftmargin=25pt, topsep=4pt, itemsep=2pt]
    \item Sanayi 1.0
    \item Sanayi 2.0
    \item Sanayi 3.0
    \item Sanayi 4.0
\end{enumerate}
\vspace{4pt}
\begin{tcolorbox}[title=Cevap ve Açıklama, colback=gray!3!white, colframe=primarycolor!90!black, arc=2mm, fonttitle=\bfseries\sffamily, top=6pt, bottom=6pt]
    \textbf{Doğru Cevap: D}
    \vspace{2pt}
    Sanayi 4.0; bilgi, veri, siber-fiziksel sistemler ve IoT entegrasyonu (MES ve SCADA sistemleri) üzerine kuruludur.
\end{tcolorbox}
\vspace{10pt}

\needspace{5cm}
\vspace{8pt}
\noindent\textbf{\large\sffamily Soru 67.} \ Sayısallaştırılmış bilginin iş süreçlerini kolaylaştırmak ve veri üzerinden değer yaratmak için kullanılması hangisidir?
\begin{enumerate}[label=\Alph*), leftmargin=25pt, topsep=4pt, itemsep=2pt]
    \item Sayısallaştırma
    \item Dijitalleşme (Digitalization)
    \item Sanal Modelleme
    \item Kültürel Asimilasyon
\end{enumerate}
\vspace{4pt}
\begin{tcolorbox}[title=Cevap ve Açıklama, colback=gray!3!white, colframe=primarycolor!90!black, arc=2mm, fonttitle=\bfseries\sffamily, top=6pt, bottom=6pt]
    \textbf{Doğru Cevap: B}
    \vspace{2pt}
    Sayısallaştırılmış verinin süreç verimliliğini artırmak amacıyla iş süreçlerinde kullanılması dijitalleşmeyi (digitalization) tanımlar.
\end{tcolorbox}
\vspace{10pt}

\needspace{5cm}
\vspace{8pt}
\noindent\textbf{\large\sffamily Soru 68.} \ Dijital dönüşümün, sayısallaştırma ve dijitalleşmeden ayrılan en temel farkı aşağıdakilerden hangisidir?
\begin{enumerate}[label=\Alph*), leftmargin=25pt, topsep=4pt, itemsep=2pt]
    \item Sadece bilgisayar donanımı gerektirmesi
    \item Şirket kültürünü, iş modellerini ve değer zincirini kökten dönüştürmesi
    \item Daha ucuz bir yöntem olması
    \item Sadece kağıt tasarrufu sağlaması
\end{enumerate}
\vspace{4pt}
\begin{tcolorbox}[title=Cevap ve Açıklama, colback=gray!3!white, colframe=primarycolor!90!black, arc=2mm, fonttitle=\bfseries\sffamily, top=6pt, bottom=6pt]
    \textbf{Doğru Cevap: B}
    \vspace{2pt}
    Dijital dönüşüm, süreç bazlı iyileştirmelerin (dijitalleşme) ötesinde, organizasyonel kültür ve iş modelini kökten değiştiren stratejik bir süreçtir.
\end{tcolorbox}
\vspace{10pt}

\needspace{5cm}
\vspace{8pt}
\noindent\textbf{\large\sffamily Soru 69.} \ Sanayi 4.0 kapsamında kullanılan 'Sanal Modeller ve Simülasyonlar' üretim sürecinde hangi somut yararı sağlar?
\begin{enumerate}[label=\Alph*), leftmargin=25pt, topsep=4pt, itemsep=2pt]
    \item Fiziksel üretime geçilmeden önce sistemin test edilmesini ve olası hataların önlenmesini
    \item Elektrik tüketiminin tamamen sıfırlanmasını
    \item Lojistik maliyetlerinin sıfıra inmesini
    \item Müşteri anketlerinin kaldırılmasını
\end{enumerate}
\vspace{4pt}
\begin{tcolorbox}[title=Cevap ve Açıklama, colback=gray!3!white, colframe=primarycolor!90!black, arc=2mm, fonttitle=\bfseries\sffamily, top=6pt, bottom=6pt]
    \textbf{Doğru Cevap: A}
    \vspace{2pt}
    Simülasyonlar ve sanal modeller (dijital ikizler), üretim hatlarındaki hataları ve darboğazları fiziksel kuruluma geçmeden analiz etme imkanı verir.
\end{tcolorbox}
\vspace{10pt}

\needspace{5cm}
\vspace{8pt}
\noindent\textbf{\large\sffamily Soru 70.} \ Aşağıdakilerden hangisi analog verinin 0 ve 1 kodlarından oluşan dijital forma dönüştürülmesini ifade eder?
\begin{enumerate}[label=\Alph*), leftmargin=25pt, topsep=4pt, itemsep=2pt]
    \item Dijital Dönüşüm
    \item Dijitalleşme
    \item Sayısallaştırma (Digitization)
    \item Yapay Zekâ Fabrikası
\end{enumerate}
\vspace{4pt}
\begin{tcolorbox}[title=Cevap ve Açıklama, colback=gray!3!white, colframe=primarycolor!90!black, arc=2mm, fonttitle=\bfseries\sffamily, top=6pt, bottom=6pt]
    \textbf{Doğru Cevap: C}
    \vspace{2pt}
    Sayısallaştırma (digitization), analog veya fiziksel girdilerin doğrudan dijital (binary) biçime dönüştürülmesidir.
\end{tcolorbox}
\vspace{10pt}

\needspace{5cm}
\vspace{8pt}
\noindent\textbf{\large\sffamily Soru 71.} \ Analog VHS video kasetlerindeki kayıtların bilgisayara aktarılarak MP4 formatında saklanması hangi kavramla açıklanır?
\begin{enumerate}[label=\Alph*), leftmargin=25pt, topsep=4pt, itemsep=2pt]
    \item Dijital Dönüşüm
    \item Sayısallaştırma (Digitization)
    \item Dijitalleşme
    \item Sanal Gerçeklik
\end{enumerate}
\vspace{4pt}
\begin{tcolorbox}[title=Cevap ve Açıklama, colback=gray!3!white, colframe=primarycolor!90!black, arc=2mm, fonttitle=\bfseries\sffamily, top=6pt, bottom=6pt]
    \textbf{Doğru Cevap: B}
    \vspace{2pt}
    VHS kasetteki analog video verisinin bilgisayara MP4 olarak kaydedilmesi sayısallaştırma (digitization) örneğidir.
\end{tcolorbox}
\vspace{10pt}

\needspace{5cm}
\vspace{8pt}
\noindent\textbf{\large\sffamily Soru 72.} \ Hangi sanayi devrimi veri ve bilgiye olan bağımlılığı nedeniyle diğer sanayi devrimlerinden kökten ayrılmaktadır?
\begin{enumerate}[label=\Alph*), leftmargin=25pt, topsep=4pt, itemsep=2pt]
    \item Sanayi 1.0
    \item Sanayi 2.0
    \item Sanayi 3.0
    \item Sanayi 4.0
\end{enumerate}
\vspace{4pt}
\begin{tcolorbox}[title=Cevap ve Açıklama, colback=gray!3!white, colframe=primarycolor!90!black, arc=2mm, fonttitle=\bfseries\sffamily, top=6pt, bottom=6pt]
    \textbf{Doğru Cevap: D}
    \vspace{2pt}
    Dördüncü Sanayi Devrimi (Sanayi 4.0), siber-fiziksel sistemler, veri analitiği ve internet teknolojilerine bağımlı olmasıyla önceki makineleşme/elektriklendirme devrimlerinden ayrılır.
\end{tcolorbox}
\vspace{10pt}

\needspace{5cm}
\vspace{8pt}
\noindent\textbf{\large\sffamily Soru 73.} \ Fiziksel belgelerin taranıp buluta aktarılmasından sonra, bu belgelere dayalı olarak çalışanların uzaktan iş akışlarını yönetebilmesi hangisidir?
\begin{enumerate}[label=\Alph*), leftmargin=25pt, topsep=4pt, itemsep=2pt]
    \item Sayısallaştırma
    \item Dijitalleşme (Digitalization)
    \item Manuel Süreç
    \item Sanayi 1.0
\end{enumerate}
\vspace{4pt}
\begin{tcolorbox}[title=Cevap ve Açıklama, colback=gray!3!white, colframe=primarycolor!90!black, arc=2mm, fonttitle=\bfseries\sffamily, top=6pt, bottom=6pt]
    \textbf{Doğru Cevap: B}
    \vspace{2pt}
    Sayısallaştırılmış belgelerin uzaktan erişim ve iş akışı yönetimi amacıyla kullanılması bir süreç iyileştirmesidir, yani dijitalleşmedir.
\end{tcolorbox}
\vspace{10pt}

\needspace{5cm}
\vspace{8pt}
\noindent\textbf{\large\sffamily Soru 74.} \ Dördüncü Sanayi Devrimi (Sanayi 4.0) kavramı ilk kez nerede ve hangi yılda resmi olarak ortaya atılmıştır?
\begin{enumerate}[label=\Alph*), leftmargin=25pt, topsep=4pt, itemsep=2pt]
    \item ABD - Silikon Vadisi (2005)
    \item Almanya - Hannover Fuarı (2011)
    \item İsviçre - Davos Zirvesi (2015)
    \item Japonya - Tokyo Teknoloji Forumu (2018)
\end{enumerate}
\vspace{4pt}
\begin{tcolorbox}[title=Cevap ve Açıklama, colback=gray!3!white, colframe=primarycolor!90!black, arc=2mm, fonttitle=\bfseries\sffamily, top=6pt, bottom=6pt]
    \textbf{Doğru Cevap: B}
    \vspace{2pt}
    Sanayi 4.0 (Industrie 4.0) terimi, ilk olarak 2011 yılında Almanya'daki Hannover Fuarı'nda Alman hükümetinin yüksek teknoloji strateji grubu tarafından sunulmuştur.
\end{tcolorbox}
\vspace{10pt}

\needspace{5cm}
\vspace{8pt}
\noindent\textbf{\large\sffamily Soru 75.} \ Dijital dönüşümün başarıyla hayata geçirilmesinde organizasyonel düzeyde en büyük engel genellikle hangisidir?
\begin{enumerate}[label=\Alph*), leftmargin=25pt, topsep=4pt, itemsep=2pt]
    \item Bilgisayar işlemci güçlerinin yetersizliği
    \item Kültürel direnç ve değişim yönetimi eksikliği
    \item İnternet altyapısının bulunmaması
    \item Kağıt üretiminin ucuzlaması
\end{enumerate}
\vspace{4pt}
\begin{tcolorbox}[title=Cevap ve Açıklama, colback=gray!3!white, colframe=primarycolor!90!black, arc=2mm, fonttitle=\bfseries\sffamily, top=6pt, bottom=6pt]
    \textbf{Doğru Cevap: B}
    \vspace{2pt}
    Dijital dönüşüm projelerinde en büyük zorluk teknolojik olmaktan ziyade, çalışanların ve yönetimin eski alışkanlıkları sürdürme eğilimi (kültürel direnç) ve liderlik eksikliğidir.
\end{tcolorbox}
\vspace{10pt}

\needspace{5cm}
\vspace{8pt}
\noindent\textbf{\large\sffamily Soru 76.} \ Kuruluşların yeteneklerinin belirli süreçler dahilinde zaman içerisindeki gelişim durumunu tanımlayan Pullen (2007) olgunluk modelini nasıl betimler?
\begin{enumerate}[label=\Alph*), leftmargin=25pt, topsep=4pt, itemsep=2pt]
    \item Sadece teknolojik araçların listesidir
    \item Gelişim sürecinin farklı aşamalarındaki etkili işlemlerin özelliklerini tanımlayan yapılandırılmış öğeler koleksiyonu
    \item Şirketin yıllık finansal bütçesidir
    \item Sadece yazılımcıların performans kriteridir
\end{enumerate}
\vspace{4pt}
\begin{tcolorbox}[title=Cevap ve Açıklama, colback=gray!3!white, colframe=primarycolor!90!black, arc=2mm, fonttitle=\bfseries\sffamily, top=6pt, bottom=6pt]
    \textbf{Doğru Cevap: B}
    \vspace{2pt}
    Pullen (2007) olgunluk modelini, gelişim sürecinin farklı aşamalarında etkili işlemlerin özelliklerini tanımlayan yapılandırılmış öğeler koleksiyonu olarak tanımlamaktadır.
\end{tcolorbox}
\vspace{10pt}

\needspace{5cm}
\vspace{8pt}
\noindent\textbf{\large\sffamily Soru 77.} \ Kuruluşların yazılım geliştirme yeteneklerini ve olgunluk seviyelerini ölçmek amacıyla Paulk ve ark. (1993) tarafından geliştirilen model hangisidir?
\begin{enumerate}[label=\Alph*), leftmargin=25pt, topsep=4pt, itemsep=2pt]
    \item IBM Olgunluk Modeli
    \item CMM (Capability Maturity Model)
    \item PUKÖ Kalite Modeli
    \item Dalenogare Modeli
\end{enumerate}
\vspace{4pt}
\begin{tcolorbox}[title=Cevap ve Açıklama, colback=gray!3!white, colframe=primarycolor!90!black, arc=2mm, fonttitle=\bfseries\sffamily, top=6pt, bottom=6pt]
    \textbf{Doğru Cevap: B}
    \vspace{2pt}
    Paulk ve ark. (1993) yazılım süreçlerini iyileştirmek ve yeteneklerini ölçmek amacıyla Yetenek Olgunluk Modeli'ni (CMM - Capability Maturity Model) geliştirmişlerdir.
\end{tcolorbox}
\vspace{10pt}

\needspace{5cm}
\vspace{8pt}
\noindent\textbf{\large\sffamily Soru 78.} \ Deming (2018) tarafından tanımlanan Kalite Yönetim Sistemleri döngüsünün temel adımları hangi sıra ile uygulanmalıdır?
\begin{enumerate}[label=\Alph*), leftmargin=25pt, topsep=4pt, itemsep=2pt]
    \item Planla - Kontrol Et - Uygula - Önlem Al
    \item Planla - Uygula - Kontrol Et - Önlem Al (PUKÖ)
    \item Önlem Al - Planla - Uygula - Kontrol Et
    \item Uygula - Kontrol Et - Önlem Al - Planla
\end{enumerate}
\vspace{4pt}
\begin{tcolorbox}[title=Cevap ve Açıklama, colback=gray!3!white, colframe=primarycolor!90!black, arc=2mm, fonttitle=\bfseries\sffamily, top=6pt, bottom=6pt]
    \textbf{Doğru Cevap: B}
    \vspace{2pt}
    Kalite ve olgunluk süreçlerinin temeli olan PUKÖ döngüsü: Planla (Plan), Uygula (Do), Kontrol Et (Check), Önlem Al (Act) sıralamasını esas alır.
\end{tcolorbox}
\vspace{10pt}

\needspace{5cm}
\vspace{8pt}
\noindent\textbf{\large\sffamily Soru 79.} \ IBM Dijital Olgunluk Modeli'nde (Berman ve Bell, 2011) yer alan 'Süreç Entegrasyonu' boyutu aşağıdakilerden hangisini hedefler?
\begin{enumerate}[label=\Alph*), leftmargin=25pt, topsep=4pt, itemsep=2pt]
    \item Şirket içindeki tüm çalışanları işten çıkarmayı
    \item Dijital olanaklar kullanılarak süreçlerin optimize edilmesi ve şirketler arasında entegre edilmesini
    \item Şirketin web sitesinin tasarımını değiştirmeyi
    \item Sosyal medya hesaplarının kapatılmasını
\end{enumerate}
\vspace{4pt}
\begin{tcolorbox}[title=Cevap ve Açıklama, colback=gray!3!white, colframe=primarycolor!90!black, arc=2mm, fonttitle=\bfseries\sffamily, top=6pt, bottom=6pt]
    \textbf{Doğru Cevap: B}
    \vspace{2pt}
    Süreç Entegrasyonu, dijital araçlarla iş süreçlerinin verimli hale getirilmesini ve iş ortaklarıyla süreç entegrasyonu kurulmasını kapsar.
\end{tcolorbox}
\vspace{10pt}

\needspace{5cm}
\vspace{8pt}
\noindent\textbf{\large\sffamily Soru 80.} \ IBM Dijital Olgunluk Modeli'nin 'İş Modeli İnovasyonu' boyutu neye odaklanmaktadır?
\begin{enumerate}[label=\Alph*), leftmargin=25pt, topsep=4pt, itemsep=2pt]
    \item Yeni fiziksel fabrikaların kurulmasına
    \item Her seviyede müşteri değerinin oluşturulmasını hedefleyen yeni iş modelleri tasarlamaya
    \item Sadece bilgisayar sayısını artırmaya
    \item Çalışanların maaşlarını düşürmeye
\end{enumerate}
\vspace{4pt}
\begin{tcolorbox}[title=Cevap ve Açıklama, colback=gray!3!white, colframe=primarycolor!90!black, arc=2mm, fonttitle=\bfseries\sffamily, top=6pt, bottom=6pt]
    \textbf{Doğru Cevap: B}
    \vspace{2pt}
    İş Modeli İnovasyonu boyutu, dijital yeteneklerle yeni gelir modelleri ve değer önerileri oluşturarak iş modellerini yenilemeye odaklanır.
\end{tcolorbox}
\vspace{10pt}

\needspace{5cm}
\vspace{8pt}
\noindent\textbf{\large\sffamily Soru 81.} \ Aşağıdakilerden hangisi IBM Dijital Olgunluk Modeli'nde yer alan 6 ana boyuttan biri değildir?
\begin{enumerate}[label=\Alph*), leftmargin=25pt, topsep=4pt, itemsep=2pt]
    \item Analitik Görü
    \item Süreç Entegrasyonu
    \item Müşteri ve Toplum İş Birliği
    \item Tersine Lojistik ve İmhalamalar
\end{enumerate}
\vspace{4pt}
\begin{tcolorbox}[title=Cevap ve Açıklama, colback=gray!3!white, colframe=primarycolor!90!black, arc=2mm, fonttitle=\bfseries\sffamily, top=6pt, bottom=6pt]
    \textbf{Doğru Cevap: D}
    \vspace{2pt}
    Berman ve Bell (2011) tarafından geliştirilen IBM modelinde Analitik Görü, Süreç Entegrasyonu ve Müşteri İş Birliği yer alır; Tersine Lojistik bu olgunluk boyutları arasında değildir.
\end{tcolorbox}
\vspace{10pt}

\needspace{5cm}
\vspace{8pt}
\noindent\textbf{\large\sffamily Soru 82.} \ Olgunluk modelini 'bir işletmede yer alan bir varlığın zaman içindeki gelişim süreci' olarak tanımlayan akademik yazar kimdir?
\begin{enumerate}[label=\Alph*), leftmargin=25pt, topsep=4pt, itemsep=2pt]
    \item Paulk (1993)
    \item Klimko (2001)
    \item Pullen (2007)
    \item Hardin (1968)
\end{enumerate}
\vspace{4pt}
\begin{tcolorbox}[title=Cevap ve Açıklama, colback=gray!3!white, colframe=primarycolor!90!black, arc=2mm, fonttitle=\bfseries\sffamily, top=6pt, bottom=6pt]
    \textbf{Doğru Cevap: B}
    \vspace{2pt}
    Klimko (2001) olgunluk durumunu, bir işletmede yer alan fiziksel veya örgütsel bir varlığın zaman içindeki gelişim süreci olarak açıklar.
\end{tcolorbox}
\vspace{10pt}

\needspace{5cm}
\vspace{8pt}
\noindent\textbf{\large\sffamily Soru 83.} \ PUKÖ döngüsünde 'Kontrol Et' (Check) aşaması hangi faaliyeti temsil eder?
\begin{enumerate}[label=\Alph*), leftmargin=25pt, topsep=4pt, itemsep=2pt]
    \item Hedeflerin planlanması sürecini
    \item Yapılan uygulamanın planlanan hedeflerle uyumunun ve performansının ölçülmesini
    \item Hataların giderilmesi için aksiyon alınmasını
    \item Projenin tamamen iptal edilmesini
\end{enumerate}
\vspace{4pt}
\begin{tcolorbox}[title=Cevap ve Açıklama, colback=gray!3!white, colframe=primarycolor!90!black, arc=2mm, fonttitle=\bfseries\sffamily, top=6pt, bottom=6pt]
    \textbf{Doğru Cevap: B}
    \vspace{2pt}
    Kontrol Et aşaması, hedefler ile gerçekleşen sonuçların karşılaştırılması ve performansın ölçülmesidir.
\end{tcolorbox}
\vspace{10pt}

\needspace{5cm}
\vspace{8pt}
\noindent\textbf{\large\sffamily Soru 84.} \ IBM Olgunluk Modeli'nde yer alan 'Müşteri ve Toplum İş Birliği' boyutu neyi hedefler?
\begin{enumerate}[label=\Alph*), leftmargin=25pt, topsep=4pt, itemsep=2pt]
    \item Müşterilerin şirket yönetimine doğrudan atanmasını
    \item Müşteri odaklılığın tüm departmanlara yayılmasını ve sosyal ağların entegre edilmesini
    \item Müşterilere ücretsiz ürün dağıtılmasını
    \item Sadece pazarlama reklamlarının sayısını artırmayı
\end{enumerate}
\vspace{4pt}
\begin{tcolorbox}[title=Cevap ve Açıklama, colback=gray!3!white, colframe=primarycolor!90!black, arc=2mm, fonttitle=\bfseries\sffamily, top=6pt, bottom=6pt]
    \textbf{Doğru Cevap: B}
    \vspace{2pt}
    Müşteri ve Toplum İş Birliği boyutu, müşteri deneyimini geliştirmek için sosyal mecralardan ve veri ağlarından iş birliği odaklı yararlanmayı ifade eder.
\end{tcolorbox}
\vspace{10pt}

\needspace{5cm}
\vspace{8pt}
\noindent\textbf{\large\sffamily Soru 85.} \ CMMI (Yetenek Olgunluk Model Entegrasyonu) kapsamında en yüksek olgunluk seviyesi (Seviye 5) aşağıdakilerden hangisidir?
\begin{enumerate}[label=\Alph*), leftmargin=25pt, topsep=4pt, itemsep=2pt]
    \item Tanımlanmış (Defined)
    \item Yönetilen (Managed)
    \item Optimize Edilen (Optimizing)
    \item Başlangıç (Initial)
\end{enumerate}
\vspace{4pt}
\begin{tcolorbox}[title=Cevap ve Açıklama, colback=gray!3!white, colframe=primarycolor!90!black, arc=2mm, fonttitle=\bfseries\sffamily, top=6pt, bottom=6pt]
    \textbf{Doğru Cevap: C}
    \vspace{2pt}
    CMMI modelinde 5. Seviye 'Optimize Edilen' (Optimizing) seviyedir; sürekli iyileştirme ve yeniliklerin süreçlere adapte edilmesini ifade eder.
\end{tcolorbox}
\vspace{10pt}

\needspace{5cm}
\vspace{8pt}
\noindent\textbf{\large\sffamily Soru 86.} \ Yazılım olgunluk seviyelerini ölçmede Seviye 1 (Initial) neyi ifade eder?
\begin{enumerate}[label=\Alph*), leftmargin=25pt, topsep=4pt, itemsep=2pt]
    \item Süreçlerin tamamen standartlaştırıldığını
    \item Süreçlerin plansız, kaotik ve kişilere bağımlı yürütüldüğünü
    \item Sürekli optimizasyon yapıldığını
    \item Süreçlerin dış denetimden başarıyla geçtiğini
\end{enumerate}
\vspace{4pt}
\begin{tcolorbox}[title=Cevap ve Açıklama, colback=gray!3!white, colframe=primarycolor!90!black, arc=2mm, fonttitle=\bfseries\sffamily, top=6pt, bottom=6pt]
    \textbf{Doğru Cevap: B}
    \vspace{2pt}
    Olgunluk modellerinde başlangıç seviyesi (Initial/Level 1), süreçlerin tanımlanmadığı, kontrolsüz ve tamamen kahraman bireylere bağımlı olduğu aşamadır.
\end{tcolorbox}
\vspace{10pt}

\needspace{5cm}
\vspace{8pt}
\noindent\textbf{\large\sffamily Soru 87.} \ Kalite yönetim sistemlerinin babası kabul edilen ve PUKÖ döngüsünü popülerleştiren bilim insanı kimdir?
\begin{enumerate}[label=\Alph*), leftmargin=25pt, topsep=4pt, itemsep=2pt]
    \item W. Edwards Deming
    \item Joseph Juran
    \item Philip Crosby
    \item Garrett Hardin
\end{enumerate}
\vspace{4pt}
\begin{tcolorbox}[title=Cevap ve Açıklama, colback=gray!3!white, colframe=primarycolor!90!black, arc=2mm, fonttitle=\bfseries\sffamily, top=6pt, bottom=6pt]
    \textbf{Doğru Cevap: A}
    \vspace{2pt}
    Edwards Deming, sürekli kalite iyileştirme döngüsü olan PUKÖ'yü (Shewhart döngüsünden geliştirerek) popülerleştiren kalite öncüsüdür.
\end{tcolorbox}
\vspace{10pt}

\needspace{5cm}
\vspace{8pt}
\noindent\textbf{\large\sffamily Soru 88.} \ Kuruluşların dijital olgunluğunu ölçerken 'teknoloji' boyutunun yanı sıra değerlendirilmesi gereken en kritik iki boyut hangisidir?
\begin{enumerate}[label=\Alph*), leftmargin=25pt, topsep=4pt, itemsep=2pt]
    \item Finans ve Vergiler
    \item İnsan Kaynakları (Kültür) ve Süreçler
    \item Ofis binalarının sayısı ve konumu
    \item Rakip firmaların hisse değerleri
\end{enumerate}
\vspace{4pt}
\begin{tcolorbox}[title=Cevap ve Açıklama, colback=gray!3!white, colframe=primarycolor!90!black, arc=2mm, fonttitle=\bfseries\sffamily, top=6pt, bottom=6pt]
    \textbf{Doğru Cevap: B}
    \vspace{2pt}
    Dijital olgunluk sadece teknolojiyle ölçülmez; organizasyonel yapı, süreçler, liderlik ve kültür (insan) boyutlarının da değerlendirilmesi zorunludur.
\end{tcolorbox}
\vspace{10pt}

\needspace{5cm}
\vspace{8pt}
\noindent\textbf{\large\sffamily Soru 89.} \ IBM Olgunluk Modeli'nde 'Analitik Görü' seviyesi en yüksek olan bir firma kararlarını neye dayandırır?
\begin{enumerate}[label=\Alph*), leftmargin=25pt, topsep=4pt, itemsep=2pt]
    \item Yöneticilerin hislerine ve tecrübelerine
    \item Tahminleme modellerine ve ileri veri analitiğine
    \item Rakip şirketlerin yaptıkları açıklamalara
    \item Sadece geçmiş yılın ciro raporlarına
\end{enumerate}
\vspace{4pt}
\begin{tcolorbox}[title=Cevap ve Açıklama, colback=gray!3!white, colframe=primarycolor!90!black, arc=2mm, fonttitle=\bfseries\sffamily, top=6pt, bottom=6pt]
    \textbf{Doğru Cevap: B}
    \vspace{2pt}
    Analitik görü olgunluğu yüksek firmalar, sadece geçmişe bakmakla (descriptive) kalmaz, makine öğrenmesi ve yapay zeka ile geleceği tahminleyerek (predictive) karar alır.
\end{tcolorbox}
\vspace{10pt}

\needspace{5cm}
\vspace{8pt}
\noindent\textbf{\large\sffamily Soru 90.} \ PUKÖ döngüsünün 'Önlem Al' (Act) aşaması aşağıdakilerden hangisini içerir?
\begin{enumerate}[label=\Alph*), leftmargin=25pt, topsep=4pt, itemsep=2pt]
    \item Performans ölçüm sonuçlarına göre kalıcı iyileştirmelerin standartlaştırılmasını ve yaygınlaştırılmasını
    \item Sadece projenin bütçesini kontrol etmeyi
    \item Yeni bir proje fikri tasarlamayı
    \item Rakip firmaları analiz etmeyi
\end{enumerate}
\vspace{4pt}
\begin{tcolorbox}[title=Cevap ve Açıklama, colback=gray!3!white, colframe=primarycolor!90!black, arc=2mm, fonttitle=\bfseries\sffamily, top=6pt, bottom=6pt]
    \textbf{Doğru Cevap: A}
    \vspace{2pt}
    Önlem Al aşaması, test edilen ve başarılı olan iyileştirmelerin kurumsal standart haline getirilmesini ve hataların kalıcı olarak çözülmesini kapsar.
\end{tcolorbox}
\vspace{10pt}

\needspace{5cm}
\vspace{8pt}
\noindent\textbf{\large\sffamily Soru 91.} \ Türkiye'de sanayide dijital dönüşümün pilot uygulama süreçlerinin yürütülmesi görevi hangi bakanlık birimine verilmiştir?
\begin{enumerate}[label=\Alph*), leftmargin=25pt, topsep=4pt, itemsep=2pt]
    \item Ar-Ge Teşvikleri Genel Müdürlüğü
    \item Sanayi ve Verimlilik Genel Müdürlüğü
    \item KOSGEB Başkanlığı
    \item TÜBİTAK TEYDEB
\end{enumerate}
\vspace{4pt}
\begin{tcolorbox}[title=Cevap ve Açıklama, colback=gray!3!white, colframe=primarycolor!90!black, arc=2mm, fonttitle=\bfseries\sffamily, top=6pt, bottom=6pt]
    \textbf{Doğru Cevap: B}
    \vspace{2pt}
    Türkiye'deki sanayi dijitalleşme politikalarında 'Pilot Uygulama Süreçleri' Sanayi ve Verimlilik Genel Müdürlüğü uhdesindedir.
\end{tcolorbox}
\vspace{10pt}

\needspace{5cm}
\vspace{8pt}
\noindent\textbf{\large\sffamily Soru 92.} \ Türkiye'de dijitalleşme stratejilerini geliştirmek üzere TOBB, TİM, TÜSİAD, MÜSİAD, YASED ve TTGV katılımıyla kurulan çalışma gruplarının sayısı kaçtır?
\begin{enumerate}[label=\Alph*), leftmargin=25pt, topsep=4pt, itemsep=2pt]
    \item 3
    \item 6
    \item 10
    \item 12
\end{enumerate}
\vspace{4pt}
\begin{tcolorbox}[title=Cevap ve Açıklama, colback=gray!3!white, colframe=primarycolor!90!black, arc=2mm, fonttitle=\bfseries\sffamily, top=6pt, bottom=6pt]
    \textbf{Doğru Cevap: B}
    \vspace{2pt}
    Sanayinin dijital dönüşümü yol haritasını belirlemek amacıyla iş dünyası temsilcileriyle birlikte 6 ana çalışma grubu kurulmuştur.
\end{tcolorbox}
\vspace{10pt}

\needspace{5cm}
\vspace{8pt}
\noindent\textbf{\large\sffamily Soru 93.} \ Sanayide dijital dönüşüm kapsamında kurulan 6 çalışma grubundan hangisi dijital dönüşümün insan kaynağı boyutunu ele almaktadır?
\begin{enumerate}[label=\Alph*), leftmargin=25pt, topsep=4pt, itemsep=2pt]
    \item Altyapı Çalışma Grubu
    \item Eğitim Çalışma Grubu
    \item Açık İnovasyon Çalışma Grubu
    \item Mevzuat Çalışma Grubu
\end{enumerate}
\vspace{4pt}
\begin{tcolorbox}[title=Cevap ve Açıklama, colback=gray!3!white, colframe=primarycolor!90!black, arc=2mm, fonttitle=\bfseries\sffamily, top=6pt, bottom=6pt]
    \textbf{Doğru Cevap: B}
    \vspace{2pt}
    Eğitim Çalışma Grubu, dijital dönüşümün gerektirdiği nitelikli iş gücü, mühendislik ve teknik yetkinliklerin kazandırılmasını koordine eder.
\end{tcolorbox}
\vspace{10pt}

\needspace{5cm}
\vspace{8pt}
\noindent\textbf{\large\sffamily Soru 94.} \ Türkiye'de sanayicilere uygulamalı yalın üretim ve dijital dönüşüm eğitimleri vermek amacıyla ilk olarak Ankara ve Bursa'da kurulan merkezlere ne ad verilir?
\begin{enumerate}[label=\Alph*), leftmargin=25pt, topsep=4pt, itemsep=2pt]
    \item Teknoloji Geliştirme Bölgeleri
    \item Dijital Model Fabrikalar
    \item KOSGEB Laboratuvarları
    \item Organize Sanayi Merkezleri
\end{enumerate}
\vspace{4pt}
\begin{tcolorbox}[title=Cevap ve Açıklama, colback=gray!3!white, colframe=primarycolor!90!black, arc=2mm, fonttitle=\bfseries\sffamily, top=6pt, bottom=6pt]
    \textbf{Doğru Cevap: B}
    \vspace{2pt}
    Model Fabrikalar (Yetkinlik ve Dijital Dönüşüm Merkezleri), sanayicilere yaşayarak öğrenme ve dijitalleşme yetenekleri kazandıran pilot merkezlerdir.
\end{tcolorbox}
\vspace{10pt}

\needspace{5cm}
\vspace{8pt}
\noindent\textbf{\large\sffamily Soru 95.} \ Türkiye'de kamunun ve sanayinin dijital dönüşüm koordinasyonunu sağlamak amacıyla Cumhurbaşkanlığı bünyesinde kurulan ofis hangisidir?
\begin{enumerate}[label=\Alph*), leftmargin=25pt, topsep=4pt, itemsep=2pt]
    \item İnsan Kaynakları Ofisi
    \item Dijital Dönüşüm Ofisi
    \item Yatırım Ofisi
    \item Finans Ofisi
\end{enumerate}
\vspace{4pt}
\begin{tcolorbox}[title=Cevap ve Açıklama, colback=gray!3!white, colframe=primarycolor!90!black, arc=2mm, fonttitle=\bfseries\sffamily, top=6pt, bottom=6pt]
    \textbf{Doğru Cevap: B}
    \vspace{2pt}
    Cumhurbaşkanlığı Dijital Dönüşüm Ofisi, Türkiye'nin e-devlet, siber güvenlik, yapay zeka ve dijital teknolojilerdeki vizyonunu ve dönüşümünü koordine eder.
\end{tcolorbox}
\vspace{10pt}

\needspace{5cm}
\vspace{8pt}
\noindent\textbf{\large\sffamily Soru 96.} \ Aşağıdakilerden hangisi Türkiye'de dijitalleşme projelerinin finansmanında rol alan resmi kaynak kurumlardan biri değildir?
\begin{enumerate}[label=\Alph*), leftmargin=25pt, topsep=4pt, itemsep=2pt]
    \item TÜBİTAK
    \item KOSGEB
    \item AB ve Dış İlişkiler Genel Müdürlüğü
    \item Türkiye Odalar ve Borsalar Birliği (TOBB)
\end{enumerate}
\vspace{4pt}
\begin{tcolorbox}[title=Cevap ve Açıklama, colback=gray!3!white, colframe=primarycolor!90!black, arc=2mm, fonttitle=\bfseries\sffamily, top=6pt, bottom=6pt]
    \textbf{Doğru Cevap: D}
    \vspace{2pt}
    TOBB bir finansman sağlayıcı resmi devlet bütçe mercii değil; iş dünyasını temsil eden bir sivil toplum kuruluşu/meslek örgütüdür. Finansal kaynaklar TÜBİTAK, KOSGEB ve Bakanlık fonlarıdır.
\end{tcolorbox}
\vspace{10pt}

\needspace{5cm}
\vspace{8pt}
\noindent\textbf{\large\sffamily Soru 97.} \ Türkiye'de sanayinin dijital olgunluk düzeyini ölçmek amacıyla işletmelere yönelik yapılan somut çalışma hangisidir?
\begin{enumerate}[label=\Alph*), leftmargin=25pt, topsep=4pt, itemsep=2pt]
    \item Vergi denetimleri yapılması
    \item Sanayide Dijital Olgunluk Seviyesi anket ve ölçümlerinin gerçekleştirilmesi
    \item Tüm fabrikalara bilgisayar dağıtılması
    \item İthalat kotası konulması
\end{enumerate}
\vspace{4pt}
\begin{tcolorbox}[title=Cevap ve Açıklama, colback=gray!3!white, colframe=primarycolor!90!black, arc=2mm, fonttitle=\bfseries\sffamily, top=6pt, bottom=6pt]
    \textbf{Doğru Cevap: B}
    \vspace{2pt}
    Bakanlık ve paydaşlar tarafından sanayinin mevcut durumunu analiz etmek için 'Sanayide Dijital Olgunluk Seviyesi' anket ve saha çalışmaları yapılmıştır.
\end{tcolorbox}
\vspace{10pt}

\needspace{5cm}
\vspace{8pt}
\noindent\textbf{\large\sffamily Soru 98.} \ Sanayide Dijital Dönüşüm Platformu kapsamında kurulan 6 çalışma grubundan hangisi standartların uyumlulaştırılması ve yasal süreçlerle ilgilenir?
\begin{enumerate}[label=\Alph*), leftmargin=25pt, topsep=4pt, itemsep=2pt]
    \item Altyapı Çalışma Grubu
    \item Mevzuat ve Standartlar Çalışma Grubu
    \item Eğitim Çalışma Grubu
    \item Dijital Teknolojiler Çalışma Grubu
\end{enumerate}
\vspace{4pt}
\begin{tcolorbox}[title=Cevap ve Açıklama, colback=gray!3!white, colframe=primarycolor!90!black, arc=2mm, fonttitle=\bfseries\sffamily, top=6pt, bottom=6pt]
    \textbf{Doğru Cevap: B}
    \vspace{2pt}
    Mevzuat ve Standartlar Çalışma Grubu, siber güvenlik standartları, veri paylaşımı regülasyonları ve yasal altyapıyla ilgilenir.
\end{tcolorbox}
\vspace{10pt}

\needspace{5cm}
\vspace{8pt}
\noindent\textbf{\large\sffamily Soru 99.} \ Türkiye'de Model Fabrika uygulamalarında katılımcılara öğretilen temel üretim felsefesi aşağıdakilerden hangisidir?
\begin{enumerate}[label=\Alph*), leftmargin=25pt, topsep=4pt, itemsep=2pt]
    \item Seri ve Tek Tip Üretim
    \item Yalın Üretim (Lean Production)
    \item Manuel El Sanatları Üretimi
    \item Stok Odaklı İtme Üretimi
\end{enumerate}
\vspace{4pt}
\begin{tcolorbox}[title=Cevap ve Açıklama, colback=gray!3!white, colframe=primarycolor!90!black, arc=2mm, fonttitle=\bfseries\sffamily, top=6pt, bottom=6pt]
    \textbf{Doğru Cevap: B}
    \vspace{2pt}
    Model Fabrikalarda öncelikle israfların önlenmesini hedefleyen 'Yalın Üretim' (Lean) öğretilir, ardından bu yalın süreçler dijitalleştirilir.
\end{tcolorbox}
\vspace{10pt}

\needspace{5cm}
\vspace{8pt}
\noindent\textbf{\large\sffamily Soru 100.} \ Türkiye'nin yapay zeka alanındaki stratejik hedeflerini belirleyen ve 2021 yılında yayınlanan resmi belgenin adı nedir?
\begin{enumerate}[label=\Alph*), leftmargin=25pt, topsep=4pt, itemsep=2pt]
    \item Ulusal Yapay Zeka Stratejisi (2021-2025)
    \item Dijital Yapay Zeka Kanunu
    \item Sanayi 4.0 Eylem Planı
    \item Kalkınma Yapay Zeka Planı
\end{enumerate}
\vspace{4pt}
\begin{tcolorbox}[title=Cevap ve Açıklama, colback=gray!3!white, colframe=primarycolor!90!black, arc=2mm, fonttitle=\bfseries\sffamily, top=6pt, bottom=6pt]
    \textbf{Doğru Cevap: A}
    \vspace{2pt}
    Sanayi ve Teknoloji Bakanlığı ile Cumhurbaşkanlığı Dijital Dönüşüm Ofisi iş birliğiyle hazırlanan 'Ulusal Yapay Zeka Stratejisi (2021-2025)' Türkiye'nin resmi yapay zeka yol haritasıdır.
\end{tcolorbox}
\vspace{10pt}

\needspace{5cm}
\vspace{8pt}
\noindent\textbf{\large\sffamily Soru 101.} \ Türkiye'de dijitalleşme çalışmalarında 'Politika Geliştirme Süreci' resmi olarak hangi birimin sorumluluğundadır?
\begin{enumerate}[label=\Alph*), leftmargin=25pt, topsep=4pt, itemsep=2pt]
    \item TÜBİTAK Başkanlığı
    \item Sanayi ve Teknoloji Bakanlığı Ar-Ge Teşvikleri Genel Müdürlüğü
    \item KOSGEB Destek Dairesi
    \item Ticaret Bakanlığı
\end{enumerate}
\vspace{4pt}
\begin{tcolorbox}[title=Cevap ve Açıklama, colback=gray!3!white, colframe=primarycolor!90!black, arc=2mm, fonttitle=\bfseries\sffamily, top=6pt, bottom=6pt]
    \textbf{Doğru Cevap: B}
    \vspace{2pt}
    Dijitalleşme politikalarının resmi koordinasyon ve politika geliştirme bacağı Ar-Ge Teşvikleri Genel Müdürlüğü tarafından yürütülmektedir.
\end{tcolorbox}
\vspace{10pt}

\needspace{5cm}
\vspace{8pt}
\noindent\textbf{\large\sffamily Soru 102.} \ Türkiye'deki Model Fabrikalar projesinde devlet kurumlarının yanı sıra iş birliği yapılan uluslararası kalkınma kuruluşu hangisidir?
\begin{enumerate}[label=\Alph*), leftmargin=25pt, topsep=4pt, itemsep=2pt]
    \item Birleşmiş Milletler Kalkınma Programı (UNDP)
    \item Dünya Bankası (World Bank)
    \item IMF
    \item UNESCO
\end{enumerate}
\vspace{4pt}
\begin{tcolorbox}[title=Cevap ve Açıklama, colback=gray!3!white, colframe=primarycolor!90!black, arc=2mm, fonttitle=\bfseries\sffamily, top=6pt, bottom=6pt]
    \textbf{Doğru Cevap: A}
    \vspace{2pt}
    Model Fabrikaların kurulumu ve yaygınlaştırılması süreci Türkiye Cumhuriyeti Sanayi Bakanlığı ile UNDP (Birleşmiş Milletler Kalkınma Programı) iş birliğiyle yürütülmektedir.
\end{tcolorbox}
\vspace{10pt}

\needspace{5cm}
\vspace{8pt}
\noindent\textbf{\large\sffamily Soru 103.} \ Aşağıdakilerden hangisi Türkiye'nin sanayide dijital dönüşüm yol haritasında kurulan 6 çalışma grubundan biri değildir?
\begin{enumerate}[label=\Alph*), leftmargin=25pt, topsep=4pt, itemsep=2pt]
    \item İleri Üretim Teknikleri
    \item Dijital Teknolojiler
    \item Fosil Yakıt Teşvikleri
    \item Açık İnovasyon
\end{enumerate}
\vspace{4pt}
\begin{tcolorbox}[title=Cevap ve Açıklama, colback=gray!3!white, colframe=primarycolor!90!black, arc=2mm, fonttitle=\bfseries\sffamily, top=6pt, bottom=6pt]
    \textbf{Doğru Cevap: C}
    \vspace{2pt}
    Fosil Yakıt Teşvikleri gibi çevreye zararlı bir başlık çalışma grupları arasında yer almaz. Çalışma grupları ileri teknoloji ve sürdürülebilirlik odaklıdır.
\end{tcolorbox}
\vspace{10pt}

\needspace{5cm}
\vspace{8pt}
\noindent\textbf{\large\sffamily Soru 104.} \ Türkiye'nin dijital dönüşüm yetkinliğini artırmak için 'bilgi sermayesi' oluşturmak adına Bakanlık kaç üniversiteyle iş birliği protokolü imzalamıştır?
\begin{enumerate}[label=\Alph*), leftmargin=25pt, topsep=4pt, itemsep=2pt]
    \item 5
    \item 10
    \item 15
    \item 50
\end{enumerate}
\vspace{4pt}
\begin{tcolorbox}[title=Cevap ve Açıklama, colback=gray!3!white, colframe=primarycolor!90!black, arc=2mm, fonttitle=\bfseries\sffamily, top=6pt, bottom=6pt]
    \textbf{Doğru Cevap: C}
    \vspace{2pt}
    Dijital Türkiye yol haritasında nitelikli insan kaynağı ve bilgi birikimi için 15 üniversite ile stratejik iş birliği protokolleri imzalanmıştır.
\end{tcolorbox}
\vspace{10pt}

\needspace{5cm}
\vspace{8pt}
\noindent\textbf{\large\sffamily Soru 105.} \ Türkiye'de KOSGEB aracılığıyla KOBİ'lerin dijitalleşmesini desteklemek amacıyla sunulan popüler hibe ve destek programının adı hangisidir?
\begin{enumerate}[label=\Alph*), leftmargin=25pt, topsep=4pt, itemsep=2pt]
    \item KOBİGEL - KOBİ Gelişim Destek Programı
    \item Tarımsal Kalkınma Programı
    \item Esnaf Kredi Desteği
    \item İthalat İkame Programı
\end{enumerate}
\vspace{4pt}
\begin{tcolorbox}[title=Cevap ve Açıklama, colback=gray!3!white, colframe=primarycolor!90!black, arc=2mm, fonttitle=\bfseries\sffamily, top=6pt, bottom=6pt]
    \textbf{Doğru Cevap: A}
    \vspace{2pt}
    KOSGEB, KOBİGEL programı kapsamında 'İmalat Sanayinde Dijitalleşme' başlıklarıyla KOBİ'lerin yazılım, donanım ve teknik dönüşümlerine hibe ve destek sağlamaktadır.
\end{tcolorbox}
\vspace{10pt}

\needspace{5cm}
\vspace{8pt}
\noindent\textbf{\large\sffamily Soru 106.} \ Modern iş dünyasında verinin önemini vurgulamak için sıkça kullanılan 'Veri yeni petroldür' (Data is the new oil) sözü kime aittir?
\begin{enumerate}[label=\Alph*), leftmargin=25pt, topsep=4pt, itemsep=2pt]
    \item Bill Gates
    \item Clive Humby
    \item Steve Jobs
    \item Garrett Hardin
\end{enumerate}
\vspace{4pt}
\begin{tcolorbox}[title=Cevap ve Açıklama, colback=gray!3!white, colframe=primarycolor!90!black, arc=2mm, fonttitle=\bfseries\sffamily, top=6pt, bottom=6pt]
    \textbf{Doğru Cevap: B}
    \vspace{2pt}
    Verinin işlenerek anlamlı hale getirilmesiyle değer kazanacağını ifade eden 'Veri yeni petroldür' sözü ilk kez 2006'da İngiliz matematikçi Clive Humby tarafından söylenmiştir.
\end{tcolorbox}
\vspace{10pt}

\needspace{5cm}
\vspace{8pt}
\noindent\textbf{\large\sffamily Soru 107.} \ Büyük Veri (Big Data) analizinde, verinin metin, ses, video veya sensör sinyali gibi farklı formatlarda bulunmasını ifade eden V hangisidir?
\begin{enumerate}[label=\Alph*), leftmargin=25pt, topsep=4pt, itemsep=2pt]
    \item Volume (Hacim)
    \item Velocity (Hız)
    \item Variety (Çeşitlilik)
    \item Veracity (Doğruluk)
\end{enumerate}
\vspace{4pt}
\begin{tcolorbox}[title=Cevap ve Açıklama, colback=gray!3!white, colframe=primarycolor!90!black, arc=2mm, fonttitle=\bfseries\sffamily, top=6pt, bottom=6pt]
    \textbf{Doğru Cevap: C}
    \vspace{2pt}
    Variety (Çeşitlilik), büyük veriyi oluşturan kaynakların yapılandırılmış (SQL tabloları) veya yapılandırılmamış (video, ses, sosyal medya metinleri) olarak çeşitlilik göstermesidir.
\end{tcolorbox}
\vspace{10pt}

\needspace{5cm}
\vspace{8pt}
\noindent\textbf{\large\sffamily Soru 108.} \ Büyük Veri analitiğinde 'Kök neden analizi' yaparak 'Bir olay neden gerçekleşti?' sorusuna yanıt arayan seviye hangisidir?
\begin{enumerate}[label=\Alph*), leftmargin=25pt, topsep=4pt, itemsep=2pt]
    \item Tanımlayıcı Analitik (Descriptive)
    \item Tanısal Analitik (Diagnostic)
    \item Öngörücü Analitik (Predictive)
    \item Yönlendirici Analitik (Prescriptive)
\end{enumerate}
\vspace{4pt}
\begin{tcolorbox}[title=Cevap ve Açıklama, colback=gray!3!white, colframe=primarycolor!90!black, arc=2mm, fonttitle=\bfseries\sffamily, top=6pt, bottom=6pt]
    \textbf{Doğru Cevap: B}
    \vspace{2pt}
    Tanısal analitik (Diagnostic), geçmiş verileri inceleyerek sistemlerdeki bir arızanın, satış düşüşünün veya olayın neden kaynaklandığını bulmaya çalışır.
\end{tcolorbox}
\vspace{10pt}

\needspace{5cm}
\vspace{8pt}
\noindent\textbf{\large\sffamily Soru 109.} \ Büyük Veri analitiğinde 'Gelecekte ne olacak?' sorusunu makine öğrenmesi modelleriyle tahmin etmeye çalışan seviye hangisidir?
\begin{enumerate}[label=\Alph*), leftmargin=25pt, topsep=4pt, itemsep=2pt]
    \item Tanımlayıcı Analitik
    \item Tanısal Analitik
    \item Öngörücü Analitik (Predictive)
    \item Yönlendirici Analitik
\end{enumerate}
\vspace{4pt}
\begin{tcolorbox}[title=Cevap ve Açıklama, colback=gray!3!white, colframe=primarycolor!90!black, arc=2mm, fonttitle=\bfseries\sffamily, top=6pt, bottom=6pt]
    \textbf{Doğru Cevap: C}
    \vspace{2pt}
    Öngörücü analitik (Predictive), geçmiş veri kalıplarından yararlanarak gelecekteki talebi, arıza ihtimalini veya müşteri davranışlarını tahmin eder.
\end{tcolorbox}
\vspace{10pt}

\needspace{5cm}
\vspace{8pt}
\noindent\textbf{\large\sffamily Soru 110.} \ Siemens tarafından geliştirilen ve 1.3 milyondan fazla bağlı cihazdan gelen büyük veriyi analiz eden bulut tabanlı IoT platformunun adı nedir?
\begin{enumerate}[label=\Alph*), leftmargin=25pt, topsep=4pt, itemsep=2pt]
    \item Watson AI
    \item MindSphere
    \item Data Lake
    \item Predix
\end{enumerate}
\vspace{4pt}
\begin{tcolorbox}[title=Cevap ve Açıklama, colback=gray!3!white, colframe=primarycolor!90!black, arc=2mm, fonttitle=\bfseries\sffamily, top=6pt, bottom=6pt]
    \textbf{Doğru Cevap: B}
    \vspace{2pt}
    Siemens'in endüstriyel bulut ve IoT işletim sistemi olan MindSphere, milyonlarca bağlı makineden veri toplayıp optimizasyon sağlar.
\end{tcolorbox}
\vspace{10pt}

\needspace{5cm}
\vspace{8pt}
\noindent\textbf{\large\sffamily Soru 111.} \ Arçelik bünyesinde kalite zekası, akıllı servis yönlendirmeleri ve üretim verimliliği için oluşturulan ortak veri havuzu projesine ne ad verilir?
\begin{enumerate}[label=\Alph*), leftmargin=25pt, topsep=4pt, itemsep=2pt]
    \item MindSphere
    \item Data Lake (Veri Gölü)
    \item Hadoop Server
    \item Enterprise Cloud
\end{enumerate}
\vspace{4pt}
\begin{tcolorbox}[title=Cevap ve Açıklama, colback=gray!3!white, colframe=primarycolor!90!black, arc=2mm, fonttitle=\bfseries\sffamily, top=6pt, bottom=6pt]
    \textbf{Doğru Cevap: B}
    \vspace{2pt}
    Arçelik, farklı departmanlardan akan yapısal ve yapısal olmayan verileri tek bir yerde analiz etmek için 'Data Lake' (Veri Gölü) mimarisini kurmuştur.
\end{tcolorbox}
\vspace{10pt}

\needspace{5cm}
\vspace{8pt}
\noindent\textbf{\large\sffamily Soru 112.} \ Büyük Verinin 3V'sinden 'Velocity' (Hız) aşağıdakilerden hangisini tanımlar?
\begin{enumerate}[label=\Alph*), leftmargin=25pt, topsep=4pt, itemsep=2pt]
    \item Verilerin gigabayt veya terabayt cinsinden boyutunu
    \item Verilerin ne kadar hızlı üretildiği ve ne kadar hızlı işlenmesi gerektiğini
    \item Verilerin saklandığı sabit disklerin dönüş hızını
    \item İnternet kablolarının veri iletim hızını
\end{enumerate}
\vspace{4pt}
\begin{tcolorbox}[title=Cevap ve Açıklama, colback=gray!3!white, colframe=primarycolor!90!black, arc=2mm, fonttitle=\bfseries\sffamily, top=6pt, bottom=6pt]
    \textbf{Doğru Cevap: B}
    \vspace{2pt}
    Velocity (Hız), verinin üretim sıklığı ve gerçek zamanlı analiz edilme hızını (örneğin anlık finansal işlemler veya sensör verileri) temsil eder.
\end{tcolorbox}
\vspace{10pt}

\needspace{5cm}
\vspace{8pt}
\noindent\textbf{\large\sffamily Soru 113.} \ Şirketlerde veri ambarlarının ve veriden değer üretme süreçlerinin yönetiminden sorumlu olan üst düzey liderlik pozisyonu hangisidir?
\begin{enumerate}[label=\Alph*), leftmargin=25pt, topsep=4pt, itemsep=2pt]
    \item CFO (Finans Başkanı)
    \item CDO (Chief Data Officer / Veri Başkanı)
    \item CMO (Pazarlama Başkanı)
    \item CEO (Genel Müdür)
\end{enumerate}
\vspace{4pt}
\begin{tcolorbox}[title=Cevap ve Açıklama, colback=gray!3!white, colframe=primarycolor!90!black, arc=2mm, fonttitle=\bfseries\sffamily, top=6pt, bottom=6pt]
    \textbf{Doğru Cevap: B}
    \vspace{2pt}
    CDO (Chief Data Officer), verinin toplanması, kalitesi, analitiği ve stratejik kararlara dönüştürülmesinden sorumlu üst düzey yöneticidir.
\end{tcolorbox}
\vspace{10pt}

\needspace{5cm}
\vspace{8pt}
\noindent\textbf{\large\sffamily Soru 114.} \ Büyük Veri analitiğinde 'Ne yapmalıyız?' sorusuna yanıt veren, kararı optimize eden ve en iyi aksiyon planını sunan en üst analitik seviyesi hangisidir?
\begin{enumerate}[label=\Alph*), leftmargin=25pt, topsep=4pt, itemsep=2pt]
    \item Tanımlayıcı Analitik
    \item Tanısal Analitik
    \item Öngörücü Analitik
    \item Yönlendirici Analitik (Prescriptive)
\end{enumerate}
\vspace{4pt}
\begin{tcolorbox}[title=Cevap ve Açıklama, colback=gray!3!white, colframe=primarycolor!90!black, arc=2mm, fonttitle=\bfseries\sffamily, top=6pt, bottom=6pt]
    \textbf{Doğru Cevap: D}
    \vspace{2pt}
    Yönlendirici analitik (Prescriptive), sadece tahminde bulunmaz; simülasyon ve optimizasyon algoritmalarıyla en doğru kararın ne olması gerektiğini önerir.
\end{tcolorbox}
\vspace{10pt}

\needspace{5cm}
\vspace{8pt}
\noindent\textbf{\large\sffamily Soru 115.} \ Büyük veri platformlarında sıklıkla kullanılan ve verileri satır/sütun kısıtlaması olmadan ham biçimde depolayan havuzlara ne ad verilir?
\begin{enumerate}[label=\Alph*), leftmargin=25pt, topsep=4pt, itemsep=2pt]
    \item İlişkisel Veritabanı (SQL)
    \item Veri Gölü (Data Lake)
    \item Sunucu Diski
    \item Excel Sayfası
\end{enumerate}
\vspace{4pt}
\begin{tcolorbox}[title=Cevap ve Açıklama, colback=gray!3!white, colframe=primarycolor!90!black, arc=2mm, fonttitle=\bfseries\sffamily, top=6pt, bottom=6pt]
    \textbf{Doğru Cevap: B}
    \vspace{2pt}
    Veri Gölü (Data Lake), yapısal (relational), yarı yapısal (JSON, XML) ve yapısal olmayan (video, ses, PDF) ham verileri devasa boyutta saklayan havuzdur.
\end{tcolorbox}
\vspace{10pt}

\needspace{5cm}
\vspace{8pt}
\noindent\textbf{\large\sffamily Soru 116.} \ Büyük Verinin temel özelliklerine sonradan eklenen ve verinin 'güvenilirliği ve doğruluğu' ile ilgili olan V hangisidir?
\begin{enumerate}[label=\Alph*), leftmargin=25pt, topsep=4pt, itemsep=2pt]
    \item Volume
    \item Variety
    \item Veracity (Doğruluk/Güvenilirlik)
    \item Value
\end{enumerate}
\vspace{4pt}
\begin{tcolorbox}[title=Cevap ve Açıklama, colback=gray!3!white, colframe=primarycolor!90!black, arc=2mm, fonttitle=\bfseries\sffamily, top=6pt, bottom=6pt]
    \textbf{Doğru Cevap: C}
    \vspace{2pt}
    Veracity (Doğruluk/Güvenilirlik), toplanan büyük verinin ne kadar temiz, gürültüsüz ve analiz için güvenilir olduğunu ifade eder.
\end{tcolorbox}
\vspace{10pt}

\needspace{5cm}
\vspace{8pt}
\noindent\textbf{\large\sffamily Soru 117.} \ Aşağıdakilerden hangisi tanımlayıcı (descriptive) analitiğe örnek olarak gösterilebilir?
\begin{enumerate}[label=\Alph*), leftmargin=25pt, topsep=4pt, itemsep=2pt]
    \item Geçen ayın elektrik tüketim raporunun incelenmesi
    \item Fabrikadaki bir makinenin 3 gün sonra arızalanacağının tahmin edilmesi
    \item En az karbon salınımı yapacak lojistik rotanın otomatik hesaplanması
    \item Müşteri kaybının (churn) nedenlerinin regresyonla bulunması
\end{enumerate}
\vspace{4pt}
\begin{tcolorbox}[title=Cevap ve Açıklama, colback=gray!3!white, colframe=primarycolor!90!black, arc=2mm, fonttitle=\bfseries\sffamily, top=6pt, bottom=6pt]
    \textbf{Doğru Cevap: A}
    \vspace{2pt}
    Geçen ayın elektrik tüketim raporu geçmişte 'Ne oldu?' sorusuna cevap verdiğinden tanımlayıcı (descriptive) analitik örneğidir.
\end{tcolorbox}
\vspace{10pt}

\needspace{5cm}
\vspace{8pt}
\noindent\textbf{\large\sffamily Soru 118.} \ Sensörlerden akan sıcaklık verilerini anlık izleyip 'Sıcaklık 90 dereceye ulaştı' uyarısı veren bir IoT sistemi hangi analitik seviyesindedir?
\begin{enumerate}[label=\Alph*), leftmargin=25pt, topsep=4pt, itemsep=2pt]
    \item Tanımlayıcı Analitik (Descriptive)
    \item Öngörücü Analitik (Predictive)
    \item Yönlendirici Analitik (Prescriptive)
    \item Zayıf Yapay Zeka
\end{enumerate}
\vspace{4pt}
\begin{tcolorbox}[title=Cevap ve Açıklama, colback=gray!3!white, colframe=primarycolor!90!black, arc=2mm, fonttitle=\bfseries\sffamily, top=6pt, bottom=6pt]
    \textbf{Doğru Cevap: A}
    \vspace{2pt}
    Anlık durumu veya geçmiş durumu raporlayan ve gösteren sistemler tanımlayıcı (descriptive) seviyededir.
\end{tcolorbox}
\vspace{10pt}

\needspace{5cm}
\vspace{8pt}
\noindent\textbf{\large\sffamily Soru 119.} \ Hava tahmini verilerini, geçmiş tarım verimini ve toprak nemini analiz ederek çiftçilere 'Bu hafta ekim yapmalısınız' tavsiyesi veren sistem hangisidir?
\begin{enumerate}[label=\Alph*), leftmargin=25pt, topsep=4pt, itemsep=2pt]
    \item Tanımlayıcı Analitik
    \item Tanısal Analitik
    \item Yönlendirici Analitik (Prescriptive)
    \item Sayısallaştırma
\end{enumerate}
\vspace{4pt}
\begin{tcolorbox}[title=Cevap ve Açıklama, colback=gray!3!white, colframe=primarycolor!90!black, arc=2mm, fonttitle=\bfseries\sffamily, top=6pt, bottom=6pt]
    \textbf{Doğru Cevap: C}
    \vspace{2pt}
    Sistem karar vericiye en doğru aksiyonu tavsiye ettiği için yönlendirici analitik (Prescriptive) seviyesindedir.
\end{tcolorbox}
\vspace{10pt}

\needspace{5cm}
\vspace{8pt}
\noindent\textbf{\large\sffamily Soru 120.} \ 2025 yılı için öngörülen 175 Zettabyte veri hacmi, büyük verinin 3V prensibinden en çok hangisiyle ilişkilidir?
\begin{enumerate}[label=\Alph*), leftmargin=25pt, topsep=4pt, itemsep=2pt]
    \item Velocity (Hız)
    \item Variety (Çeşitlilik)
    \item Volume (Hacim)
    \item Veracity (Güvenilirlik)
\end{enumerate}
\vspace{4pt}
\begin{tcolorbox}[title=Cevap ve Açıklama, colback=gray!3!white, colframe=primarycolor!90!black, arc=2mm, fonttitle=\bfseries\sffamily, top=6pt, bottom=6pt]
    \textbf{Doğru Cevap: C}
    \vspace{2pt}
    Üretilen verinin boyutsal büyüklüğü (Zettabyte düzeyleri) doğrudan Volume (Hacim) bileşeniyle ilgilidir.
\end{tcolorbox}
\vspace{10pt}

\needspace{5cm}
\vspace{8pt}
\noindent\textbf{\large\sffamily Soru 121.} \ Veri yönetiminde 'Veri Ambarı' (Data Warehouse) ile 'Veri Gölü' (Data Lake) arasındaki en temel fark hangisidir?
\begin{enumerate}[label=\Alph*), leftmargin=25pt, topsep=4pt, itemsep=2pt]
    \item Veri ambarında sadece videolar saklanır
    \item Veri ambarı yapılandırılmış ve önceden işlenmiş verileri saklarken, veri gölü ham ve yapısal olmayan verileri saklar
    \item Veri gölü sadece Excel dosyalarını saklar
    \item Aralarında hiçbir fark yoktur
\end{enumerate}
\vspace{4pt}
\begin{tcolorbox}[title=Cevap ve Açıklama, colback=gray!3!white, colframe=primarycolor!90!black, arc=2mm, fonttitle=\bfseries\sffamily, top=6pt, bottom=6pt]
    \textbf{Doğru Cevap: B}
    \vspace{2pt}
    Data Warehouse yapılandırılmış, temizlenmiş ve şeması belirlenmiş verileri raporlama için saklar; Data Lake ise her türlü ham veriyi saklar.
\end{tcolorbox}
\vspace{10pt}

\needspace{5cm}
\vspace{8pt}
\noindent\textbf{\large\sffamily Soru 122.} \ Büyük veri analitiğinde, verinin kalitesini artırmak ve analizlerin yanlış sonuç vermesini önlemek için yapılan ilk işlem hangisidir?
\begin{enumerate}[label=\Alph*), leftmargin=25pt, topsep=4pt, itemsep=2pt]
    \item Veri Görselleştirme
    \item Veri Temizleme ve Ön İşleme (Data Cleaning)
    \item Yapay Zekâ Eğitimi
    \item Raporu sunucuya kaydetme
\end{enumerate}
\vspace{4pt}
\begin{tcolorbox}[title=Cevap ve Açıklama, colback=gray!3!white, colframe=primarycolor!90!black, arc=2mm, fonttitle=\bfseries\sffamily, top=6pt, bottom=6pt]
    \textbf{Doğru Cevap: B}
    \vspace{2pt}
    Büyük veri gürültülü ve eksik olabileceğinden, analizden önce boş değerlerin doldurulması ve hatalı kayıtların silinmesi (veri temizleme) gerekir.
\end{tcolorbox}
\vspace{10pt}

\needspace{5cm}
\vspace{8pt}
\noindent\textbf{\large\sffamily Soru 123.} \ Hangi kavram, işletmelerin karar alma süreçlerini tamamen verilere ve algoritmalara dayandırarak objektif hale getirmesini tanımlar?
\begin{enumerate}[label=\Alph*), leftmargin=25pt, topsep=4pt, itemsep=2pt]
    \item Sezgisel Yönetim
    \item Veri Odaklı Karar Alma (Data-Driven Decision Making)
    \item Geleneksel Bürokrasi
    \item Ortak Malların Yönetimi
\end{enumerate}
\vspace{4pt}
\begin{tcolorbox}[title=Cevap ve Açıklama, colback=gray!3!white, colframe=primarycolor!90!black, arc=2mm, fonttitle=\bfseries\sffamily, top=6pt, bottom=6pt]
    \textbf{Doğru Cevap: B}
    \vspace{2pt}
    Veri odaklı karar alma (DDDM), kararların varsayımlardan ziyade analiz edilmiş verilere dayandırılması sürecidir.
\end{tcolorbox}
\vspace{10pt}

\needspace{5cm}
\vspace{8pt}
\noindent\textbf{\large\sffamily Soru 124.} \ Büyük Verinin 3V'sinden 'Variety' (Çeşitlilik) işletmeler için hangi zorluğu beraberinde getirir?
\begin{enumerate}[label=\Alph*), leftmargin=25pt, topsep=4pt, itemsep=2pt]
    \item Sabit disklerin ucuzlamasını
    \item Farklı formatlardaki (video, ses, metin) verilerin birleştirilmesi ve analiz edilmesinin zor olmasını
    \item Verilerin çok hızlı silinmesini
    \item İnternet hızının düşmesini
\end{enumerate}
\vspace{4pt}
\begin{tcolorbox}[title=Cevap ve Açıklama, colback=gray!3!white, colframe=primarycolor!90!black, arc=2mm, fonttitle=\bfseries\sffamily, top=6pt, bottom=6pt]
    \textbf{Doğru Cevap: B}
    \vspace{2pt}
    Çeşitlilik nedeniyle, yapısal olmayan verileri (ses, görüntü) klasik ilişkisel veritabanlarında analiz etmek zordur; özel büyük veri teknolojileri (Hadoop, Spark vb.) gerektirir.
\end{tcolorbox}
\vspace{10pt}

\needspace{5cm}
\vspace{8pt}
\noindent\textbf{\large\sffamily Soru 125.} \ Aşağıdakilerden hangisi bir işletmenin 'veri stratejisi' geliştirmesinin temel amaçlarından biridir?
\begin{enumerate}[label=\Alph*), leftmargin=25pt, topsep=4pt, itemsep=2pt]
    \item Verileri rakiplerle ücretsiz paylaşmak
    \item Veriyi iş süreçlerini optimize etmek ve yeni iş modelleri/değer yaratmak için stratejik bir varlık haline getirmek
    \item Bütün fiziksel sunucuları imha etmek
    \item Çalışanların kişisel bilgilerini internete yüklemek
\end{enumerate}
\vspace{4pt}
\begin{tcolorbox}[title=Cevap ve Açıklama, colback=gray!3!white, colframe=primarycolor!90!black, arc=2mm, fonttitle=\bfseries\sffamily, top=6pt, bottom=6pt]
    \textbf{Doğru Cevap: B}
    \vspace{2pt}
    Veri stratejisi, işletmenin veriyi depolama, yönetme, analiz etme ve ticari değere/faydaya dönüştürme yol haritasıdır.
\end{tcolorbox}
\vspace{10pt}

\needspace{5cm}
\vspace{8pt}
\noindent\textbf{\large\sffamily Soru 126.} \ Alibaba bünyesindeki Ant Financial firmasının, geleneksel bankalara kıyasla çok az çalışanla 1 milyar müşteriye hizmet vermesini sağlayan yapı hangisidir?
\begin{enumerate}[label=\Alph*), leftmargin=25pt, topsep=4pt, itemsep=2pt]
    \item Çok büyük çağrı merkezleri
    \item Yapay Zekâ Karar Fabrikası (AI Decision Factory)
    \item Ucuz iş gücü kullanımı
    \item Şube sayısının çok fazla olması
\end{enumerate}
\vspace{4pt}
\begin{tcolorbox}[title=Cevap ve Açıklama, colback=gray!3!white, colframe=primarycolor!90!black, arc=2mm, fonttitle=\bfseries\sffamily, top=6pt, bottom=6pt]
    \textbf{Doğru Cevap: B}
    \vspace{2pt}
    Ant Financial, insan müdahalesi olmadan kredi onayı ve operasyonları yürüten algoritmik bir yapay zekâ fabrikası kurmuştur.
\end{tcolorbox}
\vspace{10pt}

\needspace{5cm}
\vspace{8pt}
\noindent\textbf{\large\sffamily Soru 127.} \ Yapay Zekâ Fabrikası'nın (AI Factory) dört temel bileşeninden hangisi yeni algoritmaların performansını gerçek kullanıcılar üzerinde test eder?
\begin{enumerate}[label=\Alph*), leftmargin=25pt, topsep=4pt, itemsep=2pt]
    \item Veri Hattı (Data Pipeline)
    \item Kestirimci Karar Motoru (Predictive Engine)
    \item Deney Platformu (Experimentation Platform)
    \item Yazılım Entegrasyonu
\end{enumerate}
\vspace{4pt}
\begin{tcolorbox}[title=Cevap ve Açıklama, colback=gray!3!white, colframe=primarycolor!90!black, arc=2mm, fonttitle=\bfseries\sffamily, top=6pt, bottom=6pt]
    \textbf{Doğru Cevap: C}
    \vspace{2pt}
    Deney Platformu (A/B testing platformu), yeni algoritmaların veya işlevlerin gerçek dünyada test edilmesini sağlar.
\end{tcolorbox}
\vspace{10pt}

\needspace{5cm}
\vspace{8pt}
\noindent\textbf{\large\sffamily Soru 128.} \ Sadece tek bir dar alanda uzmanlaşmış (örneğin sadece kanserli hücre tespiti yapan veya satranç oynayan) yapay zekâ türü hangisidir?
\begin{enumerate}[label=\Alph*), leftmargin=25pt, topsep=4pt, itemsep=2pt]
    \item Güçlü Yapay Zeka (Strong AI)
    \item Zayıf Yapay Zeka (Narrow / Weak AI)
    \item Genel Yapay Zeka (AGI)
    \item Süper Yapay Zeka
\end{enumerate}
\vspace{4pt}
\begin{tcolorbox}[title=Cevap ve Açıklama, colback=gray!3!white, colframe=primarycolor!90!black, arc=2mm, fonttitle=\bfseries\sffamily, top=6pt, bottom=6pt]
    \textbf{Doğru Cevap: B}
    \vspace{2pt}
    Dar veya Zayıf Yapay Zeka (Narrow AI), sadece belirli bir sınırlandırılmış görevde insan seviyesinde veya üzerinde performans gösteren sistemlerdir.
\end{tcolorbox}
\vspace{10pt}

\needspace{5cm}
\vspace{8pt}
\noindent\textbf{\large\sffamily Soru 129.} \ Yapay Zekâ Fabrikasının girdisini oluşturan ve verilerin otomatik toplanmasını sağlayan bileşen hangisidir?
\begin{enumerate}[label=\Alph*), leftmargin=25pt, topsep=4pt, itemsep=2pt]
    \item Algoritmik Karar Motoru
    \item Veri Hattı (Data Pipeline)
    \item A/B Test Sistemi
    \item Mobil Arayüz
\end{enumerate}
\vspace{4pt}
\begin{tcolorbox}[title=Cevap ve Açıklama, colback=gray!3!white, colframe=primarycolor!90!black, arc=2mm, fonttitle=\bfseries\sffamily, top=6pt, bottom=6pt]
    \textbf{Doğru Cevap: B}
    \vspace{2pt}
    Veri Hattı (Data Pipeline), büyük verileri toplayıp, temizleyip algoritmaların kullanacağı hale getiren otomatik altyapıdır.
\end{tcolorbox}
\vspace{10pt}

\needspace{5cm}
\vspace{8pt}
\noindent\textbf{\large\sffamily Soru 130.} \ Yapay Zekâ Fabrikasının kalbi sayılan ve girdilere göre tahminlerde bulunan modelleri barındıran bileşen hangisidir?
\begin{enumerate}[label=\Alph*), leftmargin=25pt, topsep=4pt, itemsep=2pt]
    \item Veri Hattı
    \item Kestirimci Karar Motoru (Predictive Engine)
    \item Deney Platformu
    \item Bulut Depolama
\end{enumerate}
\vspace{4pt}
\begin{tcolorbox}[title=Cevap ve Açıklama, colback=gray!3!white, colframe=primarycolor!90!black, arc=2mm, fonttitle=\bfseries\sffamily, top=6pt, bottom=6pt]
    \textbf{Doğru Cevap: B}
    \vspace{2pt}
    Kestirimci Karar Motoru (Predictive Engine), makine öğrenmesi modelleriyle girdi verilerinden kararlar ve tahminler üretir.
\end{tcolorbox}
\vspace{10pt}

\needspace{5cm}
\vspace{8pt}
\noindent\textbf{\large\sffamily Soru 131.} \ Avrupa Birliği genelinde yapay zekânın etik ve güvenli kullanımını sağlamak amacıyla yürürlüğe giren yasal regülasyon hangisidir?
\begin{enumerate}[label=\Alph*), leftmargin=25pt, topsep=4pt, itemsep=2pt]
    \item GDPR
    \item EU AI Act (AB Yapay Zeka Yasası)
    \item Kyoto Protokolü
    \item Brundtland Kanunu
\end{enumerate}
\vspace{4pt}
\begin{tcolorbox}[title=Cevap ve Açıklama, colback=gray!3!white, colframe=primarycolor!90!black, arc=2mm, fonttitle=\bfseries\sffamily, top=6pt, bottom=6pt]
    \textbf{Doğru Cevap: B}
    \vspace{2pt}
    EU AI Act, yapay zeka teknolojilerini risk seviyelerine göre sınıflandırıp denetleyen ilk kapsamlı uluslararası yapay zeka kanunudur.
\end{tcolorbox}
\vspace{10pt}

\needspace{5cm}
\vspace{8pt}
\noindent\textbf{\large\sffamily Soru 132.} \ İnsanın sahip olduğu genel bilişsel yeteneklere sahip, her türlü değişen koşula adapte olabilen yapay zekâ seviyesi hangisidir?
\begin{enumerate}[label=\Alph*), leftmargin=25pt, topsep=4pt, itemsep=2pt]
    \item Zayıf Yapay Zeka
    \item Güçlü / Genel Yapay Zeka (AGI / Strong AI)
    \item Karar Motoru
    \item Uzman Sistemler
\end{enumerate}
\vspace{4pt}
\begin{tcolorbox}[title=Cevap ve Açıklama, colback=gray!3!white, colframe=primarycolor!90!black, arc=2mm, fonttitle=\bfseries\sffamily, top=6pt, bottom=6pt]
    \textbf{Doğru Cevap: B}
    \vspace{2pt}
    Güçlü veya Genel Yapay Zeka (AGI - Artificial General Intelligence), insanın yapabildiği her türlü zihinsel görevi yerine getirebilen teorik yapay zekadır.
\end{tcolorbox}
\vspace{10pt}

\needspace{5cm}
\vspace{8pt}
\noindent\textbf{\large\sffamily Soru 133.} \ Ant Financial firmasının 2018 sonlarında ulaştığı 150 milyar dolarlık piyasa değeri, ABD'nin hangi dev bankasının yaklaşık yarısına denk gelmektedir?
\begin{enumerate}[label=\Alph*), leftmargin=25pt, topsep=4pt, itemsep=2pt]
    \item Bank of America
    \item JPMorgan Chase
    \item Wells Fargo
    \item Goldman Sachs
\end{enumerate}
\vspace{4pt}
\begin{tcolorbox}[title=Cevap ve Açıklama, colback=gray!3!white, colframe=primarycolor!90!black, arc=2mm, fonttitle=\bfseries\sffamily, top=6pt, bottom=6pt]
    \textbf{Doğru Cevap: B}
    \vspace{2pt}
    Ant Financial, çok az fiziki varlığı ve çalışanı olmasına rağmen yapay zeka gücüyle JPMorgan Chase'in (ABD'nin en büyük bankası) yarısı kadar piyasa değerine ulaşmıştır.
\end{tcolorbox}
\vspace{10pt}

\needspace{5cm}
\vspace{8pt}
\noindent\textbf{\large\sffamily Soru 134.} \ Yapay zekâ algoritmalarının eğitilmesi ve karar alma yeteneklerinin geliştirilmesi için en kritik girdi aşağıdakilerden hangisidir?
\begin{enumerate}[label=\Alph*), leftmargin=25pt, topsep=4pt, itemsep=2pt]
    \item Ofis binalarının konforu
    \item Büyük ve kaliteli veri kümeleri (Data)
    \item Sunucuların dış kasa tasarımı
    \item Şirket hisselerinin fiyatı
\end{enumerate}
\vspace{4pt}
\begin{tcolorbox}[title=Cevap ve Açıklama, colback=gray!3!white, colframe=primarycolor!90!black, arc=2mm, fonttitle=\bfseries\sffamily, top=6pt, bottom=6pt]
    \textbf{Doğru Cevap: B}
    \vspace{2pt}
    Yapay zeka modellerinin (özellikle derin öğrenme) başarısı, eğitildikleri veri kümesinin büyüklüğüne ve kalitesine doğrudan bağlıdır.
\end{tcolorbox}
\vspace{10pt}

\needspace{5cm}
\vspace{8pt}
\noindent\textbf{\large\sffamily Soru 135.} \ Yapay Zekâ Fabrikası'nın son bileşeni olan ve üretilen tahminleri/kararları son kullanıcı uygulamalarına entegre eden yapı hangisidir?
\begin{enumerate}[label=\Alph*), leftmargin=25pt, topsep=4pt, itemsep=2pt]
    \item Veri Hattı
    \item Deney Platformu
    \item Yazılım Entegrasyonu ve Altyapı
    \item Veri Gölü
\end{enumerate}
\vspace{4pt}
\begin{tcolorbox}[title=Cevap ve Açıklama, colback=gray!3!white, colframe=primarycolor!90!black, arc=2mm, fonttitle=\bfseries\sffamily, top=6pt, bottom=6pt]
    \textbf{Doğru Cevap: C}
    \vspace{2pt}
    Yazılım Entegrasyonu ve Altyapı, algoritmaların çıktılarını web/mobil uygulamalar gibi arayüzlere gömerek müşteriye ulaştırır.
\end{tcolorbox}
\vspace{10pt}

\needspace{5cm}
\vspace{8pt}
\noindent\textbf{\large\sffamily Soru 136.} \ Aşağıdakilerden hangisi zayıf yapay zekâya (Narrow AI) bir örnek değildir?
\begin{enumerate}[label=\Alph*), leftmargin=25pt, topsep=4pt, itemsep=2pt]
    \item E-ticaret sitelerindeki ürün tavsiye algoritmaları
    \item İnsan gibi duygusal bilince sahip olan ve her konuda muhakeme yapabilen bir android robot
    \item Spam e-postaları filtreleyen sistemler
    \item Fotoğraftaki nesneleri etiketleyen görüntü işleme yazılımları
\end{enumerate}
\vspace{4pt}
\begin{tcolorbox}[title=Cevap ve Açıklama, colback=gray!3!white, colframe=primarycolor!90!black, arc=2mm, fonttitle=\bfseries\sffamily, top=6pt, bottom=6pt]
    \textbf{Doğru Cevap: B}
    \vspace{2pt}
    Duygusal bilince sahip, genel muhakeme yeteneği olan robot 'Güçlü/Genel Yapay Zeka' (AGI) sınıfına girer; günümüzde henüz örneği yoktur.
\end{tcolorbox}
\vspace{10pt}

\needspace{5cm}
\vspace{8pt}
\noindent\textbf{\large\sffamily Soru 137.} \ Yapay zekânın kurumlara sağladığı operasyonel ölçeklenebilirlik avantajı neyi ifade eder?
\begin{enumerate}[label=\Alph*), leftmargin=25pt, topsep=4pt, itemsep=2pt]
    \item Şirketin fiziksel büyüklüğünü artırmasını
    \item Müşteri sayısı artsa bile operasyonel maliyetlerin doğrusal artmamasını, algoritmaların aynı anda milyonlarca işleme bakabilmesini
    \item Şirketin hisse senetlerini satmasını
    \item Tüm işleri yöneticilerin yapmasını
\end{enumerate}
\vspace{4pt}
\begin{tcolorbox}[title=Cevap ve Açıklama, colback=gray!3!white, colframe=primarycolor!90!black, arc=2mm, fonttitle=\bfseries\sffamily, top=6pt, bottom=6pt]
    \textbf{Doğru Cevap: B}
    \vspace{2pt}
    Operasyonel ölçeklenebilirlik, insan gücünün aksine algoritmaların marjinal maliyeti neredeyse sıfıra yakın olacak şekilde milyonlarca kullanıcıya aynı anda hizmet verebilmesidir.
\end{tcolorbox}
\vspace{10pt}

\needspace{5cm}
\vspace{8pt}
\noindent\textbf{\large\sffamily Soru 138.} \ AB Yapay Zeka Yasası (EU AI Act) yapay zeka sistemlerini risk derecelerine göre sınıflandırırken, 'sosyal puanlama' yapan sistemleri hangi risk grubuna dahil etmiştir?
\begin{enumerate}[label=\Alph*), leftmargin=25pt, topsep=4pt, itemsep=2pt]
    \item Düşük Risk
    \item Kabul Edilemez Risk (Yasaklı sistemler)
    \item Yüksek Risk
    \item Sınırlı Risk
\end{enumerate}
\vspace{4pt}
\begin{tcolorbox}[title=Cevap ve Açıklama, colback=gray!3!white, colframe=primarycolor!90!black, arc=2mm, fonttitle=\bfseries\sffamily, top=6pt, bottom=6pt]
    \textbf{Doğru Cevap: B}
    \vspace{2pt}
    Vatandaşların davranışlarına göre puanlandığı 'sosyal puanlama' sistemleri ve biyometrik gözetim sistemleri AB Yapay Zeka Yasası'na göre 'Kabul Edilemez Risk' sınıfında olup yasaklanmıştır.
\end{tcolorbox}
\vspace{10pt}

\needspace{5cm}
\vspace{8pt}
\noindent\textbf{\large\sffamily Soru 139.} \ Yapay zeka modellerinde 'aşırı öğrenme' (overfitting) ne anlama gelmektedir?
\begin{enumerate}[label=\Alph*), leftmargin=25pt, topsep=4pt, itemsep=2pt]
    \item Modelin veriyi çok yavaş öğrenmesi
    \item Modelin eğitim verilerini ezberlemesi ancak yeni/farklı verilerde başarısız olması
    \item Modelin tüm dilleri öğrenmesi
    \item Modelin çalışmayı durdurması
\end{enumerate}
\vspace{4pt}
\begin{tcolorbox}[title=Cevap ve Açıklama, colback=gray!3!white, colframe=primarycolor!90!black, arc=2mm, fonttitle=\bfseries\sffamily, top=6pt, bottom=6pt]
    \textbf{Doğru Cevap: B}
    \vspace{2pt}
    Overfitting, modelin eğitim veri kümesindeki gürültüleri ezberlemesi ve gerçek hayattaki genel verilere uygulandığında tahmin gücünün düşmesidir.
\end{tcolorbox}
\vspace{10pt}

\needspace{5cm}
\vspace{8pt}
\noindent\textbf{\large\sffamily Soru 140.} \ Yapay zekânın 'karar destek sistemi' olarak kullanılması ne demektir?
\begin{enumerate}[label=\Alph*), leftmargin=25pt, topsep=4pt, itemsep=2pt]
    \item Kararları tamamen yöneticilerin hislerine bırakması
    \item Kararın yapay zeka analiz ve önerileri doğrultusunda insan yönetici tarafından verilmesi
    \item Kararı sadece yazılımcıların alması
    \item Karar alma sürecinin tamamen kaldırılması
\end{enumerate}
\vspace{4pt}
\begin{tcolorbox}[title=Cevap ve Açıklama, colback=gray!3!white, colframe=primarycolor!90!black, arc=2mm, fonttitle=\bfseries\sffamily, top=6pt, bottom=6pt]
    \textbf{Doğru Cevap: B}
    \vspace{2pt}
    Karar destek sistemlerinde yapay zeka büyük veriyi işleyip seçenekler sunar ve kararı optimize eder; son onayı insan yönetici verir.
\end{tcolorbox}
\vspace{10pt}

\needspace{5cm}
\vspace{8pt}
\noindent\textbf{\large\sffamily Soru 141.} \ Sürücüsüz otonom araçların kaza anında vereceği kararların ahlaki boyutu yapay zekânın hangi alanıyla ilgilidir?
\begin{enumerate}[label=\Alph*), leftmargin=25pt, topsep=4pt, itemsep=2pt]
    \item Donanım Mühendisliği
    \item Yapay Zekâ Etiği
    \item Veritabanı Yönetimi
    \item Ağ Altyapısı
\end{enumerate}
\vspace{4pt}
\begin{tcolorbox}[title=Cevap ve Açıklama, colback=gray!3!white, colframe=primarycolor!90!black, arc=2mm, fonttitle=\bfseries\sffamily, top=6pt, bottom=6pt]
    \textbf{Doğru Cevap: B}
    \vspace{2pt}
    Yapay zeka etiği (AI Ethics), algoritmik kararların adil, güvenli, şeffaf olmasını ve ahlaki değerlerle çelişmemesini inceleyen disiplindir.
\end{tcolorbox}
\vspace{10pt}

\needspace{5cm}
\vspace{8pt}
\noindent\textbf{\large\sffamily Soru 142.} \ Yapay Zeka Fabrikası mimarisinde 'A/B Testi' ne amaçla kullanılır?
\begin{enumerate}[label=\Alph*), leftmargin=25pt, topsep=4pt, itemsep=2pt]
    \item Verilerin yedeklenmesi için
    \item İki farklı algoritmanın/tasarımın kullanıcı grubu üzerinde test edilip hangisinin daha başarılı olduğunu belirlemek için
    \item Çalışanların performansını ölçmek için
    \item Finansal bütçeyi denetlemek için
\end{enumerate}
\vspace{4pt}
\begin{tcolorbox}[title=Cevap ve Açıklama, colback=gray!3!white, colframe=primarycolor!90!black, arc=2mm, fonttitle=\bfseries\sffamily, top=6pt, bottom=6pt]
    \textbf{Doğru Cevap: B}
    \vspace{2pt}
    A/B testi, sistemdeki bir değişikliğin veya yeni bir modelin (A sürümü vs B sürümü) kullanıcılar üzerindeki etkisini bilimsel olarak ölçmeye yarar.
\end{tcolorbox}
\vspace{10pt}

\needspace{5cm}
\vspace{8pt}
\noindent\textbf{\large\sffamily Soru 143.} \ Büyük dil modellerinin (LLM) eğitilmesinde kullanılan ve insan beynindeki sinapsları taklit eden yapılara ne ad verilir?
\begin{enumerate}[label=\Alph*), leftmargin=25pt, topsep=4pt, itemsep=2pt]
    \item Optik Kablolar
    \item Yapay Sinir Ağları (Artificial Neural Networks)
    \item Karar Ağaçları
    \item Veri Tabanı İndeksleri
\end{enumerate}
\vspace{4pt}
\begin{tcolorbox}[title=Cevap ve Açıklama, colback=gray!3!white, colframe=primarycolor!90!black, arc=2mm, fonttitle=\bfseries\sffamily, top=6pt, bottom=6pt]
    \textbf{Doğru Cevap: B}
    \vspace{2pt}
    Yapay Sinir Ağları (ANN), katmanlar halinde birbirine bağlı yapay nöronlardan oluşan ve karmaşık örüntüleri öğrenen algoritmik yapılardır.
\end{tcolorbox}
\vspace{10pt}

\needspace{5cm}
\vspace{8pt}
\noindent\textbf{\large\sffamily Soru 144.} \ Yapay zekâ projelerinde 'kara kutu' (black box) problemi neyi ifade eder?
\begin{enumerate}[label=\Alph*), leftmargin=25pt, topsep=4pt, itemsep=2pt]
    \item Sunucuların siyah renkte olmasını
    \item Özellikle derin öğrenme modellerinin kararları nasıl aldığının (hangi mantığa dayandırdığının) insanlar tarafından şeffaf olarak görülememesi
    \item Verilerin şifrelenip silinmesini
    \item Uçaklardaki kara kutunun bozulmasını
\end{enumerate}
\vspace{4pt}
\begin{tcolorbox}[title=Cevap ve Açıklama, colback=gray!3!white, colframe=primarycolor!90!black, arc=2mm, fonttitle=\bfseries\sffamily, top=6pt, bottom=6pt]
    \textbf{Doğru Cevap: B}
    \vspace{2pt}
    Kara kutu problemi, karmaşık modellerin (derin öğrenme) karar mekanizmasının açıklanabilir olmaması (şeffaflık eksikliği) sorunudur.
\end{tcolorbox}
\vspace{10pt}

\needspace{5cm}
\vspace{8pt}
\noindent\textbf{\large\sffamily Soru 145.} \ Görsel veya metinsel girdilerden yola çıkarak sıfırdan yeni içerik (resim, kod, makale) üretebilen yapay zeka türü hangisidir?
\begin{enumerate}[label=\Alph*), leftmargin=25pt, topsep=4pt, itemsep=2pt]
    \item Tanımlayıcı Yapay Zeka
    \item Üretken Yapay Zekâ (Generative AI)
    \item Kural Tabanlı Sistemler
    \item Otonom Lojistik Algoritması
\end{enumerate}
\vspace{4pt}
\begin{tcolorbox}[title=Cevap ve Açıklama, colback=gray!3!white, colframe=primarycolor!90!black, arc=2mm, fonttitle=\bfseries\sffamily, top=6pt, bottom=6pt]
    \textbf{Doğru Cevap: B}
    \vspace{2pt}
    Üretken Yapay Zekâ (Generative AI - ChatGPT, Midjourney vb.), eğitildiği verilerden yeni özgün çıktılar üreten yapay zeka alanıdır.
\end{tcolorbox}
\vspace{10pt}

\needspace{5cm}
\vspace{8pt}
\noindent\textbf{\large\sffamily Soru 146.} \ Aeroflot havayolu şirketinin dijital dönüşüm başarısı, dijitalleşmeyle ilgili hangi efsanenin yanlış olduğunu kanıtlar?
\begin{enumerate}[label=\Alph*), leftmargin=25pt, topsep=4pt, itemsep=2pt]
    \item Startup satın almanın zorunlu olduğunu
    \item Dijital dönüşümün sadece devasa yatırımlar ve iş modelinde radikal yıkımlar gerektirdiği efsanesini
    \item Fiziksel kanalların tamamen kalkacağını
    \item Altyapıyı toptan yenilemek gerektiğini
\end{enumerate}
\vspace{4pt}
\begin{tcolorbox}[title=Cevap ve Açıklama, colback=gray!3!white, colframe=primarycolor!90!black, arc=2mm, fonttitle=\bfseries\sffamily, top=6pt, bottom=6pt]
    \textbf{Doğru Cevap: B}
    \vspace{2pt}
    Aeroflot, iş modelini yıkmadan süreçlerini dijitalleştirip raporlamayı düzelterek ve operasyonel verim kazanarak NPS skorunu \%44 artırmış, yıkım efsanesini çürütmüştür.
\end{tcolorbox}
\vspace{10pt}

\needspace{5cm}
\vspace{8pt}
\noindent\textbf{\large\sffamily Soru 147.} \ Galeries Lafayette Champs-Élysées mağazası, fiziksel akıllı mağaza tasarımıyla hangi dijital dönüşüm gerçeğini temsil eder?
\begin{enumerate}[label=\Alph*), leftmargin=25pt, topsep=4pt, itemsep=2pt]
    \item Fiziksel mağazaların tamamen öleceğini
    \item Dijital teknolojiler ile fiziksel deneyimin (stilistler aracılığıyla kurulan duygusal bağ) harmanlandığı hibrit modeli
    \item Müşterilerin robotlarla alışveriş yapacağını
    \item E-ticaretin yasaklanacağını
\end{enumerate}
\vspace{4pt}
\begin{tcolorbox}[title=Cevap ve Açıklama, colback=gray!3!white, colframe=primarycolor!90!black, arc=2mm, fonttitle=\bfseries\sffamily, top=6pt, bottom=6pt]
    \textbf{Doğru Cevap: B}
    \vspace{2pt}
    Akıllı lüks mağaza, dijital veri toplama ile stilistlerin fiziksel temasını harmanlayarak müşterilerle güçlü bir duygusal bağ kurmuştur.
\end{tcolorbox}
\vspace{10pt}

\needspace{5cm}
\vspace{8pt}
\noindent\textbf{\large\sffamily Soru 148.} \ Avnet firmasının Hackster.io startup'ını satın alıp onu kendi hiyerarşisinde boğmak yerine yarı bağımsız bırakması hangi efsaneyi çürütür?
\begin{enumerate}[label=\Alph*), leftmargin=25pt, topsep=4pt, itemsep=2pt]
    \item Şirketlerin startup satın almaması gerektiği efsanesini
    \item Satın alınan startup'ların ana firmanın bürokratik yapısına tamamen entegre edilmesi gerektiği efsanesini
    \item Startup'ların vergi ödemediği efsanesini
    \item Teknolojinin gereksiz olduğu efsanesini
\end{enumerate}
\vspace{4pt}
\begin{tcolorbox}[title=Cevap ve Açıklama, colback=gray!3!white, colframe=primarycolor!90!black, arc=2mm, fonttitle=\bfseries\sffamily, top=6pt, bottom=6pt]
    \textbf{Doğru Cevap: B}
    \vspace{2pt}
    Büyük şirketlerin startup'ları entegrasyon adı altında bürokraside boğduğu bilinir. Avnet, Hackster.io'yu bağımsız bırakarak sinerji yaratılmasını sağlamıştır.
\end{tcolorbox}
\vspace{10pt}

\needspace{5cm}
\vspace{8pt}
\noindent\textbf{\large\sffamily Soru 149.} \ Adobe firmasının çalışanlarına fikir üretme fırsatı tanıyan 'Turuncu Kutu' (Orange Box) projesi, dönüşümün hangi boyutunun önemini vurgular?
\begin{enumerate}[label=\Alph*), leftmargin=25pt, topsep=4pt, itemsep=2pt]
    \item Sadece sunucu donanımının
    \item İnsan kaynağı, kurum kültürü ve çalışan katılımının
    \item Ofis malzemelerinin kalitesinin
    \item Reklam bütçelerinin büyüklüğünün
\end{enumerate}
\vspace{4pt}
\begin{tcolorbox}[title=Cevap ve Açıklama, colback=gray!3!white, colframe=primarycolor!90!black, arc=2mm, fonttitle=\bfseries\sffamily, top=6pt, bottom=6pt]
    \textbf{Doğru Cevap: B}
    \vspace{2pt}
    Adobe Turuncu Kutu projesi, dijital dönüşümün sadece teknoloji değil, insan ve kültür değişimiyle gerçekleşeceğini gösteren bir vaka örneğidir.
\end{tcolorbox}
\vspace{10pt}

\needspace{5cm}
\vspace{8pt}
\noindent\textbf{\large\sffamily Soru 150.} \ Marriott otel zincirinin Expedia rekabetinde altyapıyı toptan değiştirmek yerine aşamalı ilerlemesi hangi dönüşüm efsanesine yanıttır?
\begin{enumerate}[label=\Alph*), leftmargin=25pt, topsep=4pt, itemsep=2pt]
    \item Otelciliğin tamamen biteceği efsanesine
    \item Eski bilgi sistemleri altyapısını tamamen sıfırlayarak işe başlamak gerektiği efsanesine (bunun yerine aşamalı test önerilir)
    \item Müşterilerin sadece internetten rezervasyon yapacağı efsanesine
    \item Fiyatların sabit kalması gerektiğine
\end{enumerate}
\vspace{4pt}
\begin{tcolorbox}[title=Cevap ve Açıklama, colback=gray!3!white, colframe=primarycolor!90!black, arc=2mm, fonttitle=\bfseries\sffamily, top=6pt, bottom=6pt]
    \textbf{Doğru Cevap: B}
    \vspace{2pt}
    Altyapıyı bir gecede değiştirmek (legacy system replacement) yüksek risklidir. Marriott, aşamalı geçiş yaparak riskleri azaltmıştır.
\end{tcolorbox}
\vspace{10pt}

\needspace{5cm}
\vspace{8pt}
\noindent\textbf{\large\sffamily Soru 151.} \ Maersk firmasının konteyner taşımacılığında IoT sensörler kullanarak tedarik zincirinde şeffaflık sağlaması hangi dönüşüm gerçeğini gösterir?
\begin{enumerate}[label=\Alph*), leftmargin=25pt, topsep=4pt, itemsep=2pt]
    \item Nakliye şirketlerinin tamamen yazılım şirketine dönüşmesi gerektiğini
    \item Mevcut çekirdek değer önermesinin (nakliye) dijital araçlarla daha şeffaf ve verimli sunulabileceğini
    \item Gemilerin artık kaptansız çalışacağını
    \item Limanların kapatılacağını
\end{enumerate}
\vspace{4pt}
\begin{tcolorbox}[title=Cevap ve Açıklama, colback=gray!3!white, colframe=primarycolor!90!black, arc=2mm, fonttitle=\bfseries\sffamily, top=6pt, bottom=6pt]
    \textbf{Doğru Cevap: B}
    \vspace{2pt}
    Maersk, ana işini (konteyner taşımacılığını) IoT sensör verileriyle şeffaflaştırarak müşteriye daha yüksek değer sunmuştur.
\end{tcolorbox}
\vspace{10pt}

\needspace{5cm}
\vspace{8pt}
\noindent\textbf{\large\sffamily Soru 152.} \ Yöneticilerin 'Dijital dönüşüm fiziksel kanalları tamamen ortadan kaldıracak' inancı neden bir efsanedir?
\begin{enumerate}[label=\Alph*), leftmargin=25pt, topsep=4pt, itemsep=2pt]
    \item Fiziksel kanallar daha ucuz olduğu için
    \item İnsanların güven ve duygusal bağ kurmak için fiziksel temas/deneyim aramaya devam etmesi ve hibrit modelin başarısı nedeniyle
    \item İnternetin yakında çökecek olması nedeniyle
    \item Devletlerin e-ticareti sınırlandırması yüzünden
\end{enumerate}
\vspace{4pt}
\begin{tcolorbox}[title=Cevap ve Açıklama, colback=gray!3!white, colframe=primarycolor!90!black, arc=2mm, fonttitle=\bfseries\sffamily, top=6pt, bottom=6pt]
    \textbf{Doğru Cevap: B}
    \vspace{2pt}
    Saf dijital kanalların müşteri bağlılığı yaratmada zorlandığı, bu yüzden fiziksel ve dijitalin birleştiği 'phygital' hibrit modellerin yükseldiği görülmektedir.
\end{tcolorbox}
\vspace{10pt}

\needspace{5cm}
\vspace{8pt}
\noindent\textbf{\large\sffamily Soru 153.} \ Paris merkezli G7 Taksi şirketinin Uber rekabetine karşı kendi çağırma uygulamasını geliştirmesi hangi stratejiye örnektir?
\begin{enumerate}[label=\Alph*), leftmargin=25pt, topsep=4pt, itemsep=2pt]
    \item Şirketi tamamen kapatıp kaçma
    \item Teknolojik tehdide karşı dijital araçları hızla kendi süreçlerine entegre ederek direnç oluşturma
    \item Uber'i mahkemeye vererek yasaklatma
    \item Fiyatları tamamen bedava yapma
\end{enumerate}
\vspace{4pt}
\begin{tcolorbox}[title=Cevap ve Açıklama, colback=gray!3!white, colframe=primarycolor!90!black, arc=2mm, fonttitle=\bfseries\sffamily, top=6pt, bottom=6pt]
    \textbf{Doğru Cevap: B}
    \vspace{2pt}
    G7 Taksi, dijital yıkım tehdidine karşı yenilikçi bir rezervasyon algoritması geliştirerek pazar payını korumayı başarmıştır.
\end{tcolorbox}
\vspace{10pt}

\needspace{5cm}
\vspace{8pt}
\noindent\textbf{\large\sffamily Soru 154.} \ Dijital dönüşümde 'eski sistemlerin' (legacy systems) tamamen yenilenmesi girişimleri neden sıklıkla başarısızlıkla sonuçlanır?
\begin{enumerate}[label=\Alph*), leftmargin=25pt, topsep=4pt, itemsep=2pt]
    \item Yeni yazılımların çok çirkin olması nedeniyle
    \item Büyük bütçeli, uzun süren, iş süreçlerini kesintiye uğratan yüksek riskler ve adaptasyon zorlukları nedeniyle
    \item Bilgisayar virüsleri yüzünden
    \item Yöneticilerin emekli olması nedeniyle
\end{enumerate}
\vspace{4pt}
\begin{tcolorbox}[title=Cevap ve Açıklama, colback=gray!3!white, colframe=primarycolor!90!black, arc=2mm, fonttitle=\bfseries\sffamily, top=6pt, bottom=6pt]
    \textbf{Doğru Cevap: B}
    \vspace{2pt}
    Legacy sistemleri bir kerede değiştirmek (big bang migration), veri kaybı ve iş durması gibi büyük riskler taşır. Aşamalı dönüşüm tercih edilmelidir.
\end{tcolorbox}
\vspace{10pt}

\needspace{5cm}
\vspace{8pt}
\noindent\textbf{\large\sffamily Soru 155.} \ Startup satın alan büyük ölçekli geleneksel şirketlerin yaptığı en yaygın yönetim hatası aşağıdakilerden hangisidir?
\begin{enumerate}[label=\Alph*), leftmargin=25pt, topsep=4pt, itemsep=2pt]
    \item Startup'a çok yüksek maaş ödemek
    \item Startup'ı kendi kurumsal bürokrasisi, ağır onay mekanizmaları ve hiyerarşisiyle boğmak
    \item Startup çalışanlarını eğitime göndermek
    \item Şirketin logosunu değiştirmek
\end{enumerate}
\vspace{4pt}
\begin{tcolorbox}[title=Cevap ve Açıklama, colback=gray!3!white, colframe=primarycolor!90!black, arc=2mm, fonttitle=\bfseries\sffamily, top=6pt, bottom=6pt]
    \textbf{Doğru Cevap: B}
    \vspace{2pt}
    Geleneksel kurumsal yapı, startup'ın çevik ve hızlı karar alma kültürünü yok ederse satın alma başarısız olur. Yarı bağımsız yapı korunmalıdır.
\end{tcolorbox}
\vspace{10pt}

\needspace{5cm}
\vspace{8pt}
\noindent\textbf{\large\sffamily Soru 156.} \ Yıkıcı inovasyon ve dijital dönüşüm süreçlerinde 'organizasyonel çeviklik' (agility) neyi ifade eder?
\begin{enumerate}[label=\Alph*), leftmargin=25pt, topsep=4pt, itemsep=2pt]
    \item Şirket çalışanlarının spor yapmasını
    \item Değişen pazar koşullarına, teknolojilere ve müşteri taleplerine şirketin hızla adapte olabilme yeteneğini
    \item Şirketin hisse değerinin hızla yükselmesini
    \item Sunucu sayısının hızlıca artırılmasını
\end{enumerate}
\vspace{4pt}
\begin{tcolorbox}[title=Cevap ve Açıklama, colback=gray!3!white, colframe=primarycolor!90!black, arc=2mm, fonttitle=\bfseries\sffamily, top=6pt, bottom=6pt]
    \textbf{Doğru Cevap: B}
    \vspace{2pt}
    Çeviklik, bürokratik yavaşlığı azaltarak pazardaki fırsat veya tehditlere hızlı cevap verebilme kabiliyetidir.
\end{tcolorbox}
\vspace{10pt}

\needspace{5cm}
\vspace{8pt}
\noindent\textbf{\large\sffamily Soru 157.} \ Maersk vaka analizinde, IoT sensör verilerinin müşterilerle paylaşılması şirkette hangi kavramın gelişimini sağlamıştır?
\begin{enumerate}[label=\Alph*), leftmargin=25pt, topsep=4pt, itemsep=2pt]
    \item Gizlilik ihlali
    \item Tedarik zincirinde şeffaflık ve güven (transparency)
    \item Müşteri şikayetlerinin artışı
    \item Fiyat artışları
\end{enumerate}
\vspace{4pt}
\begin{tcolorbox}[title=Cevap ve Açıklama, colback=gray!3!white, colframe=primarycolor!90!black, arc=2mm, fonttitle=\bfseries\sffamily, top=6pt, bottom=6pt]
    \textbf{Doğru Cevap: B}
    \vspace{2pt}
    Verilerin şeffaf paylaşımı, müşterilerin yüklerinin nerede olduğunu ve sıcaklığını anlık görmesini sağlayarak güven ilişkisini pekiştirmiştir.
\end{tcolorbox}
\vspace{10pt}

\needspace{5cm}
\vspace{8pt}
\noindent\textbf{\large\sffamily Soru 158.} \ Aeroflot havayolu firmasının Net Tavsiye Skorunu (NPS) artırmasında rol oynayan dijital veri kullanımı hangisidir?
\begin{enumerate}[label=\Alph*), leftmargin=25pt, topsep=4pt, itemsep=2pt]
    \item Müşterilerin sosyal medya şifrelerini toplamak
    \item Operasyonel aksaklıkların, uçuş gecikmelerinin ve müşteri geri bildirimlerinin anlık veri analitiğiyle izlenip iyileştirilmesi
    \item Uçak biletlerinin fiyatlarını gizlemek
    \item Sadece business class müşterilerine anket yapmak
\end{enumerate}
\vspace{4pt}
\begin{tcolorbox}[title=Cevap ve Açıklama, colback=gray!3!white, colframe=primarycolor!90!black, arc=2mm, fonttitle=\bfseries\sffamily, top=6pt, bottom=6pt]
    \textbf{Doğru Cevap: B}
    \vspace{2pt}
    Aeroflot, uçuş operasyon verilerini ve müşteri şikayetlerini dijital entegrasyonla takip edip hızlı aksiyon alarak müşteri memnuniyetini (NPS) artırmıştır.
\end{tcolorbox}
\vspace{10pt}

\needspace{5cm}
\vspace{8pt}
\noindent\textbf{\large\sffamily Soru 159.} \ Dijital dönüşüm stratejilerinde 'müşteri odaklılık' (customer-centricity) neyi önceler?
\begin{enumerate}[label=\Alph*), leftmargin=25pt, topsep=4pt, itemsep=2pt]
    \item Şirketin kârını her koşulda maksimize etmeyi
    \item Müşteri ihtiyaçlarını, beklentilerini ve deneyimini dijital süreçlerin merkezine koyarak değer üretmeyi
    \item Sadece yeni müşteriler bulmayı
    \item Ürün fiyatlarını artırmayı
\end{enumerate}
\vspace{4pt}
\begin{tcolorbox}[title=Cevap ve Açıklama, colback=gray!3!white, colframe=primarycolor!90!black, arc=2mm, fonttitle=\bfseries\sffamily, top=6pt, bottom=6pt]
    \textbf{Doğru Cevap: B}
    \vspace{2pt}
    Müşteri odaklılık, dijital dönüşüm teknolojilerini kurarken müşteri deneyimini kolaylaştırmaya ve değer sunmaya odaklanmaktır.
\end{tcolorbox}
\vspace{10pt}

\needspace{5cm}
\vspace{8pt}
\noindent\textbf{\large\sffamily Soru 160.} \ Adobe'un 'Turuncu Kutu' projesinde çalışanlara fikir geliştirirken bütçe ve süre vermesi hangi kültürel mesajı taşır?
\begin{enumerate}[label=\Alph*), leftmargin=25pt, topsep=4pt, itemsep=2pt]
    \item Şirketin parasının çok olduğunu
    \item Yenilikçi fikirlerin denenebilmesi için kontrollü başarısızlığa ve denemeye alan açılması gerektiğini (inovasyon kültürü)
    \item Çalışanların çok çalışması gerektiğini
    \item Şirketin yazılımlarını satamadığını
\end{enumerate}
\vspace{4pt}
\begin{tcolorbox}[title=Cevap ve Açıklama, colback=gray!3!white, colframe=primarycolor!90!black, arc=2mm, fonttitle=\bfseries\sffamily, top=6pt, bottom=6pt]
    \textbf{Doğru Cevap: B}
    \vspace{2pt}
    Turuncu Kutu, çalışanların risk almasını ve yeni dijital fikirleri test etmesini teşvik ederek inovasyonu tabana yayan bir kültür çalışmasıdır.
\end{tcolorbox}
\vspace{10pt}

\needspace{5cm}
\vspace{8pt}
\noindent\textbf{\large\sffamily Soru 161.} \ Garrett Hardin'in 'Ortak Malların Trajedisi' makalesinde kullandığı klasik merada aşırı otlatma örneğinde, bireyin rasyonel kararı topluluk için nasıl bir sonuç doğurur?
\begin{enumerate}[label=\Alph*), leftmargin=25pt, topsep=4pt, itemsep=2pt]
    \item Meraların daha çok yeşillenmesini sağlar
    \item Bireyin kısa vadeli çıkarı (+1 hayvan), uzun vadede kaynağın (meranın) yok olmasına ve herkesin kaybetmesine yol açar
    \item Devletin meraları satın almasıyla sonuçlanır
    \item Hayvan fiyatlarının düşmesini sağlar
\end{enumerate}
\vspace{4pt}
\begin{tcolorbox}[title=Cevap ve Açıklama, colback=gray!3!white, colframe=primarycolor!90!black, arc=2mm, fonttitle=\bfseries\sffamily, top=6pt, bottom=6pt]
    \textbf{Doğru Cevap: B}
    \vspace{2pt}
    Ortak malların trajedisinde, bireylerin kendi faydalarını sınırsızca maksimize etme çabası, paylaşılan sınırlı kaynağın yok olmasıyla felakete yol açar.
\end{tcolorbox}
\vspace{10pt}

\needspace{5cm}
\vspace{8pt}
\noindent\textbf{\large\sffamily Soru 162.} \ 3 Aralık 1984'te Hindistan'ın Bhopal kentinde binlerce insanın ölümüne yol açan Union Carbide böcek ilacı fabrikasından sızan zehirli gaz hangisidir?
\begin{enumerate}[label=\Alph*), leftmargin=25pt, topsep=4pt, itemsep=2pt]
    \item Karbondioksit
    \item Metil İzosiyanat (MIC)
    \item Sülfürdioksit
    \item Karbonmonoksit
\end{enumerate}
\vspace{4pt}
\begin{tcolorbox}[title=Cevap ve Açıklama, colback=gray!3!white, colframe=primarycolor!90!black, arc=2mm, fonttitle=\bfseries\sffamily, top=6pt, bottom=6pt]
    \textbf{Doğru Cevap: B}
    \vspace{2pt}
    Bhopal felaketine yol açan ve tarihin en ölümcül endüstriyel gaz sızıntısı kabul edilen madde Metil İzosiyanat (MIC) gazıdır.
\end{tcolorbox}
\vspace{10pt}

\needspace{5cm}
\vspace{8pt}
\noindent\textbf{\large\sffamily Soru 163.} \ 1996 yılında Nike firmasının Pakistan'daki tedarik zincirinde çocuk işçi çalıştırdığının ortaya çıkması, kurumsal yönetimde hangi kavramın doğuşunu hızlandırmıştır?
\begin{enumerate}[label=\Alph*), leftmargin=25pt, topsep=4pt, itemsep=2pt]
    \item Tedarik Zinciri Otomasyonu
    \item Kurumsal Sosyal Sorumluluk (KSS) ve sosyal denetim standartları
    \item İthalat Vergileri
    \item Reklam Bütçesi Planlaması
\end{enumerate}
\vspace{4pt}
\begin{tcolorbox}[title=Cevap ve Açıklama, colback=gray!3!white, colframe=primarycolor!90!black, arc=2mm, fonttitle=\bfseries\sffamily, top=6pt, bottom=6pt]
    \textbf{Doğru Cevap: B}
    \vspace{2pt}
    Nike skandalı, küresel markaların sadece kendi fabrikalarından değil, tüm tedarik zincirindeki sosyal ve etik koşullardan (KSS) sorumlu olduğunu göstermiştir.
\end{tcolorbox}
\vspace{10pt}

\needspace{5cm}
\vspace{8pt}
\noindent\textbf{\large\sffamily Soru 164.} \ Gelecek nesillerin kendi ihtiyaçlarını karşılama yeteneğinden ödün vermeden bugünün ihtiyaçlarını karşılamayı hedefleyen kavram hangisidir?
\begin{enumerate}[label=\Alph*), leftmargin=25pt, topsep=4pt, itemsep=2pt]
    \item Ekonomik Büyüme
    \item Sürdürülebilirlik (Sustainability)
    \item Yaratıcı Yıkım
    \item Dijital Olgunluk
\end{enumerate}
\vspace{4pt}
\begin{tcolorbox}[title=Cevap ve Açıklama, colback=gray!3!white, colframe=primarycolor!90!black, arc=2mm, fonttitle=\bfseries\sffamily, top=6pt, bottom=6pt]
    \textbf{Doğru Cevap: B}
    \vspace{2pt}
    Brundtland Raporu'nda yapılan bu tanım sürdürülebilirlik (sustainability) kavramının en yaygın kabul gören tanımıdır.
\end{tcolorbox}
\vspace{10pt}

\needspace{5cm}
\vspace{8pt}
\noindent\textbf{\large\sffamily Soru 165.} \ Sürdürülebilirliğin değerlendirildiği üç ana boyuttan biri olan 'Sosyal Sürdürülebilirlik' aşağıdakilerden hangisini kapsar?
\begin{enumerate}[label=\Alph*), leftmargin=25pt, topsep=4pt, itemsep=2pt]
    \item Şirketin kâr marjlarını
    \item İnsan hakları, çalışma koşulları, fırsat eşitliği ve toplumsal refahı
    \item Fabrikadaki robot sayısını
    \item Karbon salınım oranını
\end{enumerate}
\vspace{4pt}
\begin{tcolorbox}[title=Cevap ve Açıklama, colback=gray!3!white, colframe=primarycolor!90!black, arc=2mm, fonttitle=\bfseries\sffamily, top=6pt, bottom=6pt]
    \textbf{Doğru Cevap: B}
    \vspace{2pt}
    Sosyal sürdürülebilirlik; adalet, çalışan hakları, topluluk refahı ve insan sağlığı gibi sosyal adalet konularını ele alır.
\end{tcolorbox}
\vspace{10pt}

\needspace{5cm}
\vspace{8pt}
\noindent\textbf{\large\sffamily Soru 166.} \ Aşağıdakilerden hangisi 'Ortak Malların Trajedisi'ne günümüz küresel çevre sorunlarından bir örnek olarak gösterilebilir?
\begin{enumerate}[label=\Alph*), leftmargin=25pt, topsep=4pt, itemsep=2pt]
    \item Özel bir tarlada domates yetiştirilmesi
    \item Okyanuslardaki aşırı avlanma ve deniz kaynaklarının tükenmesi
    \item Şirketlerin ofis yazılımları satın alması
    \item Bireylerin evlerinde su tasarrufu yapması
\end{enumerate}
\vspace{4pt}
\begin{tcolorbox}[title=Cevap ve Açıklama, colback=gray!3!white, colframe=primarycolor!90!black, arc=2mm, fonttitle=\bfseries\sffamily, top=6pt, bottom=6pt]
    \textbf{Doğru Cevap: B}
    \vspace{2pt}
    Okyanuslar ve denizler kamuya açık, ücretsiz ortak kaynaklardır. Sınırsız ve denetimsiz balık avcılığı bu kaynağı yok ederek trajediyi doğrular.
\end{tcolorbox}
\vspace{10pt}

\needspace{5cm}
\vspace{8pt}
\noindent\textbf{\large\sffamily Soru 167.} \ Garrett Hardin'in 'Ortak Malların Trajedisi' teorisine göre ortak kaynakların yok olmasını engellemek için önerilen iki temel çözüm yolu hangisidir?
\begin{enumerate}[label=\Alph*), leftmargin=25pt, topsep=4pt, itemsep=2pt]
    \item Kaynağın tamamen kullanıma kapatılması veya bedava yapılması
    \item Özel mülkiyet haline getirme (özelleştirme) veya devlet eliyle düzenleme/denetleme (yasal kısıtlamalar)
    \item Hayvan sayısının sınırsız artırılması veya meraların yakılması
    \item Hiçbir çözüm yolu yoktur
\end{enumerate}
\vspace{4pt}
\begin{tcolorbox}[title=Cevap ve Açıklama, colback=gray!3!white, colframe=primarycolor!90!black, arc=2mm, fonttitle=\bfseries\sffamily, top=6pt, bottom=6pt]
    \textbf{Doğru Cevap: B}
    \vspace{2pt}
    Hardin, ortak malların korunması için ya mülkiyetin özelleştirilmesini ya da devlet denetimi ve hukuki sınırlandırmalar (vergiler, kotalar) getirilmesini önermiştir.
\end{tcolorbox}
\vspace{10pt}

\needspace{5cm}
\vspace{8pt}
\noindent\textbf{\large\sffamily Soru 168.} \ Kurumsal Sosyal Sorumluluk (KSS) kapsamında bir şirketin çevreye verdiği zararları en aza indirmesi hangi sürdürülebilirlik boyutuyla ilgilidir?
\begin{enumerate}[label=\Alph*), leftmargin=25pt, topsep=4pt, itemsep=2pt]
    \item Ekonomik boyut
    \item Çevresel boyut
    \item Sosyal boyut
    \item Kültürel boyut
\end{enumerate}
\vspace{4pt}
\begin{tcolorbox}[title=Cevap ve Açıklama, colback=gray!3!white, colframe=primarycolor!90!black, arc=2mm, fonttitle=\bfseries\sffamily, top=6pt, bottom=6pt]
    \textbf{Doğru Cevap: B}
    \vspace{2pt}
    Çevresel sürdürülebilirlik; karbon salınımının azaltılması, atık yönetimi ve doğal kaynakların korunmasını hedefler.
\end{tcolorbox}
\vspace{10pt}

\needspace{5cm}
\vspace{8pt}
\noindent\textbf{\large\sffamily Soru 169.} \ Bhopal felaketinin ardından küresel ölçekte kimya endüstrisinde güvenlik standartlarını artırmak amacıyla başlatılan gönüllü programın adı nedir?
\begin{enumerate}[label=\Alph*), leftmargin=25pt, topsep=4pt, itemsep=2pt]
    \item Yeşil Barış (Greenpeace)
    \item Üçlü Sorumluluk (Responsible Care)
    \item Kyoto Protokolü
    \item ISO 9001
\end{enumerate}
\vspace{4pt}
\begin{tcolorbox}[title=Cevap ve Açıklama, colback=gray!3!white, colframe=primarycolor!90!black, arc=2mm, fonttitle=\bfseries\sffamily, top=6pt, bottom=6pt]
    \textbf{Doğru Cevap: B}
    \vspace{2pt}
    Responsible Care (Üçlü Sorumluluk), kimya sanayinin çevre, iş sağlığı ve güvenliği performansını sürekli iyileştirmeyi taahhüt ettiği küresel bir girişimdir.
\end{tcolorbox}
\vspace{10pt}

\needspace{5cm}
\vspace{8pt}
\noindent\textbf{\large\sffamily Soru 170.} \ Nike firmasının 1990'lardaki çocuk işçi krizinden sonra tedarikçilerini denetlemek için kurduğu bağımsız yapı hangisidir?
\begin{enumerate}[label=\Alph*), leftmargin=25pt, topsep=4pt, itemsep=2pt]
    \item FLA (Fair Labor Association / Adil Emeği Koruma Derneği)
    \item FIFA Denetim Kurulu
    \item Birleşmiş Milletler
    \item Dünya Ticaret Örgütü
\end{enumerate}
\vspace{4pt}
\begin{tcolorbox}[title=Cevap ve Açıklama, colback=gray!3!white, colframe=primarycolor!90!black, arc=2mm, fonttitle=\bfseries\sffamily, top=6pt, bottom=6pt]
    \textbf{Doğru Cevap: A}
    \vspace{2pt}
    Nike ve diğer giyim markaları, çalışma standartlarını denetlemek ve işçi haklarını korumak için FLA'nın kurulmasına öncülük etmiştir.
\end{tcolorbox}
\vspace{10pt}

\needspace{5cm}
\vspace{8pt}
\noindent\textbf{\large\sffamily Soru 171.} \ Geleneksel 'kâr odaklı' kurumsal yönetim anlayışının yerine, şirketin Çevresel, Sosyal ve Yönetişimsel (ESG) performansını da ölçen yaklaşım hangisidir?
\begin{enumerate}[label=\Alph*), leftmargin=25pt, topsep=4pt, itemsep=2pt]
    \item Tek Yönlü Finans
    \item Üçlü Bilanço (Triple Bottom Line - People, Planet, Profit)
    \item Schumpeterci Analiz
    \item IBM Olgunluk Analizi
\end{enumerate}
\vspace{4pt}
\begin{tcolorbox}[title=Cevap ve Açıklama, colback=gray!3!white, colframe=primarycolor!90!black, arc=2mm, fonttitle=\bfseries\sffamily, top=6pt, bottom=6pt]
    \textbf{Doğru Cevap: B}
    \vspace{2pt}
    Triple Bottom Line (Üçlü Bilanço), şirketlerin sadece finansal kârını (profit) değil, insan (people) ve çevre (planet) üzerindeki etkilerini de raporlamasını savunur.
\end{tcolorbox}
\vspace{10pt}

\needspace{5cm}
\vspace{8pt}
\noindent\textbf{\large\sffamily Soru 172.} \ Atmosferdeki karbondioksit miktarını artırarak küresel ısınmaya yol açan sera gazlarının en temel kaynağı aşağıdakilerden hangisidir?
\begin{enumerate}[label=\Alph*), leftmargin=25pt, topsep=4pt, itemsep=2pt]
    \item Rüzgar türbinleri
    \item Fosil yakıtların (kömür, petrol, doğalgaz) enerji ve üretimde aşırı tüketimi
    \item Tarımda damlama sulama yapılması
    \item Dijital kitap okunması
\end{enumerate}
\vspace{4pt}
\begin{tcolorbox}[title=Cevap ve Açıklama, colback=gray!3!white, colframe=primarycolor!90!black, arc=2mm, fonttitle=\bfseries\sffamily, top=6pt, bottom=6pt]
    \textbf{Doğru Cevap: B}
    \vspace{2pt}
    Sera gazı salınımlarının (karbon ayak izi) en büyük nedeni sanayide, elektrik üretiminde ve ulaşımda kullanılan fosil yakıtlardır.
\end{tcolorbox}
\vspace{10pt}

\needspace{5cm}
\vspace{8pt}
\noindent\textbf{\large\sffamily Soru 173.} \ Ortak malların trajedisinde meraların aşırı otlatılması gibi bir durumun önüne geçmek için uygulanan 'kota sistemi' neyi ifade eder?
\begin{enumerate}[label=\Alph*), leftmargin=25pt, topsep=4pt, itemsep=2pt]
    \item Meraların tamamen kapatılmasını
    \item Her çobanın meraya salabileceği maksimum hayvan sayısının yasal olarak sınırlandırılmasını
    \item Her çobana ücretsiz yem verilmesini
    \item Hayvanların vergilendirilmesini
\end{enumerate}
\vspace{4pt}
\begin{tcolorbox}[title=Cevap ve Açıklama, colback=gray!3!white, colframe=primarycolor!90!black, arc=2mm, fonttitle=\bfseries\sffamily, top=6pt, bottom=6pt]
    \textbf{Doğru Cevap: B}
    \vspace{2pt}
    Kota sistemi, kaynağın kendini yenileyebilme kapasitesini aşmamak için kullanıcılara getirilen yasal tüketim sınırlandırmasıdır.
\end{tcolorbox}
\vspace{10pt}

\needspace{5cm}
\vspace{8pt}
\noindent\textbf{\large\sffamily Soru 174.} \ Aşağıdakilerden hangisi bir işletmenin 'sosyal sürdürülebilirlik' ilkelerine aykırı hareket ettiğini gösterir?
\begin{enumerate}[label=\Alph*), leftmargin=25pt, topsep=4pt, itemsep=2pt]
    \item Çalışanlarına sendikalaşma ve toplu sözleşme hakkı tanıması
    \item Kadın ve erkek çalışanlara aynı iş için farklı ücretler (ücret adaletsizliği) ödemesi
    \item Çalışma saatlerini yasal sınırlarda tutması
    \item İş sağlığı ve güvenliği eğitimleri vermesi
\end{enumerate}
\vspace{4pt}
\begin{tcolorbox}[title=Cevap ve Açıklama, colback=gray!3!white, colframe=primarycolor!90!black, arc=2mm, fonttitle=\bfseries\sffamily, top=6pt, bottom=6pt]
    \textbf{Doğru Cevap: B}
    \vspace{2pt}
    Cinsiyete dayalı ücret ayrımcılığı fırsat eşitliği ve sosyal adalet ilkelerine aykırı olup, sosyal sürdürülebilirliği zedeler.
\end{tcolorbox}
\vspace{10pt}

\needspace{5cm}
\vspace{8pt}
\noindent\textbf{\large\sffamily Soru 175.} \ Ortak Malların Trajedisi kavramının yaratıcısı Garrett Hardin hangi bilim dalında çalışmalar yapmıştır?
\begin{enumerate}[label=\Alph*), leftmargin=25pt, topsep=4pt, itemsep=2pt]
    \item Yazılım Mühendisliği
    \item Ekoloji ve Biyoloji
    \item Finansal Muhasebe
    \item Siyaset Bilimi
\end{enumerate}
\vspace{4pt}
\begin{tcolorbox}[title=Cevap ve Açıklama, colback=gray!3!white, colframe=primarycolor!90!black, arc=2mm, fonttitle=\bfseries\sffamily, top=6pt, bottom=6pt]
    \textbf{Doğru Cevap: B}
    \vspace{2pt}
    Garrett Hardin Amerikalı bir ekolog ve biyologdur. Çevresel kıtlıklar ve nüfus artışı üzerine çalışmalar yürütmüştür.
\end{tcolorbox}
\vspace{10pt}

\needspace{5cm}
\vspace{8pt}
\noindent\textbf{\large\sffamily Soru 176.} \ Lojistikte, ürünlerin değer kazanımı, geri kazanımı, yenilenmesi veya uygun şekilde imha edilmesi amacıyla tüketim noktasından çıkış noktasına geri akışını yöneten süreç hangisidir?
\begin{enumerate}[label=\Alph*), leftmargin=25pt, topsep=4pt, itemsep=2pt]
    \item İleri Lojistik
    \item Tersine Lojistik (Reverse Logistics)
    \item Intermodal Lojistik
    \item Soğuk Zincir Lojistiği
\end{enumerate}
\vspace{4pt}
\begin{tcolorbox}[title=Cevap ve Açıklama, colback=gray!3!white, colframe=primarycolor!90!black, arc=2mm, fonttitle=\bfseries\sffamily, top=6pt, bottom=6pt]
    \textbf{Doğru Cevap: B}
    \vspace{2pt}
    Tersine lojistik, geri dönüşüm, tamirat, iade yönetimi ve yeniden üretim süreçleri için ürünlerin tüketimden kaynağa doğru ters yönlü akışıdır.
\end{tcolorbox}
\vspace{10pt}

\needspace{5cm}
\vspace{8pt}
\noindent\textbf{\large\sffamily Soru 177.} \ Yeniden üretim (re-manufacturing) süreçlerinin, hammadde kullanarak sıfırdan üretime kıyasla sağladığı enerji tasarrufu oranı en fazla yüzde kaçtır?
\begin{enumerate}[label=\Alph*), leftmargin=25pt, topsep=4pt, itemsep=2pt]
    \item \%10
    \item \%50
    \item \%85
    \item Tasarruf sağlamaz
\end{enumerate}
\vspace{4pt}
\begin{tcolorbox}[title=Cevap ve Açıklama, colback=gray!3!white, colframe=primarycolor!90!black, arc=2mm, fonttitle=\bfseries\sffamily, top=6pt, bottom=6pt]
    \textbf{Doğru Cevap: C}
    \vspace{2pt}
    Yeniden üretim süreçleri, maden çıkarma ve ilk işleme adımlarını atladığı için sıfırdan üretime göre \%85'e varan enerji tasarrufu sağlar.
\end{tcolorbox}
\vspace{10pt}

\needspace{5cm}
\vspace{8pt}
\noindent\textbf{\large\sffamily Soru 178.} \ Geri dönüştürülmüş alüminyum üretmek, sıfırdan boksit madeninden alüminyum üretmek için gereken enerjinin yalnızca yüzde kaçını tüketir?
\begin{enumerate}[label=\Alph*), leftmargin=25pt, topsep=4pt, itemsep=2pt]
    \item \%5'ini
    \item \%25'ini
    \item \%50'sini
    \item \%85'ini
\end{enumerate}
\vspace{4pt}
\begin{tcolorbox}[title=Cevap ve Açıklama, colback=gray!3!white, colframe=primarycolor!90!black, arc=2mm, fonttitle=\bfseries\sffamily, top=6pt, bottom=6pt]
    \textbf{Doğru Cevap: A}
    \vspace{2pt}
    Maden işleme çok yüksek elektrik gerektirir. Alüminyumu geri dönüştürmek ise bu enerjinin sadece \%5'i ile (yani \%95 tasarrufla) gerçekleştirilir.
\end{tcolorbox}
\vspace{10pt}

\needspace{5cm}
\vspace{8pt}
\noindent\textbf{\large\sffamily Soru 179.} \ 3M firmasının lojistik araçlarında yükleme yüksekliğini ayarlayan raflar geliştirerek araç doluluk oranlarında elde ettiği artış oranı yüzde kaçtır?
\begin{enumerate}[label=\Alph*), leftmargin=25pt, topsep=4pt, itemsep=2pt]
    \item \%10
    \item \%20
    \item \%40
    \item \%80
\end{enumerate}
\vspace{4pt}
\begin{tcolorbox}[title=Cevap ve Açıklama, colback=gray!3!white, colframe=primarycolor!90!black, arc=2mm, fonttitle=\bfseries\sffamily, top=6pt, bottom=6pt]
    \textbf{Doğru Cevap: C}
    \vspace{2pt}
    3M, araç doluluk oranını \%40 oranında artırarak sefer sayısını azaltmış ve lojistik kaynaklı karbon emisyonlarını düşürmüştür.
\end{tcolorbox}
\vspace{10pt}

\needspace{5cm}
\vspace{8pt}
\noindent\textbf{\large\sffamily Soru 180.} \ İstanbul, Trieste ve Mannheim arasında karayolu yerine Ro-Ro gemisi ve tren hatlarını birleştirerek karbon salınımını azaltan Ekol Lojistik modeli hangisidir?
\begin{enumerate}[label=\Alph*), leftmargin=25pt, topsep=4pt, itemsep=2pt]
    \item Havayolu Taşımacılığı
    \item İntermodal Taşımacılık (Intermodal Transport)
    \item Karayolu Kamyon Lojistiği
    \item Kurye Dağıtımı
\end{enumerate}
\vspace{4pt}
\begin{tcolorbox}[title=Cevap ve Açıklama, colback=gray!3!white, colframe=primarycolor!90!black, arc=2mm, fonttitle=\bfseries\sffamily, top=6pt, bottom=6pt]
    \textbf{Doğru Cevap: B}
    \vspace{2pt}
    İntermodal taşımacılık, yükün araç değiştirmeden birden fazla taşıma türüyle (deniz, demir, kara) taşınarak çevre dostu lojistik kurulmasıdır.
\end{tcolorbox}
\vspace{10pt}

\needspace{5cm}
\vspace{8pt}
\noindent\textbf{\large\sffamily Soru 181.} \ Çevre dostu binaların (yeşil binalar) geleneksel binalara kıyasla sağladığı enerji tasarrufu aralığı aşağıdakilerden hangisidir?
\begin{enumerate}[label=\Alph*), leftmargin=25pt, topsep=4pt, itemsep=2pt]
    \item \%5 - \%10
    \item \%24 - \%50
    \item \%80 - \%95
    \item Tasarruf sağlamaz
\end{enumerate}
\vspace{4pt}
\begin{tcolorbox}[title=Cevap ve Açıklama, colback=gray!3!white, colframe=primarycolor!90!black, arc=2mm, fonttitle=\bfseries\sffamily, top=6pt, bottom=6pt]
    \textbf{Doğru Cevap: B}
    \vspace{2pt}
    Yeşil binalar yalıtım, güneş paneli ve akıllı otomasyon sistemleri sayesinde geleneksel binalara göre \%24 ile \%50 arasında enerji tasarrufu sunar.
\end{tcolorbox}
\vspace{10pt}

\needspace{5cm}
\vspace{8pt}
\noindent\textbf{\large\sffamily Soru 182.} \ Tedarik zincirinde paketleme malzemelerinin biyolojik olarak doğada çözünebilen karton veya nişasta bazlı malzemelerden seçilmesi hangisidir?
\begin{enumerate}[label=\Alph*), leftmargin=25pt, topsep=4pt, itemsep=2pt]
    \item Yeşil Satın Alma
    \item Yeşil Paketleme (Green Packaging)
    \item Tersine Lojistik
    \item İntermodal Dağıtım
\end{enumerate}
\vspace{4pt}
\begin{tcolorbox}[title=Cevap ve Açıklama, colback=gray!3!white, colframe=primarycolor!90!black, arc=2mm, fonttitle=\bfseries\sffamily, top=6pt, bottom=6pt]
    \textbf{Doğru Cevap: B}
    \vspace{2pt}
    Plastik yerine doğa dostu, çözünür ve tekrar kullanılabilir paketleme malzemelerinin kullanılması Yeşil Paketleme (Green Packaging) kapsamındadır.
\end{tcolorbox}
\vspace{10pt}

\needspace{5cm}
\vspace{8pt}
\noindent\textbf{\large\sffamily Soru 183.} \ Lojistikte araçların motor teknolojilerini iyileştirerek (Euro 5/6 standardı) veya elektrikli araçlarla karbon salınımını düşürmeyi hedefleyen yeşil strateji hangisidir?
\begin{enumerate}[label=\Alph*), leftmargin=25pt, topsep=4pt, itemsep=2pt]
    \item Daha Az Taşı
    \item Daha Kısa Taşı
    \item Daha Temiz Taşı (Transport Cleaner)
    \item Tersine Taşı
\end{enumerate}
\vspace{4pt}
\begin{tcolorbox}[title=Cevap ve Açıklama, colback=gray!3!white, colframe=primarycolor!90!black, arc=2mm, fonttitle=\bfseries\sffamily, top=6pt, bottom=6pt]
    \textbf{Doğru Cevap: C}
    \vspace{2pt}
    Daha Temiz Taşı stratejisi, taşıma araçlarının yakıt verimliliğini artırmayı, temiz enerji ve Euro 5/6 motor standartlarını kullanmayı hedefler.
\end{tcolorbox}
\vspace{10pt}

\needspace{5cm}
\vspace{8pt}
\noindent\textbf{\large\sffamily Soru 184.} \ Lojistik süreçleri planlarken rota optimizasyonu yazılımları kullanarak kat edilen toplam mesafeyi azaltmayı hedefleyen yeşil lojistik stratejisi hangisidir?
\begin{enumerate}[label=\Alph*), leftmargin=25pt, topsep=4pt, itemsep=2pt]
    \item Daha Az Taşı
    \item Daha Kısa Taşı (Transport Shorter)
    \item Daha Temiz Taşı
    \item Tersine Taşı
\end{enumerate}
\vspace{4pt}
\begin{tcolorbox}[title=Cevap ve Açıklama, colback=gray!3!white, colframe=primarycolor!90!black, arc=2mm, fonttitle=\bfseries\sffamily, top=6pt, bottom=6pt]
    \textbf{Doğru Cevap: B}
    \vspace{2pt}
    Daha Kısa Taşı stratejisi, akıllı rota planlama (routing) ve depoların doğru konumlandırılmasıyla kat edilen kilometreyi ve emisyonları azaltmayı amaçlar.
\end{tcolorbox}
\vspace{10pt}

\needspace{5cm}
\vspace{8pt}
\noindent\textbf{\large\sffamily Soru 185.} \ Tedarik zincirinde tedarikçi seçim kriterlerine 'karbon ayak izi', 'çevre sertifikaları (ISO 14001)' gibi maddelerin eklenmesi hangisidir?
\begin{enumerate}[label=\Alph*), leftmargin=25pt, topsep=4pt, itemsep=2pt]
    \item Yeşil Satın Alma (Green Procurement)
    \item Tersine Dağıtım
    \item Fiyat Odaklı Satın Alma
    \item Tedarikçi Entegrasyonu
\end{enumerate}
\vspace{4pt}
\begin{tcolorbox}[title=Cevap ve Açıklama, colback=gray!3!white, colframe=primarycolor!90!black, arc=2mm, fonttitle=\bfseries\sffamily, top=6pt, bottom=6pt]
    \textbf{Doğru Cevap: A}
    \vspace{2pt}
    Yeşil Satın Alma, hammadde ve hizmet alımlarında çevresel performans kriterlerinin (ISO 14001, geri dönüştürülebilir malzeme vb.) önceliklendirilmesidir.
\end{tcolorbox}
\vspace{10pt}

\needspace{5cm}
\vspace{8pt}
\noindent\textbf{\large\sffamily Soru 186.} \ Yeşil binaların tescillenmesinde dünyada en yaygın kullanılan sertifikasyon sistemlerinden biri aşağıdakilerden hangisidir?
\begin{enumerate}[label=\Alph*), leftmargin=25pt, topsep=4pt, itemsep=2pt]
    \item ISO 9001
    \item LEED (Leadership in Energy and Environmental Design)
    \item CMMI Level 5
    \item Brundtland Belgesi
\end{enumerate}
\vspace{4pt}
\begin{tcolorbox}[title=Cevap ve Açıklama, colback=gray!3!white, colframe=primarycolor!90!black, arc=2mm, fonttitle=\bfseries\sffamily, top=6pt, bottom=6pt]
    \textbf{Doğru Cevap: B}
    \vspace{2pt}
    LEED sertifikası, binaların enerji verimliliği, su tasarrufu, malzeme seçimi gibi çevresel sürdürülebilirlik kriterlerini ölçen küresel derecelendirme sistemidir.
\end{tcolorbox}
\vspace{10pt}

\needspace{5cm}
\vspace{8pt}
\noindent\textbf{\large\sffamily Soru 187.} \ Tersine lojistikte 'yeniden kullanım' (reuse) ile 'geri dönüşüm' (recycling) arasındaki fark hangisidir?
\begin{enumerate}[label=\Alph*), leftmargin=25pt, topsep=4pt, itemsep=2pt]
    \item Geri dönüşüm daha ucuzdur
    \item Yeniden kullanım ürünü yapısını bozmadan tekrar kullanırken, geri dönüşüm ürünü eriterek hammaddesine dönüştürür
    \item Yeniden kullanım sadece plastiklerde geçerlidir
    \item Hiçbir fark yoktur
\end{enumerate}
\vspace{4pt}
\begin{tcolorbox}[title=Cevap ve Açıklama, colback=gray!3!white, colframe=primarycolor!90!black, arc=2mm, fonttitle=\bfseries\sffamily, top=6pt, bottom=6pt]
    \textbf{Doğru Cevap: B}
    \vspace{2pt}
    Yeniden kullanımda (örneğin cam şişelerin yıkanıp tekrar doldurulması) enerji harcanmaz. Geri dönüşümde ise ürün (örneğin camın eritilmesi) hammaddeye dönüştürülür ve enerji harcanır.
\end{tcolorbox}
\vspace{10pt}

\needspace{5cm}
\vspace{8pt}
\noindent\textbf{\large\sffamily Soru 188.} \ Tedarik zincirinde 'Karbon Ayak İzi' (Carbon Footprint) neyi temsil eder?
\begin{enumerate}[label=\Alph*), leftmargin=25pt, topsep=4pt, itemsep=2pt]
    \item Fabrikada çalışan işçi sayısını
    \item Üretim ve lojistik faaliyetleri sonucu doğrudan veya dolaylı olarak atmosfere salınan sera gazlarının karbondioksit (CO2) cinsinden toplam miktarını
    \item Kömür depolarının fiziksel alanını
    \item Fabrika bacalarının uzunluğunu
\end{enumerate}
\vspace{4pt}
\begin{tcolorbox}[title=Cevap ve Açıklama, colback=gray!3!white, colframe=primarycolor!90!black, arc=2mm, fonttitle=\bfseries\sffamily, top=6pt, bottom=6pt]
    \textbf{Doğru Cevap: B}
    \vspace{2pt}
    Karbon ayak izi, bir kurumun veya ürünün yaşam döngüsü boyunca iklim değişikliğine neden olan sera gazı emisyonlarının ton CO2 eşdeğeri olarak ölçülmesidir.
\end{tcolorbox}
\vspace{10pt}

\needspace{5cm}
\vspace{8pt}
\noindent\textbf{\large\sffamily Soru 189.} \ Lojistikte 'Daha Az Taşı' (Transport Less) yeşil stratejisi aşağıdakilerden hangisiyle gerçekleştirilebilir?
\begin{enumerate}[label=\Alph*), leftmargin=25pt, topsep=4pt, itemsep=2pt]
    \item Üretimi tamamen durdurarak
    \item Paketleme boyutlarını küçülterek ve araç doluluk oranlarını (konsolidasyon) maksimize ederek sefer sayılarını azaltmak
    \item Kamyon hızlarını düşürerek
    \item Sadece yerel pazarda satış yapmak
\end{enumerate}
\vspace{4pt}
\begin{tcolorbox}[title=Cevap ve Açıklama, colback=gray!3!white, colframe=primarycolor!90!black, arc=2mm, fonttitle=\bfseries\sffamily, top=6pt, bottom=6pt]
    \textbf{Doğru Cevap: B}
    \vspace{2pt}
    Daha Az Taşı stratejisi, yükleme planlama algoritmaları ve paketleme optimizasyonu ile gereksiz taşımaları (boş git-gelleri) önlemeyi amaçlar.
\end{tcolorbox}
\vspace{10pt}

\needspace{5cm}
\vspace{8pt}
\noindent\textbf{\large\sffamily Soru 190.} \ Geri kazanım süreçlerinde 'upcycling' (ileri dönüşüm) kavramı neyi ifade eder?
\begin{enumerate}[label=\Alph*), leftmargin=25pt, topsep=4pt, itemsep=2pt]
    \item Ürünün eritilip kalitesiz plastik yapılması
    \item Atık veya eski malzemelerin orijinalinden daha yüksek değere veya kaliteye sahip yeni bir ürüne dönüştürülmesi
    \item Atıkların yakılarak yok edilmesi
    \item Ürünün yurtdışına satılması
\end{enumerate}
\vspace{4pt}
\begin{tcolorbox}[title=Cevap ve Açıklama, colback=gray!3!white, colframe=primarycolor!90!black, arc=2mm, fonttitle=\bfseries\sffamily, top=6pt, bottom=6pt]
    \textbf{Doğru Cevap: B}
    \vspace{2pt}
    İleri dönüşüm (upcycling), atık malzemelerin yaratıcı fikirlerle (örneğin araba lastiğinden lüks mobilya tasarlamak) katma değerli hale getirilmesidir.
\end{tcolorbox}
\vspace{10pt}

\needspace{5cm}
\vspace{8pt}
\noindent\textbf{\large\sffamily Soru 191.} \ Ekolojik faktörleri korumaya odaklanan 'kentsel sürdürülebilirlik' ile teknoloji odaklı 'akıllı şehir' kavramlarını sürdürülebilir akıllı şehir çatısında birleştiren Neirotti vd. (2014) çalışmasına göre akıllı şehirleşmenin esası nedir?
\begin{enumerate}[label=\Alph*), leftmargin=25pt, topsep=4pt, itemsep=2pt]
    \item Şehir nüfusunun azaltılmasıdır
    \item Bilgi ve iletişim teknolojilerini kullanarak kaynakları etkin yönetmek, karbon ayak izini azaltmak ve yaşam kalitesini artırmaktır
    \item Şehirdeki tüm yolların paralı yapılmasıdır
    \item Sanayi tesislerinin tamamen kapatılmasıdır
\end{enumerate}
\vspace{4pt}
\begin{tcolorbox}[title=Cevap ve Açıklama, colback=gray!3!white, colframe=primarycolor!90!black, arc=2mm, fonttitle=\bfseries\sffamily, top=6pt, bottom=6pt]
    \textbf{Doğru Cevap: B}
    \vspace{2pt}
    Sürdürülebilir akıllı şehir, teknolojiyi çevre koruma, kaynak verimliliği ve kentsel yaşam kalitesini artırma aracı olarak entegre eder.
\end{tcolorbox}
\vspace{10pt}

\needspace{5cm}
\vspace{8pt}
\noindent\textbf{\large\sffamily Soru 192.} \ Küresel enerji tüketiminin yüzde 75'i ve karbondioksit (CO2) salınımının yüzde 80'i aşağıdakilerden hangisinden kaynaklanmaktadır?
\begin{enumerate}[label=\Alph*), leftmargin=25pt, topsep=4pt, itemsep=2pt]
    \item Tarım ve hayvancılık alanlarından
    \item Şehirlerden (kentsel alanlardan)
    \item Okyanus taşımacılığından
    \item Orman yangınlarından
\end{enumerate}
\vspace{4pt}
\begin{tcolorbox}[title=Cevap ve Açıklama, colback=gray!3!white, colframe=primarycolor!90!black, arc=2mm, fonttitle=\bfseries\sffamily, top=6pt, bottom=6pt]
    \textbf{Doğru Cevap: B}
    \vspace{2pt}
    Şehirler, nüfusun ve sanayinin yoğunlaşması nedeniyle küresel enerjinin \%75'ini tüketmekte ve emisyonların \%80'ine neden olmaktadır.
\end{tcolorbox}
\vspace{10pt}

\needspace{5cm}
\vspace{8pt}
\noindent\textbf{\large\sffamily Soru 193.} \ Gül ve Çobanoğlu (2017) çalışmasına göre akıllı şehir mimarisinin entegre olması gereken dört temel bileşen hangisidir?
\begin{enumerate}[label=\Alph*), leftmargin=25pt, topsep=4pt, itemsep=2pt]
    \item Finans, Politika, Reklam ve İthalat
    \item Bilgi, Teknoloji, Sürdürülebilirlik ve Katılım
    \item Binalar, Yollar, Arabalar ve Fabrikalar
    \item İnternet, Bilgisayar, Yazılım ve Donanım
\end{enumerate}
\vspace{4pt}
\begin{tcolorbox}[title=Cevap ve Açıklama, colback=gray!3!white, colframe=primarycolor!90!black, arc=2mm, fonttitle=\bfseries\sffamily, top=6pt, bottom=6pt]
    \textbf{Doğru Cevap: B}
    \vspace{2pt}
    Gül ve Çobanoğlu'na göre akıllı kent; bilginin üretildiği, teknolojinin kullanıldığı, ekolojik dengenin (sürdürülebilirlik) korunduğu ve vatandaşların yönetime katıldığı yapıdır.
\end{tcolorbox}
\vspace{10pt}

\needspace{5cm}
\vspace{8pt}
\noindent\textbf{\large\sffamily Soru 194.} \ Projeksiyonlara göre bugün dünya nüfusunun yüzde 55'i kentlerde yaşarken, 2050 yılına kadar bu oranın yüzde kaça ulaşması beklenmektedir?
\begin{enumerate}[label=\Alph*), leftmargin=25pt, topsep=4pt, itemsep=2pt]
    \item \%60
    \item \%68
    \item \%85
    \item \%95
\end{enumerate}
\vspace{4pt}
\begin{tcolorbox}[title=Cevap ve Açıklama, colback=gray!3!white, colframe=primarycolor!90!black, arc=2mm, fonttitle=\bfseries\sffamily, top=6pt, bottom=6pt]
    \textbf{Doğru Cevap: B}
    \vspace{2pt}
    Birleşmiş Milletler nüfus projeksiyonlarına göre hızlı göç ve kentleşme sonucu 2050 yılında dünya nüfusunun \%68'i şehirlerde yaşayacaktır.
\end{tcolorbox}
\vspace{10pt}

\needspace{5cm}
\vspace{8pt}
\noindent\textbf{\large\sffamily Soru 195.} \ UNEP ve UNWTO sürdürülebilir turizm ilkelerinden hangisi turizm gelirlerinin yerel halk ve yerel üreticilere doğrudan ulaşmasını hedefler?
\begin{enumerate}[label=\Alph*), leftmargin=25pt, topsep=4pt, itemsep=2pt]
    \item Ekonomik Süreklilik
    \item Yerel Kalkınma (Local Prosperity)
    \item İstihdam Kalitesi
    \item Çevresel Etki
\end{enumerate}
\vspace{4pt}
\begin{tcolorbox}[title=Cevap ve Açıklama, colback=gray!3!white, colframe=primarycolor!90!black, arc=2mm, fonttitle=\bfseries\sffamily, top=6pt, bottom=6pt]
    \textbf{Doğru Cevap: B}
    \vspace{2pt}
    Yerel kalkınma (local prosperity) ilkesi, turist harcamalarının destinasyondaki yerel esnaf, çiftçi ve üreticide kalma oranını artırarak sızıntıları (leakage) önlemeyi amaçlar.
\end{tcolorbox}
\vspace{10pt}

\needspace{5cm}
\vspace{8pt}
\noindent\textbf{\large\sffamily Soru 196.} \ Sürdürülebilir turizm ilkelerinde 'biyolojik çeşitlilik' (biodiversity) kavramı neyi önceler?
\begin{enumerate}[label=\Alph*), leftmargin=25pt, topsep=4pt, itemsep=2pt]
    \item Turistlerin her alana serbestçe girebilmesini
    \item Koruma-kullanma dengesi gözetilerek doğal ekosistemlerin, endemik flora ve faunanın korunmasını
    \item Şehirdeki hayvanat bahçesi sayısını artırmayı
    \item Tarım arazilerini turizme açmayı
\end{enumerate}
\vspace{4pt}
\begin{tcolorbox}[title=Cevap ve Açıklama, colback=gray!3!white, colframe=primarycolor!90!black, arc=2mm, fonttitle=\bfseries\sffamily, top=6pt, bottom=6pt]
    \textbf{Doğru Cevap: B}
    \vspace{2pt}
    Biyolojik çeşitlilik ilkesi, hassas ekosistemlerin (örneğin mercan resifleri, milli parklar) turizm baskısıyla tahrip edilmesini engellemeyi amaçlar.
\end{tcolorbox}
\vspace{10pt}

\needspace{5cm}
\vspace{8pt}
\noindent\textbf{\large\sffamily Soru 197.} \ Kültürel sürdürülebilirlik ve süreklilik kapsamında vurgulanmak istenen ana düşünce aşağıdakilerden hangisidir?
\begin{enumerate}[label=\Alph*), leftmargin=25pt, topsep=4pt, itemsep=2pt]
    \item Kültürün hiçbir şekilde değiştirilmeden eski haliyle dondurulması
    \item Geçmişin mimari, yerleşim ve toplumsal ilkelerinin günümüz için sürdürülebilir büyük birer değer kaynağı olduğunun fark edilmesi
    \item Küresel popüler kültürün yerel kültürü yok etmesi
    \item Tarihi binaların yıkılıp yerine akıllı gökdelenler yapılması
\end{enumerate}
\vspace{4pt}
\begin{tcolorbox}[title=Cevap ve Açıklama, colback=gray!3!white, colframe=primarycolor!90!black, arc=2mm, fonttitle=\bfseries\sffamily, top=6pt, bottom=6pt]
    \textbf{Doğru Cevap: B}
    \vspace{2pt}
    Kültürel sürdürülebilirlik gelenekçi/donmuş bir tavır değildir. Geçmiş yerleşim ve yaşam pratiklerinin (örneğin kerpiç evler, su sarnıçları gibi yeşil tekniklerin) değer kaynağı olarak günümüze aktarılmasıdır.
\end{tcolorbox}
\vspace{10pt}

\needspace{5cm}
\vspace{8pt}
\noindent\textbf{\large\sffamily Soru 198.} \ Ekonomik sürdürülebilirlik kapsamında Burns ve Holden (1995) çalışmasına göre ikame olanağı bulunmayan hangi kaynakların kullanımına sınır getirilmelidir?
\begin{enumerate}[label=\Alph*), leftmargin=25pt, topsep=4pt, itemsep=2pt]
    \item Yenilenebilir kaynaklar (güneş, rüzgar)
    \item Yenilenemeyen kaynaklar (fosil yakıtlar, madenler)
    \item Dijital veri kaynakları
    \item İnsan gücü kaynakları
\end{enumerate}
\vspace{4pt}
\begin{tcolorbox}[title=Cevap ve Açıklama, colback=gray!3!white, colframe=primarycolor!90!black, arc=2mm, fonttitle=\bfseries\sffamily, top=6pt, bottom=6pt]
    \textbf{Doğru Cevap: B}
    \vspace{2pt}
    Burns ve Holden'a göre ikamesi olmayan, tükenebilir (yenilenemeyen) kaynakların tüketimine sınırlama getirmek ekonomik sürdürülebilirliğin temel şartıdır.
\end{tcolorbox}
\vspace{10pt}

\needspace{5cm}
\vspace{8pt}
\noindent\textbf{\large\sffamily Soru 199.} \ Ateş (2018) çalışmasına göre sürdürülebilir akıllı şehirleşmede kentsel karbon ayak izini azaltmak için atılması gereken en öncelikli ulaşım adımı hangisidir?
\begin{enumerate}[label=\Alph*), leftmargin=25pt, topsep=4pt, itemsep=2pt]
    \item Bireysel otomobil kullanımını teşvik etmek
    \item Toplu taşıma sistemlerinin raylı ve elektrikli alternatiflerle güçlendirilmesi ve mikromobilite (bisiklet vb.) entegrasyonu
    \item Tüm yollara hız sınırı koymak
    \item Toplu taşımayı tamamen kaldırmak
\end{enumerate}
\vspace{4pt}
\begin{tcolorbox}[title=Cevap ve Açıklama, colback=gray!3!white, colframe=primarycolor!90!black, arc=2mm, fonttitle=\bfseries\sffamily, top=6pt, bottom=6pt]
    \textbf{Doğru Cevap: B}
    \vspace{2pt}
    Kentsel ulaşımdan kaynaklanan karbon emisyonlarını azaltmanın en etkin yolu elektrikli toplu taşıma, raylı sistemler ve bisiklet gibi çevre dostu ulaşım entegrasyonlarıdır.
\end{tcolorbox}
\vspace{10pt}

\needspace{5cm}
\vspace{8pt}
\noindent\textbf{\large\sffamily Soru 200.} \ Sürdürülebilir kalkınmanın gerçekleşmesinde ekolojik dengenin korunması neden zorunludur?
\begin{enumerate}[label=\Alph*), leftmargin=25pt, topsep=4pt, itemsep=2pt]
    \item Teknoloji fiyatlarının ucuzlaması buna bağlı olduğu için
    \item Ekonomik faaliyetlerin sürekliliğinin, ekosistemin sunduğu su, hava, toprak ve hammadde gibi kaynakların varlığına bağımlı olması nedeniyle
    \item İnsanların köylere göç etmesini sağlamak için
    \item Fabrika kârlarını artırmak için
\end{enumerate}
\vspace{4pt}
\begin{tcolorbox}[title=Cevap ve Açıklama, colback=gray!3!white, colframe=primarycolor!90!black, arc=2mm, fonttitle=\bfseries\sffamily, top=6pt, bottom=6pt]
    \textbf{Doğru Cevap: B}
    \vspace{2pt}
    Ekolojik denge bozulduğunda su kıtlığı, iklim krizleri ve kaynak tükenmesi başlar; bu da ekonomik üretimin ve kalkınmanın tamamen durmasına yol açar.
\end{tcolorbox}
\vspace{10pt}


\end{document}
